% Systèmes d'équations linéaires à deux ou trois inconnues — Mathématiques 1ère E — Cours signé OnBuch+
% Programme officiel MINESEC. Compiler avec : tectonic N02-systemes-equations.tex
\newcommand{\DOCMATIERE}{Mathématiques}
\newcommand{\DOCNIVEAU}{1ère E}
\newcommand{\DOCMODULE}{Relations et opérations fondamentales dans l'ensemble des nombres réels}
\newcommand{\DOCLECON}{N02}
\newcommand{\DOCTITRE}{Systèmes d'équations linéaires à deux ou trois inconnues}
% OnBuch+ — préambule « cours signés » ENRICHI (compilé avec tectonic / XeLaTeX)
% Système visuel « Pop » d'OnBuch+ : fond crème (jamais blanc), bordures encre
% épaisses, ombres « dures » décalées, titres Archivo Black en capitales.
% Version MATHÉMATIQUES : graphes (pgfplots/tikz), boîtes pédagogiques
% (définition, propriété, méthode, attention, le savais-tu, exercice résolu,
% exercice, TP, bilan), page de garde et signature OnBuch+ ancrée en bas.
%
% Macros attendues dans main.tex AVANT \input{preamble} :
%   \DOCMATIERE  \DOCNIVEAU  \DOCMODULE  \DOCLECON (numéro)  \DOCTITRE
\documentclass[11pt]{article}

\usepackage{fontspec}
\usepackage[french]{babel}
\usepackage[a4paper,margin=2cm,top=3.4cm,bottom=2.8cm,headheight=1.4cm,footskip=1.3cm]{geometry}
\usepackage{amsmath,amssymb,amsfonts}
\usepackage{mathtools} % \xmapsto, \coloneqq... souvent utilisés par les modèles
\usepackage{cancel} % simplifications barrées
% \abs{z} (module, valeur absolue) et \norm{u} (norme), utilisés par les modèles
\providecommand{\abs}{}\let\abs\relax\DeclarePairedDelimiter{\abs}{\lvert}{\rvert}
\providecommand{\norm}{}\let\norm\relax\DeclarePairedDelimiter{\norm}{\lVert}{\rVert}
\usepackage[table]{xcolor} % [table] : fournit aussi \rowcolors (alternance de lignes)
\usepackage{pagecolor}
\usepackage{tikz}
\usetikzlibrary{arrows.meta,positioning,calc,shapes.geometric,shapes.misc,shapes.arrows,decorations.pathmorphing,decorations.pathreplacing,decorations.markings,patterns,shadows,mindmap,trees,fit,backgrounds,matrix,intersections,angles,quotes}
\usepackage{pgfplots}
\usepackage{pgf-pie} % diagrammes circulaires \pie (statistiques)
\pgfplotsset{compat=1.18}
\usepgfplotslibrary{fillbetween,groupplots} % aires entre courbes (calcul intégral)
\usepackage{enumitem}
\usepackage{fancyhdr}
% [most] charge toutes les bibliothèques tcolorbox (listings, minted, raster,
% documentation...) et gonfle inutilement la mémoire TeX sur un document long
% (~300 boîtes) : on ne charge que ce qui est réellement utilisé.
\usepackage[skins,breakable]{tcolorbox}
\usepackage{graphicx}
\usepackage{colortbl}
\usepackage{booktabs}
\usepackage{tabularx}
\usepackage{diagbox} % case d'en-tête diagonale des tableaux à double entrée
% Colonnes extensibles centrée / gauche / droite (C, L, R), souvent utilisées par les modèles
\newcolumntype{C}{>{\centering\arraybackslash}X}
\newcolumntype{L}{>{\raggedright\arraybackslash}X}
\newcolumntype{R}{>{\raggedleft\arraybackslash}X}
\usepackage{array}
\usepackage{multicol}
\usepackage{siunitx}
% parse-numbers=false : accepte aussi \SI{2,5 \times 10^{-3}}{...}
\sisetup{locale=FR,output-decimal-marker={,},group-minimum-digits=4,parse-numbers=false}
\DeclareSIUnit\ohm{\ensuremath{\symup{\Omega}}} % U+2126 absent de la police
\usepackage{qrcode}
% \degree : commande fréquemment utilisée par les modèles pour un angle ou une
% température en dehors de \SI{}{} (non fournie par les paquets chargés).
\providecommand{\degree}{^{\circ}}
% \qed (fin de démonstration), utilisé par les modèles sans amsthm
\providecommand{\qed}{\hfill$\square$}
% \barre : double barre des valeurs interdites dans un tableau de variations (tkz-tab)
\providecommand{\barre}{\ensuremath{\|}}
% environnement proof (amsthm non chargé) utilisé par les modèles
\ifdefined\proof\else\newenvironment{proof}[1][Démonstration]{\par\noindent\textit{#1.}\ }{\qed\par}\fi
\usepackage{lastpage}
\usepackage{hyperref}
\hypersetup{hidelinks}

% ============================================================
% Palette « Pop » OnBuch+ — mêmes hex que la classe P (plus_ui.dart)
% ============================================================
\definecolor{poporange}{HTML}{F59321}
\definecolor{popdark}{HTML}{E07A0C}
\definecolor{popcream}{HTML}{FFF6E8}
\definecolor{popink}{HTML}{1C1714}
\definecolor{poppurple}{HTML}{7A5AE0}
\definecolor{popgreen}{HTML}{2E9E6A}
\definecolor{poppink}{HTML}{E0457B}
\definecolor{popblue}{HTML}{2F7DE1}
\definecolor{popgold}{HTML}{C98A12}
\definecolor{popmuted}{HTML}{8A7F76}
\definecolor{popline}{HTML}{EADFD2}
\definecolor{popgray}{HTML}{8A7F76}
\colorlet{popgrey}{popgray}
% Alias tolérants (le modèle nomme parfois les couleurs différemment)
\colorlet{popwhite}{white}
\colorlet{popblack}{popink}
\colorlet{popred}{poppink}
\colorlet{popyellow}{popgold}
% Teintes claires pour fonds de tableaux / zones
\colorlet{popblueL}{popblue!10!white}
\colorlet{poporangeL}{poporange!14!white}
\colorlet{popgreenL}{popgreen!12!white}
\colorlet{poppurpleL}{poppurple!12!white}
\colorlet{poppinkL}{poppink!12!white}
\colorlet{popgoldL}{popgold!14!white}

\pagecolor{popcream}

% ============================================================
% Typographie
% ============================================================
\defaultfontfeatures{Ligatures=TeX}
\setmainfont{PlusJakartaSans-Medium.ttf}[Path=../../../fonts/,BoldFont=PlusJakartaSans-Bold.ttf]
% Maths en sans-serif (Fira Math) avec chiffres et lettres droites de Plus
% Jakarta, pour que les formules se fondent dans le texte.
\usepackage[math-style=ISO,bold-style=ISO]{unicode-math}
\setmathfont{FiraMath-Regular.otf}[Path=../../../fonts/]
\setmathfont{PlusJakartaSans-Medium.ttf}[Path=../../../fonts/,range={up/num,up/{Latin,latin},\mathup}]
\usepackage{icomma} % virgule décimale sans espace en mode math
% Caractères mathématiques Unicode absents des polices de texte (Plus Jakarta,
% Archivo Black) : rendus par leur équivalent mathématique, en texte comme en
% formule — sinon ils s'affichent en carré vide (ex. « Arithmétique dans ℤ »).
\usepackage{newunicodechar}
\newunicodechar{ℝ}{\ensuremath{\mathbb{R}}}
\newunicodechar{ℤ}{\ensuremath{\mathbb{Z}}}
\newunicodechar{ℂ}{\ensuremath{\mathbb{C}}}
\newunicodechar{ℕ}{\ensuremath{\mathbb{N}}}
\newunicodechar{ℚ}{\ensuremath{\mathbb{Q}}}
\newunicodechar{𝔻}{\ensuremath{\mathbb{D}}}
\newunicodechar{α}{\ensuremath{\alpha}}
\newunicodechar{β}{\ensuremath{\beta}}
\newunicodechar{γ}{\ensuremath{\gamma}}
\newunicodechar{δ}{\ensuremath{\delta}}
\newunicodechar{ε}{\ensuremath{\varepsilon}}
\newunicodechar{θ}{\ensuremath{\theta}}
\newunicodechar{λ}{\ensuremath{\lambda}}
\newunicodechar{μ}{\ensuremath{\mu}}
\newunicodechar{σ}{\ensuremath{\sigma}}
\newunicodechar{τ}{\ensuremath{\tau}}
\newunicodechar{φ}{\ensuremath{\varphi}}
\newunicodechar{ψ}{\ensuremath{\psi}}
\newunicodechar{ω}{\ensuremath{\omega}}
\newunicodechar{Σ}{\ensuremath{\Sigma}}
% majuscules produites par \MakeUppercase dans les titres (α-aminés -> Α)
\newunicodechar{Α}{\ensuremath{\alpha}}
\newunicodechar{Β}{\ensuremath{\beta}}
\newunicodechar{Γ}{\ensuremath{\Gamma}}
\newunicodechar{Δ}{\ensuremath{\Delta}}
\newunicodechar{Ε}{\ensuremath{\varepsilon}}
\newunicodechar{Θ}{\ensuremath{\Theta}}
\newunicodechar{Λ}{\ensuremath{\Lambda}}
\newunicodechar{Μ}{\ensuremath{\mu}}
\newunicodechar{Τ}{\ensuremath{\tau}}
\newunicodechar{Φ}{\ensuremath{\Phi}}
\newunicodechar{Ψ}{\ensuremath{\Psi}}
\newunicodechar{Ω}{\ensuremath{\Omega}}
\newunicodechar{↦}{\ensuremath{\mapsto}}
\newunicodechar{∈}{\ensuremath{\in}}
\newunicodechar{∉}{\ensuremath{\notin}}
\newunicodechar{∧}{\ensuremath{\wedge}}
\newunicodechar{⇔}{\ensuremath{\Leftrightarrow}}
\newunicodechar{⟺}{\ensuremath{\Longleftrightarrow}}
\newunicodechar{⇒}{\ensuremath{\Rightarrow}}
\newunicodechar{⟹}{\ensuremath{\Longrightarrow}}
\newunicodechar{∘}{\ensuremath{\circ}}
\newunicodechar{∪}{\ensuremath{\cup}}
\newunicodechar{∩}{\ensuremath{\cap}}
\newunicodechar{≡}{\ensuremath{\equiv}}
\newunicodechar{⊂}{\ensuremath{\subset}}
\newunicodechar{⊥}{\ensuremath{\perp}}
\newunicodechar{∃}{\ensuremath{\exists}}
\newunicodechar{∀}{\ensuremath{\forall}}
\newunicodechar{∛}{\ensuremath{\sqrt[3]{\,}}}
\newunicodechar{∜}{\ensuremath{\sqrt[4]{\,}}}
\newunicodechar{⃗}{\ensuremath{{}^{\to}}}
\newunicodechar{ⁿ}{\ifmmode^{n}\else\textsuperscript{\ensuremath{n}}\fi}
\newunicodechar{ˣ}{\ifmmode^{x}\else\textsuperscript{\ensuremath{x}}\fi}
\newunicodechar{ᵇ}{\ifmmode^{b}\else\textsuperscript{\ensuremath{b}}\fi}
\newunicodechar{ʳ}{\ifmmode^{r}\else\textsuperscript{\ensuremath{r}}\fi}
\newunicodechar{ᵘ}{\ifmmode^{u}\else\textsuperscript{\ensuremath{u}}\fi}
\newunicodechar{ʸ}{\ifmmode^{y}\else\textsuperscript{\ensuremath{y}}\fi}
\newunicodechar{ᵃ}{\ifmmode^{a}\else\textsuperscript{\ensuremath{a}}\fi}
\newunicodechar{ᵉ}{\ifmmode^{e}\else\textsuperscript{\ensuremath{e}}\fi}
\newunicodechar{ᵏ}{\ifmmode^{k}\else\textsuperscript{\ensuremath{k}}\fi}
\newunicodechar{ᵖ}{\ifmmode^{p}\else\textsuperscript{\ensuremath{p}}\fi}
\newunicodechar{ᴺ}{\ifmmode^{N}\else\textsuperscript{\ensuremath{N}}\fi}
\newunicodechar{ᶜ}{\ifmmode^{c}\else\textsuperscript{\ensuremath{c}}\fi}
\newunicodechar{ʰ}{\ifmmode^{h}\else\textsuperscript{\ensuremath{h}}\fi}
\newunicodechar{ᵠ}{\ifmmode^{q}\else\textsuperscript{\ensuremath{q}}\fi}
\newunicodechar{ₙ}{\ifmmode_{n}\else\textsubscript{\ensuremath{n}}\fi}
\newunicodechar{ₐ}{\ifmmode_{a}\else\textsubscript{\ensuremath{a}}\fi}
\newunicodechar{ᵢ}{\ifmmode_{i}\else\textsubscript{\ensuremath{i}}\fi}
\newunicodechar{ᵣ}{\ifmmode_{r}\else\textsubscript{\ensuremath{r}}\fi}
\newunicodechar{ₖ}{\ifmmode_{k}\else\textsubscript{\ensuremath{k}}\fi}
\newunicodechar{ᵦ}{\ifmmode_{\beta}\else\textsubscript{\ensuremath{\beta}}\fi}
\newunicodechar{ₓ}{\ifmmode_{x}\else\textsubscript{\ensuremath{x}}\fi}
\newfontfamily\popdisplay{ArchivoBlack-Regular.ttf}[Path=../../../fonts/]
\newfontfamily\poplogo{SpaceGrotesk-Bold.ttf}[Path=../../../fonts/]
\newfontfamily\popsemi{PlusJakartaSans-SemiBold.ttf}[Path=../../../fonts/]
\newfontfamily\popxbold{PlusJakartaSans-ExtraBold.ttf}[Path=../../../fonts/]
\newfontfamily\popxboldls{PlusJakartaSans-ExtraBold.ttf}[Path=../../../fonts/,LetterSpace=3.0]

\setlist[enumerate]{leftmargin=1.7em,itemsep=5pt,topsep=4pt}
\setlist[itemize]{leftmargin=1.7em,itemsep=4pt,topsep=4pt,label={\color{poporange}\textbullet}}
\setlength{\parindent}{0pt}
\setlength{\parskip}{5pt}
\renewcommand{\arraystretch}{1.25}

% pgfplots : style Pop par défaut
\pgfplotsset{
  every axis/.append style={axis line style={popink,line width=1pt},tick style={popink},
    grid style={popline,line width=0.5pt},label style={font=\small},tick label style={font=\footnotesize}},
  popcurve/.style={poporange,line width=1.8pt,no markers,smooth},
}
% Même style disponible dans un \draw TikZ simple (hors pgfplots)
\tikzset{popcurve/.style={poporange,line width=1.8pt}}

% ============================================================
% Pill / squiggle
% ============================================================
\newcommand{\poppill}[3]{%
  \tikz[baseline=-0.55ex]\node[draw=popink,line width=1.2pt,rounded corners=2.6pt,%
    fill=#2,inner xsep=6.5pt,inner ysep=3pt]{%
    \popxboldls\fontsize{7}{7}\selectfont\color{#3}\MakeUppercase{#1}};%
}
\newcommand{\popsquiggle}[1]{%
  \tikz[baseline=-0.4ex]{
    \draw[#1,line width=2.6pt,line cap=round]
      plot [smooth] coordinates {(0,0.09) (0.34,0.01) (0.68,0.21) (1.02,0.01) (1.36,0.10)};
  }%
}
% Mot-clé surligné (marqueur orange sous le texte)
\newcommand{\cle}[1]{{\popxbold\color{popdark}#1}}

% ============================================================
% En-tête / pied
% ============================================================
\pagestyle{fancy}
\fancyhf{}
\renewcommand{\headrulewidth}{0pt}
\renewcommand{\footrulewidth}{0pt}
\lhead{\raisebox{-0.15ex}{\poplogo\fontsize{13}{13}\selectfont\color{popink}OnBuch\color{poporange}+}}
\rhead{\poppill{\DOCMATIERE\ \textbullet\ \DOCNIVEAU\ \textbullet\ Leçon \DOCLECON}{white}{popink}}
\lfoot{\popxbold\fontsize{7.6}{7.6}\selectfont\color{popmuted}\MakeUppercase{Cours signé OnBuch+}\ \textbullet\ {\popsemi\DOCTITRE}}
\rfoot{\tikz[baseline=-0.6ex]\node[rounded corners=8.5pt,fill=popink,minimum height=17pt,minimum width=17pt,inner xsep=5pt,inner sep=0]{\popxbold\fontsize{8}{8}\selectfont\color{popcream}\thepage\,/\,\pageref*{LastPage}};}

% ============================================================
% Titres
% ============================================================
\newcommand{\chaptitre}[2]{%
  \begin{tcolorbox}[enhanced,colback=poporange,colframe=popink,boxrule=2.6pt,arc=11pt,
      drop shadow=popink,left=15pt,right=15pt,top=12pt,bottom=11pt,before skip=3pt,after skip=18pt]
    \poppill{#2}{popink}{popcream}\par\vspace{9pt}
    {\raggedright\hyphenpenalty=10000\exhyphenpenalty=10000\popdisplay\fontsize{22}{24}\selectfont\color{popcream}\MakeUppercase{#1}\par}
  \end{tcolorbox}
}
% Section numérotée : pastille orange avec le numéro + titre Archivo
\newcounter{popsec}
\newcommand{\coursec}[1]{%
  \stepcounter{popsec}\setcounter{popsub}{0}%
  \par\vspace{16pt}\noindent
  \begin{minipage}{\linewidth}
  \tikz[baseline=-0.7ex]\node[draw=popink,line width=1.6pt,fill=poporange,rounded corners=4pt,minimum size=22pt,inner sep=0]{\popdisplay\fontsize{12}{12}\selectfont\color{popcream}\thepopsec};%
  \hspace{8pt}{\popdisplay\fontsize{15}{17}\selectfont\color{popink}\MakeUppercase{#1}}\par
  \vspace{4pt}\noindent{\color{poporange}\rule{36pt}{3.6pt}}
  \end{minipage}\par\nopagebreak\vspace{8pt}
}
\newcounter{popsub}
\newcommand{\courssub}[1]{%
  \stepcounter{popsub}%
  \par\vspace{9pt}\noindent{\popxbold\fontsize{11.5}{13}\selectfont\color{popink}{\color{poporange}\thepopsec.\thepopsub}\hspace{6pt}#1}\par\nopagebreak\vspace{4pt}
}

% ============================================================
% Boîtes pédagogiques — même recette : fond blanc, bordure épaisse colorée,
% ombre dure encre, badge plein. Titre optionnel [#1].
% ============================================================
\tcbset{popbase/.style={breakable,enhanced,colback=white,boxrule=2.3pt,arc=9pt,
  shadow={3pt}{-3pt}{0pt}{popink},left=14pt,right=14pt,top=11pt,bottom=11pt,before skip=11pt,after skip=15pt,
  fontupper=\popsemi\fontsize{10.6}{14.5}\selectfont\color{popink},
  fontlower=\popsemi\fontsize{10.6}{14.5}\selectfont\color{popink}}}
% Coche dessinée (le glyphe ✓ n'existe pas dans les polices du document)
\newcommand{\popcheck}[1]{\tikz[baseline=-0.5ex]\draw[#1,line width=1.6pt,line cap=round,line join=round](-0.13,0.0)--(-0.04,-0.09)--(0.13,0.1);}
\newcommand{\popboxhead}[3]{\poppill{#1}{#2}{white}\ifx\relax#3\relax\else\hspace{6pt}{\popxbold\fontsize{10.6}{12}\selectfont\color{popink}#3}\fi\par\vspace{7pt}}

\newtcolorbox{aretenir}[1][]{popbase,colframe=popgreen,before upper={\popboxhead{À retenir}{popgreen}{#1}}}
\newtcolorbox{exemplebox}[1][]{popbase,colframe=poporange,before upper={\popboxhead{Exemple}{poporange}{#1}}}
\newtcolorbox{definition}[1][]{popbase,colframe=popblue,before upper={\popboxhead{Définition}{popblue}{#1}}}
\newtcolorbox{propriete}[1][]{popbase,colframe=poppurple,before upper={\popboxhead{Propriété}{poppurple}{#1}}}
\newtcolorbox{methode}[1][]{popbase,colframe=popgold,before upper={\popboxhead{Méthode}{popgold}{#1}}}
\newtcolorbox{attention}[1][]{popbase,colframe=poppink,before upper={\popboxhead{Attention}{poppink}{#1}}}
\newtcolorbox{experience}[1][]{popbase,colframe=popblue,colback=popblueL,before upper={\popboxhead{Expérience}{popblue}{#1}}}
% « Le savais-tu ? » : carte violette pleine (contexte, culture, Cameroun)
\newtcolorbox{savaistu}[1][]{popbase,colback=poppurple,colframe=popink,
  fontupper=\popsemi\fontsize{10.4}{14.2}\selectfont\color{white},
  before upper={\poppill{Le savais-tu\,?}{white}{poppurple}\ifx\relax#1\relax\else\hspace{6pt}{\popxbold\color{white}#1}\fi\par\vspace{7pt}}}
% Exercice résolu : énoncé en haut, \tcblower puis la solution sur fond crème
\newcounter{popexo}\newcounter{popexr}
\newtcolorbox{exoresolu}[1][]{popbase,colframe=poporange,colbacklower=popcream,
  segmentation style={popink,line width=1.2pt,dashed},
  before upper={\stepcounter{popexr}\popboxhead{Exercice résolu \thepopexr}{poporange}{#1}},
  before lower={\poppill{Solution}{popink}{popcream}\par\vspace{6pt}}}
% Exercice (non résolu en place) — la correction est dans la partie Corrigés
\newtcolorbox{exercice}[1][]{popbase,colframe=popink,
  before upper={\stepcounter{popexo}\popboxhead{Exercice \thepopexo}{popink}{#1}}}
% Corrigé
\newtcolorbox{corrige}[1][]{popbase,colframe=popgreen,colback=popgreenL,
  before upper={\popboxhead{Corrigé}{popgreen}{#1}}}
% Objectifs / prérequis / bilan
\newtcolorbox{objectifs}{popbase,colframe=popink,colback=poporangeL,
  before upper={\popboxhead{Objectifs de la leçon}{popink}{}}}
\newtcolorbox{prerequis}{popbase,colframe=popink,colback=popblueL,
  before upper={\popboxhead{Prérequis}{popblue}{}}}
\newtcolorbox{bilan}{popbase,colframe=popink,colback=white,boxrule=2.8pt,
  before upper={\popboxhead{Fiche bilan}{poporange}{L'essentiel à connaître pour l'examen}}}
% Figure encadrée avec légende
\newtcolorbox{popfigure}[1][]{enhanced,breakable=false,colback=white,colframe=popline,boxrule=1.4pt,arc=8pt,
  left=10pt,right=10pt,top=10pt,bottom=8pt,before skip=10pt,after skip=12pt,halign=center,
  lower separated=false,halign lower=center,
  fontlower=\popsemi\fontsize{9}{11.5}\selectfont\color{popmuted},
  #1}
\newcommand{\legende}[1]{\par\vspace{6pt}{\popsemi\fontsize{9}{11.5}\selectfont\color{popmuted}#1}}

% ============================================================
% Signature OnBuch+
% ============================================================
\newcommand{\poplogoinline}[1]{{\poplogo\fontsize{#1}{#1}\selectfont\color{popink}OnBuch\color{poporange}+}}
\newcommand{\signatureonbuch}{%
  \begin{tcolorbox}[enhanced,colback=popink,colframe=popink,boxrule=2.2pt,arc=10pt,
      drop shadow=poporange,left=14pt,right=14pt,top=10pt,bottom=10pt,before skip=0pt,after skip=0pt]
    \tikz[baseline=-0.7ex]\node[circle,fill=popgreen,draw=popcream,line width=1.3pt,minimum size=22pt,inner sep=0]{\popcheck{white}};%
    \hspace{9pt}\begin{minipage}[c]{0.62\linewidth}
      {\popxbold\fontsize{12}{13}\selectfont\color{popcream}Signé\ \textbullet\ {\poplogo OnBuch\color{poporange}+}}\par\vspace{2pt}
      {\raggedright\popsemi\fontsize{8}{10.5}\selectfont\color{popline}\MakeUppercase{Équipe pédagogique OnBuch+}\\\MakeUppercase{Conforme au programme officiel MINESEC}\par}
    \end{minipage}\hfill
    {\poplogo\fontsize{22}{22}\selectfont\color{popcream}OnBuch\color{poporange}+}
  \end{tcolorbox}%
}

% ============================================================
% Page de garde
% #1 titre · #2 matière · #3 niveau · #4 module · #5 n° leçon · #6 durée ·
% #7 description / objectifs (texte ou itemize)
% ============================================================
\newcommand{\pagegarde}[7]{%
  \thispagestyle{empty}
  \newgeometry{a4paper,margin=1.7cm,top=1.6cm,bottom=1.4cm}%
  \begin{tikzpicture}[remember picture,overlay]
    \fill[poporange!18] ([xshift=-3.2cm,yshift=-3.6cm]current page.north east) circle (4.6cm);
    \fill[poppurple!14] ([xshift=2.4cm,yshift=7.6cm]current page.south west) circle (3.8cm);
  \end{tikzpicture}%
  {\poplogo\fontsize{30}{30}\selectfont\color{popink}OnBuch\color{poporange}+}\hfill
  \raisebox{0.6ex}{\poppill{Cours signé \textbullet\ Premium}{popink}{popcream}}\par
  \vspace{1pt}\popsquiggle{poppurple}\par
  \vspace{8pt}
  \begin{tcolorbox}[enhanced,colback=poporange,colframe=popink,boxrule=2.8pt,arc=13pt,
      drop shadow=popink,left=18pt,right=18pt,top=12pt,bottom=13pt,after skip=10pt]
    \poppill{#2 \textbullet\ #3}{popink}{popcream}\hspace{5pt}\poppill{#4}{white}{popink}\par\vspace{8pt}
    {\popdisplay\fontsize{14}{14}\selectfont\color{popink}LEÇON #5}\par\vspace{4pt}
    {\raggedright\hyphenpenalty=10000\exhyphenpenalty=10000\popdisplay\fontsize{25}{27}\selectfont\color{popcream}\MakeUppercase{#1}\par}
    \vspace{8pt}
    {\popxbold\fontsize{9.5}{9.5}\selectfont\color{popink}\MakeUppercase{Durée conseillée : #6}}
  \end{tcolorbox}
  \begin{tcolorbox}[enhanced,colback=white,colframe=popink,boxrule=2.2pt,arc=10pt,
      drop shadow=popink,left=16pt,right=16pt,top=10pt,bottom=10pt,after skip=8pt]
    \poppill{Dans cette leçon}{poppurple}{white}\par\vspace{5pt}
    \popsemi\fontsize{9.3}{12.6}\selectfont\color{popink}#7
  \end{tcolorbox}
  \noindent\begin{minipage}[t]{0.487\linewidth}
    \begin{tcolorbox}[enhanced,colback=popgreen,colframe=popink,boxrule=2.2pt,arc=10pt,
        drop shadow=popink,left=11pt,right=11pt,top=10pt,bottom=10pt,height=3.1cm,valign=center]
      \tcbox[colback=white,colframe=popink,boxrule=1.3pt,arc=5pt,boxsep=0pt,nobeforeafter,
        left=3pt,right=3pt,top=3pt,bottom=3pt,valign=center]{\qrcode[height=1.9cm]{https://play.google.com/store/apps/details?id=com.luvvix.onbuch&hl=fr}}%
      \hspace{8pt}\begin{minipage}[c]{0.5\linewidth}
        {\popdisplay\fontsize{11}{12.5}\selectfont\color{white}\MakeUppercase{Télécharge OnBuch}}\par\vspace{4pt}
        {\raggedright\popsemi\fontsize{8.4}{11}\selectfont\color{white}Gratuit \textbullet\ Google Play\\Tuteur IA, annales, concours, résultats\par}
      \end{minipage}
    \end{tcolorbox}
  \end{minipage}\hfill
  \begin{minipage}[t]{0.487\linewidth}
    \begin{tcolorbox}[enhanced,colback=poppurple,colframe=popink,boxrule=2.2pt,arc=10pt,
        drop shadow=popink,left=11pt,right=11pt,top=10pt,bottom=10pt,height=3.1cm,valign=center]
      \tikz[baseline=-0.7ex]\node[circle,fill=white,draw=popink,line width=1.3pt,minimum size=1.6cm,inner sep=0]{\popdisplay\fontsize{24}{24}\selectfont\color{poppurple}+};%
      \hspace{8pt}\begin{minipage}[c]{0.6\linewidth}
        {\popdisplay\fontsize{11}{12.5}\selectfont\color{white}\MakeUppercase{Passe à OnBuch+}}\par\vspace{4pt}
        {\raggedright\popsemi\fontsize{8.4}{11}\selectfont\color{white}Tous les cours signés, Tuteur IA illimité, fiches PDF — onglet OnBuch+ de l'app\par}
      \end{minipage}
    \end{tcolorbox}
  \end{minipage}
  \vspace{8pt}
  \popcontenu
  \vfill
  \signatureonbuch
  \restoregeometry
  \newpage
}

% Rangée « Ce cours contient » (page de garde)
\newcommand{\poptile}[3]{% #1 couleur #2 gros texte #3 légende
  \begin{tcolorbox}[enhanced,colback=#1,colframe=popink,boxrule=2pt,arc=9pt,shadow={2.5pt}{-2.5pt}{0pt}{popink},
      left=6pt,right=6pt,top=9pt,bottom=9pt,halign=center,valign=center,height=1.75cm,width=\linewidth,nobeforeafter]
    {\popdisplay\fontsize{12.5}{13}\selectfont\color{white}\MakeUppercase{#2}}\par\vspace{3pt}
    {\centering\popsemi\fontsize{7.8}{9.5}\selectfont\color{white}#3\par}
  \end{tcolorbox}}
\newcommand{\popcontenu}{%
  \par\noindent{\popxboldls\fontsize{7.5}{7.5}\selectfont\color{popmuted}CE COURS SIGNÉ CONTIENT}\par\vspace{7pt}
  \noindent\begin{minipage}[t]{0.235\linewidth}\poptile{popblue}{Cours}{complet, illustré, pas à pas}\end{minipage}\hfill
  \begin{minipage}[t]{0.235\linewidth}\poptile{poporange}{Méthodes}{et exercices résolus}\end{minipage}\hfill
  \begin{minipage}[t]{0.235\linewidth}\poptile{poppink}{Exercices}{type examen + corrigés}\end{minipage}\hfill
  \begin{minipage}[t]{0.235\linewidth}\poptile{popgreen}{Fiche bilan}{l'essentiel à retenir}\end{minipage}\par}

% Sommaire
\newtcolorbox{sommaire}{enhanced,colback=white,colframe=popink,boxrule=2.2pt,arc=10pt,
  drop shadow=popink,left=18pt,right=18pt,top=13pt,bottom=13pt,before skip=6pt,after skip=20pt,
  before upper={\poppill{Sommaire}{poppurple}{white}\par\vspace{9pt}\popsemi\fontsize{10.6}{14.5}\selectfont\color{popink}}}

% Signature de fin de cours — poussée tout en bas de la dernière page
\newcommand{\findecours}{%
  \par\vfill\nopagebreak
  \begin{center}\popsquiggle{poporange}\end{center}\vspace{4pt}
  \signatureonbuch
}
\AtBeginDocument{\def\llbracket{\mathopen{[\mkern-3.5mu[}}\def\rrbracket{\mathclose{]\mkern-3.5mu]}}} % intervalles d'entiers (glyphe ⟦ absent de la police)
% Glyphes absents de Fira Math : redéfinis à partir de glyphes disponibles
\AtBeginDocument{%
  \def\setminus{\mathbin{\backslash}}\let\smallsetminus\setminus
  \def\vdots{\mathord{\vbox{\baselineskip3.2pt\lineskiplimit0pt\kern4pt\hbox{.}\hbox{.}\hbox{.}}}}%
  \def\ddots{\mathinner{\mkern1mu\raise6.5pt\hbox{.}\mkern2mu\raise3.6pt\hbox{.}\mkern2mu\raise0.7pt\hbox{.}\mkern1mu}}%
  \def\triangle{\mathord{\scalebox{0.9}{$\symup{\Delta}$}}}%
  \def\nabla{\mathord{\raisebox{\depth}{\scalebox{1}[-1]{$\symup{\Delta}$}}}}%
  \def\checkmark{\popcheck{popdark}}%
  \def\star{\mathbin{\ast}}\def\bigstar{\mathord{\ast}}%
  \def\diamond{\mathbin{\scalebox{0.6}{$\Diamond$}}}%
  \def\bigcup{\mathop{\mathchoice{\raisebox{-0.3em}{\scalebox{1.7}{$\cup$}}}{\scalebox{1.25}{$\cup$}}{\cup}{\cup}}\displaylimits}%
  \def\bigcap{\mathop{\mathchoice{\raisebox{-0.3em}{\scalebox{1.7}{$\cap$}}}{\scalebox{1.25}{$\cap$}}{\cap}{\cap}}\displaylimits}%
  \def\lesseqqgtr{\mathrel{\lessgtr}}\def\bigoplus{\mathop{\oplus}\displaylimits}%
}
\providecommand{\ssi}{si et seulement si} % abréviation fréquente
\AtBeginDocument{\providecommand{\wideparen}{\overparen}} % arc de cercle

%% FIGFIX-BEGIN v5 — correctifs globaux des figures (docs/onbuch-generation/tools/figfix)
\usepackage{adjustbox}\usepackage{etoolbox}\tcbuselibrary{raster}
% toute figure TikZ/pgfplots plus large que la ligne est réduite à la largeur disponible
\BeforeBeginEnvironment{tikzpicture}{\begin{adjustbox}{max width=\linewidth}}
\AfterEndEnvironment{tikzpicture}{\end{adjustbox}}
% légendes pgfplots lisibles par-dessus les courbes
\pgfplotsset{every axis/.append style={legend style={fill=white,fill opacity=0.92,text opacity=1,draw=black!15,inner sep=3pt}}}
% étiquettes d'arêtes (syntaxe edge["..."]) avec fond blanc
\tikzset{every edge quotes/.append style={fill=white,inner sep=1.2pt}}
%% FIGFIX-END
\begin{document}

\pagegarde{Systèmes d'équations linéaires à deux ou trois inconnues}{Mathématiques}{1ère E}{Relations et opérations fondamentales dans l'ensemble des nombres réels}{N02}{6 h}{Cette leçon présente les systèmes d'équations linéaires à deux ou trois inconnues et les méthodes de résolution : substitution, combinaisons linéaires, déterminant et pivot de Gauss. Elle montre aussi comment discuter un système dépendant d'un paramètre et comment modéliser des problèmes concrets de la vie courante.
\vspace{4pt}
\begin{multicols}{2}\small
\begin{itemize}
\item À la fin de cette leçon, je sais reconnaître un système d'équations linéaires à deux ou trois inconnues et interpréter géométriquement ses solutions.
\item À la fin de cette leçon, je sais résoudre un système linéaire dans ℝ² par la méthode de substitution.
\item À la fin de cette leçon, je sais résoudre un système linéaire dans ℝ² par la méthode des combinaisons linéaires.
\item À la fin de cette leçon, je sais calculer le déterminant d'un système 2×2 et en déduire l'existence et l'unicité de la solution, puis la donner par les formules de Cramer.
\item À la fin de cette leçon, je sais résoudre et discuter un système dans ℝ² dont le second membre dépend d'un paramètre.
\item À la fin de cette leçon, je sais résoudre un système linéaire dans ℝ³ par substitution, par combinaisons linéaires ou par le pivot de Gauss.
\end{itemize}\end{multicols}}

\begin{sommaire}
\begin{multicols}{2}
\begin{enumerate}
\item Systèmes linéaires dans ℝ² : définitions et interprétation géométrique
\item Résolution dans ℝ² : substitution et combinaisons linéaires
\item Déterminant d'un système 2×2 et formules de Cramer
\item Systèmes dans ℝ² à second membre dépendant d'un paramètre
\item Systèmes linéaires dans ℝ³ : substitution, combinaisons et pivot de Gauss
\item Résolution de problèmes concrets par les systèmes
\item Activité d'intégration
\item Méthodes et exercices résolus
\item Exercices
\item Corrigés des exercices
\item Fiche bilan
\end{enumerate}
\end{multicols}
\end{sommaire}

\begin{objectifs}
\begin{itemize}
\item À la fin de cette leçon, je sais reconnaître un système d'équations linéaires à deux ou trois inconnues et interpréter géométriquement ses solutions.
\item À la fin de cette leçon, je sais résoudre un système linéaire dans ℝ² par la méthode de substitution.
\item À la fin de cette leçon, je sais résoudre un système linéaire dans ℝ² par la méthode des combinaisons linéaires.
\item À la fin de cette leçon, je sais calculer le déterminant d'un système 2×2 et en déduire l'existence et l'unicité de la solution, puis la donner par les formules de Cramer.
\item À la fin de cette leçon, je sais résoudre et discuter un système dans ℝ² dont le second membre dépend d'un paramètre.
\item À la fin de cette leçon, je sais résoudre un système linéaire dans ℝ³ par substitution, par combinaisons linéaires ou par le pivot de Gauss.
\item À la fin de cette leçon, je sais traduire un problème concret de la vie courante en système d'équations et interpréter la solution obtenue.
\end{itemize}
\end{objectifs}

\begin{prerequis}
\begin{itemize}
\item Calcul algébrique dans ℝ : développement, factorisation, résolution d'équations du premier degré à une inconnue (classe de Seconde).
\item Équation cartésienne d'une droite du plan : ax + by + c = 0 (Seconde).
\item Notion de couple et de triplet de nombres réels ; repérage dans le plan.
\item Résolution d'inéquations simples et manipulation de paramètres littéraux.
\end{itemize}
\end{prerequis}

\newpage

\coursec{Systèmes linéaires dans $\mathbb{R}^2$ : définitions et interprétation géométrique}

\textbf{Situation d'accroche.} Il est 6 heures du matin au marché Mokolo, à Yaoundé. Mama Ngo Bell décharge ses paniers devant son étal. Elle propose à ses clients deux formules bien rodées : le \emph{panier A}, qui contient 3 kg de manioc et 2 kg de plantain, et le \emph{panier B}, qui contient 1 kg de manioc et 4 kg de plantain. Ce matin-là, un restaurateur de Douala, joint par téléphone, lui passe une commande précise : il lui faut exactement 26 kg de manioc et 28 kg de plantain pour la semaine. Mama Ngo Bell sort son cahier et se demande : combien de paniers A et combien de paniers B dois-je préparer pour satisfaire exactement cette commande ?

Notons $x$ le nombre de paniers A et $y$ le nombre de paniers B. Le manioc fourni s'écrit $3x + y$ (3 kg par panier A, 1 kg par panier B), et le plantain fourni s'écrit $2x + 4y$. La commande impose donc :
\[
\begin{cases} 3x + y = 26 \\ 2x + 4y = 28 \end{cases}
\]
Deux conditions à satisfaire \emph{en même temps} : voilà l'idée d'un \cle{système d'équations}. Ce problème, comme d'innombrables situations de la vie courante (mélanges de produits, partages d'héritage, tarifs de transport, calculs de moyennes, forfaits téléphoniques), se traduit par un système d'équations linéaires. Cette leçon t'apprendra à résoudre de tels systèmes de plusieurs façons. Mais avant de résoudre, il faut comprendre \emph{ce qu'est} un système, \emph{ce que signifie} le mot « solution », et \emph{combien} de solutions on peut espérer. C'est l'objet de cette première section.

\courssub{Équation linéaire à deux inconnues}

\textbf{Intuition.} Une équation est une « balance » : ce qui est à gauche du signe $=$ doit égaler ce qui est à droite. Quand l'inconnue est le prix d'un article, on a une équation à une inconnue. Mais dans la situation de Mama Ngo Bell, \emph{deux} quantités sont inconnues : le nombre de paniers A et le nombre de paniers B. Il nous faut donc des équations à deux inconnues.

\begin{definition}[Équation linéaire à deux inconnues]
Soient $a$, $b$ et $c$ trois réels donnés, avec $(a, b) \neq (0, 0)$. On appelle \cle{équation linéaire à deux inconnues} $x$ et $y$ toute équation qui peut s'écrire sous la forme
\[
ax + by = c.
\]
Les réels $a$ et $b$ sont les \cle{coefficients} des inconnues, $c$ est le \cle{second membre}. Une \cle{solution} de cette équation est un couple $(x, y)$ de réels qui rend l'égalité vraie.
\end{definition}

\begin{exemplebox}[Exemples et contre-exemples]
\begin{itemize}
\item $3x + y = 26$ est une équation linéaire ($a = 3$, $b = 1$, $c = 26$). Le couple $(8, 2)$ est solution car $3 \times 8 + 2 = 26$. Le couple $(5, 11)$ est aussi solution car $3 \times 5 + 11 = 26$. Une équation à deux inconnues possède donc en général \emph{une infinité} de solutions.
\item $2x - 5y = 7$ est linéaire : elle s'écrit $2x + (-5)y = 7$.
\item $x^2 + y = 3$ n'est \textbf{pas} linéaire (à cause de $x^2$), pas plus que $xy = 4$ (produit des inconnues) ni $\dfrac{1}{x} + y = 2$.
\end{itemize}
\end{exemplebox}

\begin{attention}[Pourquoi exiger $(a,b) \neq (0,0)$ ?]
Si $a = 0$ et $b = 0$, l'équation devient $0x + 0y = c$, c'est-à-dire $0 = c$. Deux choses peuvent arriver :
\begin{itemize}
\item si $c \neq 0$ (par exemple $0x + 0y = 5$), l'égalité est \textbf{toujours fausse} : aucun couple $(x,y)$ n'est solution ;
\item si $c = 0$ (c'est-à-dire $0x + 0y = 0$), l'égalité est \textbf{toujours vraie} : \textbf{tous} les couples $(x,y)$ sont solutions.
\end{itemize}
Dans les deux cas, l'équation ne donne aucune information sur $x$ et $y$ : ce n'est plus une véritable équation aux inconnues $x$ et $y$. C'est pourquoi on exclut ce cas de la définition. Ne confonds jamais « aucune solution » avec « tous les réels sont solutions » !
\end{attention}

\textbf{Lien avec la géométrie.} Puisque $(a, b) \neq (0, 0)$, on peut toujours isoler l'une des inconnues. Si $b \neq 0$, alors $ax + by = c$ équivaut à $y = -\dfrac{a}{b}x + \dfrac{c}{b}$ : c'est l'équation réduite d'une \cle{droite} du plan, de coefficient directeur $-\dfrac{a}{b}$. Si $b = 0$ (alors $a \neq 0$), l'équation devient $x = \dfrac{c}{a}$ : c'est une droite \emph{verticale}, parallèle à l'axe des ordonnées. Retenons ce fait fondamental.

\begin{propriete}[Interprétation géométrique d'une équation linéaire]
Dans un repère du plan, l'ensemble des couples $(x, y)$ solutions de l'équation $ax + by = c$ (avec $(a,b) \neq (0,0)$) est exactement l'ensemble des coordonnées des points d'une \cle{droite}. On dit que $ax + by = c$ est une \cle{équation cartésienne} de cette droite.
\end{propriete}

Ainsi, dire que $(8, 2)$ est solution de $3x + y = 26$, c'est dire que le point de coordonnées $(8\,;\,2)$ appartient à la droite d'équation $3x + y = 26$.

\courssub{Système de deux équations linéaires à deux inconnues}

\begin{definition}[Système linéaire $2 \times 2$]
Soient $a$, $b$, $c$, $a'$, $b'$, $c'$ six réels donnés, avec $(a,b) \neq (0,0)$ et $(a',b') \neq (0,0)$. On appelle \cle{système de deux équations linéaires à deux inconnues} tout système de la forme
\[
(S) \quad \begin{cases} ax + by = c \\ a'x + b'y = c' \end{cases}
\]
Une \cle{solution} du système $(S)$ est un couple $(x, y)$ qui vérifie \emph{simultanément} les deux équations. L'ensemble de toutes les solutions est noté $\mathcal{S}$. \cle{Résoudre} $(S)$, c'est déterminer $\mathcal{S}$.
\end{definition}

Le mot \emph{simultanément} est le cœur de la définition : une solution doit contenter les deux équations \emph{à la fois}, comme le restaurateur de Douala exige à la fois ses 26 kg de manioc \emph{et} ses 28 kg de plantain.

\begin{exemplebox}[Vérifier qu'un couple est solution]
Vérifions que le couple $(4\,;\,7)$ est solution du système de Mama Ngo Bell :
\[
\begin{cases} 3x + y = 26 \\ 2x + 4y = 28 \end{cases}
\]
Pour $x = 4$ et $y = 7$ :
\begin{align*}
3x + y &= 3 \times 4 + 7 = 12 + 7 = 19 \neq 26.
\end{align*}
Le couple $(4\,;\,7)$ ne vérifie pas la première équation : ce n'est \textbf{pas} une solution. Essayons $(8\,;\,2)$ :
\begin{align*}
3x + y &= 3 \times 8 + 2 = 26 \quad \checkmark \\
2x + 4y &= 2 \times 8 + 4 \times 2 = 16 + 8 = 24 \neq 28.
\end{align*}
Le couple $(8\,;\,2)$ vérifie la première équation mais pas la seconde : ce n'est toujours pas une solution. On verra plus tard que la solution exacte est $(x\,;\,y) = (7{,}6\,;\,3{,}2)$ — Mama Ngo Bell devra préparer des fractions de paniers, ou ajuster sa commande ! L'essentiel est de retenir : \emph{un couple n'est solution que s'il vérifie TOUTES les équations}.
\end{exemplebox}

\begin{attention}[Pièges de rédaction]
\begin{itemize}
\item Une solution d'un système à deux inconnues est un \textbf{couple} $(x\,;\,y)$, pas deux nombres séparés. On écrit $\mathcal{S} = \{(8\,;\,2)\}$, et non « $x = 8$ ou $y = 2$ ».
\item Vérifier une seule des deux équations ne suffit jamais à conclure.
\item L'ordre dans le couple compte : $(2\,;\,8)$ n'est pas la même chose que $(8\,;\,2)$.
\end{itemize}
\end{attention}

\courssub{Interprétation géométrique : l'intersection de deux droites}

Voici l'idée qui éclaire toute la théorie. Plaçons-nous dans un repère du plan. La première équation $ax + by = c$ représente une droite $(d)$ ; la seconde $a'x + b'y = c'$ représente une droite $(d')$. Dire qu'un couple $(x, y)$ vérifie les deux équations simultanément, c'est dire que le point $M(x\,;\,y)$ appartient à $(d)$ \emph{et} à $(d')$, c'est-à-dire à leur \cle{intersection}. Par conséquent :

\begin{aretenir}[Traduction géométrique de la résolution]
Résoudre le système $\begin{cases} ax + by = c \\ a'x + b'y = c' \end{cases}$ revient à déterminer les coordonnées des points d'intersection des droites $(d) : ax + by = c$ et $(d') : a'x + b'y = c'$.
\end{aretenir}

Or, en géométrie plane, deux droites n'ont que trois positions relatives possibles. Cela donne immédiatement le théorème central de cette section.

\begin{propriete}[Les trois cas possibles — admis, justifié géométriquement]
Un système de deux équations linéaires à deux inconnues possède :
\begin{itemize}
\item soit \textbf{une unique solution} : les droites $(d)$ et $(d')$ sont \cle{sécantes}, elles se coupent en un seul point ;
\item soit \textbf{aucune solution} : les droites sont \cle{strictement parallèles} (distinctes et sans point commun) ;
\item soit \textbf{une infinité de solutions} : les droites sont \cle{confondues}, tous leurs points sont communs.
\end{itemize}
Il n'y a pas d'autre cas possible. Cette propriété est admise ; elle se justifie par la classification des positions relatives de deux droites du plan.
\end{propriete}

\begin{popfigure}
\begin{tikzpicture}
\begin{axis}[width=0.32\linewidth,scale only axis,xmin=-1,xmax=6,ymin=-1,ymax=6,axis lines=middle,xtick=\empty,ytick=\empty,title={\small Sécantes : 1 solution},title style={color=popdark,font=\small}]
\addplot[popcurve,domain=-1:6,samples=50]{-x+5};
\addplot[popcurve,poporange,domain=-1:6,samples=50]{0.5*x+1};
\addplot[only marks,mark=*,mark size=2pt,popink] coordinates {(2.666,2.333)};
\node[popink,font=\scriptsize,anchor=south west] at (axis cs:2.8,2.5) {$(x_0;y_0)$};
\end{axis}
\end{tikzpicture}\hfill
\begin{tikzpicture}
\begin{axis}[width=0.32\linewidth,scale only axis,xmin=-1,xmax=6,ymin=-1,ymax=6,axis lines=middle,xtick=\empty,ytick=\empty,title={\small Parallèles : 0 solution},title style={color=popdark,font=\small}]
\addplot[popcurve,domain=-1:6,samples=50]{-x+5};
\addplot[popcurve,poporange,domain=-1:6,samples=50]{-x+2};
\end{axis}
\end{tikzpicture}\hfill
\begin{tikzpicture}
\begin{axis}[width=0.32\linewidth,scale only axis,xmin=-1,xmax=6,ymin=-1,ymax=6,axis lines=middle,xtick=\empty,ytick=\empty,title={\small Confondues : infinité},title style={color=popdark,font=\small}]
\addplot[popcurve,domain=-1:6,samples=50]{-x+4};
\addplot[popcurve,poporange,dashed,domain=-1:6,samples=50]{-x+4};
\end{axis}
\end{tikzpicture}
\legende{Les trois positions relatives de deux droites et le nombre de solutions correspondant : droites sécantes (solution unique), droites strictement parallèles (aucune solution), droites confondues (infinité de solutions).}
\end{popfigure}

\textbf{Illustration sur la situation d'accroche.} Les deux équations du système de Mama Ngo Bell s'écrivent $y = -3x + 26$ et $y = -\dfrac{1}{2}x + 7$. Ce sont deux droites de coefficients directeurs différents ($-3$ et $-\tfrac12$) : elles sont sécantes, et leur unique point d'intersection donne l'unique solution du système.

\begin{popfigure}
\begin{tikzpicture}
\begin{axis}[width=0.8\linewidth,scale only axis,xmin=0,xmax=10,ymin=0,ymax=28,axis lines=left,xlabel={$x$ (paniers A)},ylabel={$y$ (paniers B)},xlabel style={font=\small},ylabel style={font=\small},grid=both,grid style={popline!40},legend style={font=\scriptsize,at={(0.98,0.98)},anchor=north east}]
\addplot[popcurve,domain=0:9,samples=50]{-3*x+26};
\addlegendentry{$3x+y=26$}
\addplot[popcurve,poporange,domain=0:10,samples=50]{-0.5*x+7};
\addlegendentry{$2x+4y=28$}
\addplot[only marks,mark=*,mark size=2.5pt,popink] coordinates {(7.6,3.2)};
\draw[dashed,popmuted] (axis cs:7.6,0) -- (axis cs:7.6,3.2) -- (axis cs:0,3.2);
\node[popink,font=\small,anchor=south west] at (axis cs:7.6,3.4) {Solution unique $(7{,}6\,;\,3{,}2)$};
\end{axis}
\end{tikzpicture}
\legende{Le point d'intersection des deux droites est l'unique solution du système $\begin{cases} 3x + y = 26 \\ 2x + 4y = 28 \end{cases}$ : ses coordonnées vérifient les deux équations simultanément.}
\end{popfigure}

\courssub{Lien avec les coefficients : comparaison des rapports}

Comment prévoir le nombre de solutions \emph{sans dessiner} les droites ? Supposons d'abord $b \neq 0$ et $b' \neq 0$. Les équations réduites sont
\[
(d) : y = -\frac{a}{b}x + \frac{c}{b} \qquad \text{et} \qquad (d') : y = -\frac{a'}{b'}x + \frac{c'}{b'}.
\]
Deux droites (non verticales) sont parallèles ou confondues si et seulement si elles ont le \emph{même coefficient directeur}, c'est-à-dire $-\dfrac{a}{b} = -\dfrac{a'}{b'}$, ce qui équivaut à $ab' - a'b = 0$, ou encore, lorsque $a'$ et $b'$ sont non nuls, à $\dfrac{a}{a'} = \dfrac{b}{b'}$. Elles sont confondues si, de plus, elles ont la même ordonnée à l'origine, c'est-à-dire $\dfrac{c}{b} = \dfrac{c'}{b'}$, soit $\dfrac{c}{c'} = \dfrac{b}{b'}$. D'où le critère pratique.

\begin{propriete}[Comparaison des rapports des coefficients]
Soit le système $\begin{cases} ax + by = c \\ a'x + b'y = c' \end{cases}$ où $a'$, $b'$, $c'$ sont tous non nuls.
\begin{itemize}
\item Si $\dfrac{a}{a'} \neq \dfrac{b}{b'}$ : les droites sont \textbf{sécantes}, le système a une \textbf{unique solution}.
\item Si $\dfrac{a}{a'} = \dfrac{b}{b'} \neq \dfrac{c}{c'}$ : les droites sont \textbf{strictement parallèles}, le système n'a \textbf{aucune solution}.
\item Si $\dfrac{a}{a'} = \dfrac{b}{b'} = \dfrac{c}{c'}$ : les droites sont \textbf{confondues}, le système a une \textbf{infinité de solutions}.
\end{itemize}
\end{propriete}

\begin{attention}[Précautions quand un coefficient est nul]
La comparaison des rapports exige que les dénominateurs soient non nuls ! Si $a' = 0$, le rapport $\dfrac{a}{a'}$ n'existe pas. Dans ce cas, raisonne directement : par exemple, pour $\begin{cases} 2x + 3y = 5 \\ 4x + 6y = 9 \end{cases}$, on remarque que la seconde équation est le double de la première... sauf au second membre ($10 \neq 9$) : droites strictement parallèles, aucune solution. Le réflexe sûr : deux équations dont les premiers membres sont proportionnels mais pas les seconds membres sont incompatibles ; si les seconds membres sont proportionnels avec le même coefficient, les équations sont équivalentes (droites confondues).
\end{attention}

\begin{methode}[Reconnaître le nombre de solutions d'un système $2\times 2$]
\begin{enumerate}
\item Écrire le système sous la forme $\begin{cases} ax + by = c \\ a'x + b'y = c' \end{cases}$ et identifier les six coefficients.
\item Si possible, comparer les rapports $\dfrac{a}{a'}$, $\dfrac{b}{b'}$, $\dfrac{c}{c'}$ (sinon, chercher une proportionnalité directe entre les équations).
\item Conclure : rapports $\dfrac{a}{a'} \neq \dfrac{b}{b'}$ : unique solution ; $\dfrac{a}{a'} = \dfrac{b}{b'} \neq \dfrac{c}{c'}$ : aucune solution ; $\dfrac{a}{a'} = \dfrac{b}{b'} = \dfrac{c}{c'}$ : infinité de solutions.
\item (Contrôle) On peut aussi écrire les équations réduites $y = mx + p$ et comparer les coefficients directeurs et les ordonnées à l'origine.
\end{enumerate}
\end{methode}

\begin{exoresolu}[Nombre de solutions sans résolution]
Pour chacun des systèmes suivants, déterminer le nombre de solutions sans le résoudre :
\[
(S_1) \begin{cases} 3x + y = 26 \\ 2x + 4y = 28 \end{cases} \qquad
(S_2) \begin{cases} x - 2y = 3 \\ 3x - 6y = 5 \end{cases} \qquad
(S_3) \begin{cases} 2x + 5y = 10 \\ 4x + 10y = 20 \end{cases}
\]
\tcblower
\textbf{Solution détaillée.}
\begin{itemize}
\item Pour $(S_1)$ : $\dfrac{3}{2} \neq \dfrac{1}{4}$, donc $\dfrac{a}{a'} \neq \dfrac{b}{b'}$ : les droites sont sécantes, $(S_1)$ admet une \textbf{unique solution} (c'est le système de Mama Ngo Bell).
\item Pour $(S_2)$ : $\dfrac{1}{3} = \dfrac{-2}{-6} = \dfrac{1}{3}$, mais $\dfrac{3}{5} \neq \dfrac{1}{3}$. Donc $\dfrac{a}{a'} = \dfrac{b}{b'} \neq \dfrac{c}{c'}$ : les droites sont strictement parallèles, $(S_2)$ n'a \textbf{aucune solution}. Autre lecture : le premier membre de la seconde équation est le triple du premier, mais $5 \neq 3 \times 3 = 9$.
\item Pour $(S_3)$ : $\dfrac{2}{4} = \dfrac{5}{10} = \dfrac{10}{20} = \dfrac{1}{2}$. Les trois rapports sont égaux : la seconde équation est le double de la première, les droites sont confondues, $(S_3)$ admet une \textbf{infinité de solutions} : tous les couples $(x\,;\,y)$ tels que $2x + 5y = 10$, par exemple $(0\,;\,2)$, $(5\,;\,0)$, $(10\,;\,-2)$, etc.
\end{itemize}
\end{exoresolu}

\courssub{Systèmes équivalents et opérations élémentaires}

Pour résoudre un système, la stratégie générale consiste à le \emph{transformer} en un système plus simple, dont on lit directement les solutions, sans perdre ni créer de solution. Encore faut-il savoir quelles transformations sont autorisées.

\begin{definition}[Systèmes équivalents]
Deux systèmes sont dits \cle{équivalents} lorsqu'ils ont exactement le même ensemble de solutions $\mathcal{S}$.
\end{definition}

\begin{propriete}[Opérations élémentaires — admises]
Les transformations suivantes transforment un système en un système équivalent :
\begin{enumerate}
\item \textbf{Échanger} les deux équations ;
\item \textbf{Multiplier} une équation par un réel \emph{non nul} : remplacer $(E)$ par $k(E)$ avec $k \neq 0$ ;
\item \textbf{Ajouter} à une équation un multiple de l'autre : remplacer $(E_1)$ par $(E_1) + \lambda (E_2)$, où $\lambda$ est un réel quelconque, en conservant $(E_2)$ inchangée.
\end{enumerate}
Ces opérations sont la base des méthodes de résolution (combinaisons linéaires, pivot de Gauss) étudiées dans les sections suivantes.
\end{propriete}

\textbf{Pourquoi ces opérations conservent-elles les solutions ?} Pour la multiplication par $k \neq 0$ : l'égalité $ax + by = c$ est vraie si et seulement si $k(ax + by) = kc$ est vraie (on peut « revenir en arrière » en divisant par $k$, car $k \neq 0$). Pour l'ajout : si un couple vérifie $(E_1)$ et $(E_2)$, alors en ajoutant membre à membre $\lambda$ fois l'égalité $(E_2)$ à l'égalité $(E_1)$, il vérifie $(E_1) + \lambda(E_2)$ ; et réciproquement, un couple vérifiant $(E_1) + \lambda(E_2)$ et $(E_2)$ vérifie aussi $(E_1)$, en retranchant $\lambda(E_2)$. Les solutions sont donc exactement les mêmes.

\begin{attention}[Multiplier par zéro est interdit]
Multiplier une équation par $0$ donne $0 = 0$, qui est toujours vraie : on \emph{perd} toute l'information contenue dans l'équation, et le nouveau système peut avoir plus de solutions que l'ancien. De même, dans l'opération $(E_1) \leftarrow (E_1) + \lambda(E_2)$, il faut \textbf{conserver} $(E_2)$ dans le système : remplacer les deux équations à la fois peut détruire l'équivalence.
\end{attention}

\begin{exemplebox}[Une chaîne de systèmes équivalents]
Partons du système de Mama Ngo Bell et appliquons des opérations élémentaires :
\[
\begin{cases} 3x + y = 26 \\ 2x + 4y = 28 \end{cases}
\xrightarrow[\;k=\tfrac12\;]{(E_2) \leftarrow \frac{1}{2}(E_2)}
\begin{cases} 3x + y = 26 \\ x + 2y = 14 \end{cases}
\xrightarrow[\;\lambda=-3\;]{(E_1) \leftarrow (E_1) - 3(E_2)}
\begin{cases} -5y = -16 \\ x + 2y = 14 \end{cases}
\]
Le dernier système, équivalent au premier, se lit presque tout seul : la première équation donne $y = \dfrac{16}{5} = 3{,}2$, puis la seconde donne $x = 14 - 2 \times 3{,}2 = 7{,}6$. On retrouve la solution unique $(7{,}6\,;\,3{,}2)$ annoncée par l'interprétation géométrique. Cette stratégie — transformer pour simplifier — sera développée méthodiquement dans les sections 2 et 5.
\end{exemplebox}

\begin{savaistu}[Des systèmes vieux de deux mille ans]
Les systèmes d'équations linéaires ne sont pas une invention moderne ! Le traité chinois \emph{Jiuzhang Suanshu} (« Les Neuf Chapitres sur l'art du calcul », vers 200 av. J.-C.) résout déjà des systèmes à trois inconnues à l'aide de tableaux de nombres disposés sur une surface à calculer — une méthode très proche de ce que nous appellerons le « pivot de Gauss », dix-huit siècles avant Gauss ! Au Cameroun, les systèmes linéaires sont aujourd'hui utilisés en économie (équilibre de l'offre et de la demande), en agriculture (rations alimentaires du bétail), et dans les réseaux de télécommunications pour l'optimisation des ressources.
\end{savaistu}

\begin{exercice}[À toi de jouer]
\begin{enumerate}
\item Le couple $(3\,;\,-1)$ est-il solution du système $\begin{cases} 2x + 5y = 1 \\ x - 3y = 6 \end{cases}$ ? Justifier.
\item Sans résoudre, déterminer le nombre de solutions de chaque système :
\[
(a)\ \begin{cases} x + 2y = 5 \\ 2x + 4y = 10 \end{cases} \qquad
(b)\ \begin{cases} 3x - y = 4 \\ 6x - 2y = 7 \end{cases} \qquad
(c)\ \begin{cases} x + y = 4 \\ x - y = 2 \end{cases}
\]
\item Donner une équation cartésienne de la droite passant par les points $A(1\,;\,2)$ et $B(3\,;\,6)$, puis interpréter géométriquement le système formé par cette équation et l'équation $2x - y = 0$.
\end{enumerate}
\end{exercice}

\begin{corrige}[Exercice]
\begin{enumerate}
\item Pour $(3\,;\,-1)$ : $2 \times 3 + 5 \times (-1) = 6 - 5 = 1$ \checkmark\ et $3 - 3 \times (-1) = 3 + 3 = 6$ \checkmark. Le couple vérifie les deux équations : c'est une solution du système.
\item $(a)$ $\dfrac{1}{2} = \dfrac{2}{4} = \dfrac{5}{10}$ : droites confondues, infinité de solutions. $(b)$ $\dfrac{3}{6} = \dfrac{-1}{-2} = \dfrac{1}{2} \neq \dfrac{4}{7}$ : droites strictement parallèles, aucune solution. $(c)$ $\dfrac{1}{1} \neq \dfrac{1}{-1}$ : droites sécantes, unique solution.
\item Le coefficient directeur de $(AB)$ est $\dfrac{6-2}{3-1} = 2$ ; une équation réduite est $y = 2x$, soit $2x - y = 0$. Le système formé par cette équation et $2x - y = 0$ est en fait formé de deux fois la même équation : les droites sont confondues, le système a une infinité de solutions (tous les points de la droite $(AB)$).
\end{enumerate}
\end{corrige}

\begin{aretenir}[L'essentiel de la section]
\begin{itemize}
\item Une équation linéaire à deux inconnues s'écrit $ax + by = c$ avec $(a,b) \neq (0,0)$ ; ses solutions sont les points d'une droite.
\item Un système $2 \times 2$ est l'ensemble de deux équations à satisfaire simultanément ; résoudre, c'est chercher l'intersection de deux droites.
\item Trois cas, et trois seulement : droites sécantes (solution unique), strictement parallèles (aucune solution), confondues (infinité de solutions).
\item Critère algébrique (coefficients non nuls) : comparer $\dfrac{a}{a'}$, $\dfrac{b}{b'}$, $\dfrac{c}{c'}$.
\item Les opérations élémentaires (échange, multiplication par $k \neq 0$, ajout d'un multiple d'une autre équation) transforment un système en système équivalent : c'est l'outil de base de toutes les méthodes de résolution à venir.
\end{itemize}
\end{aretenir}

\coursec{Résolution dans $\mathbb{R}^2$ : substitution et combinaisons linéaires}

Dans la section précédente, nous avons appris à \emph{reconnaître} un système linéaire de deux équations à deux inconnues et à \emph{interpréter géométriquement} ses solutions : chaque équation représente une droite du plan, et résoudre le système revient à chercher les points d'intersection de ces deux droites. Nous savons donc déjà qu'un tel système possède soit une unique solution (droites sécantes), soit aucune (droites strictement parallèles), soit une infinité (droites confondues).

Mais le dessin, aussi parlant soit-il, ne donne qu'une valeur approchée de la solution. Il nous faut maintenant des \textbf{méthodes de calcul} exactes, rapides et sûres. C'est l'objet de cette section : les deux techniques fondamentales que tu devras maîtriser parfaitement au Probatoire sont la \cle{substitution} et les \cle{combinaisons linéaires} (dites aussi méthode d'addition). Bonne nouvelle : ces deux méthodes reposent sur une seule et même idée — \emph{transformer le système en systèmes de plus en plus simples, sans jamais changer son ensemble de solutions}.

\courssub{Le principe général : les systèmes équivalents}

\begin{definition}[Systèmes équivalents]
Deux systèmes d'équations sont dits \cle{équivalents} lorsqu'ils admettent exactement le \emph{même ensemble de solutions}.
\end{definition}

Toute la stratégie de résolution tient en une phrase : on remplace le système de départ par un système équivalent plus simple, puis encore par un autre, jusqu'à obtenir un système « triangulaire » du type $\begin{cases} x = \alpha \\ y = \beta \end{cases}$ qui donne directement la solution.

\begin{propriete}[Opérations élémentaires sur les équations]
On obtient un système équivalent en effectuant l'une des opérations suivantes :
\begin{enumerate}
\item échanger deux équations ;
\item multiplier les deux membres d'une équation par un même réel \emph{non nul} ;
\item ajouter à une équation un multiple d'une autre équation (on dit : remplacer $L_2$ par $L_2 + \lambda L_1$) ;
\item remplacer une inconnue, dans une équation, par son expression tirée d'une autre équation.
\end{enumerate}
\end{propriete}

\begin{experience}[Pourquoi l'addition membre à membre conserve-t-elle les solutions ?]
Considérons le système $(S) : \begin{cases} ax + by = c \quad (L_1) \\ a'x + b'y = c' \quad (L_2) \end{cases}$ et formons le système $(S')$ obtenu en remplaçant $(L_2)$ par $(L_2 + \lambda L_1)$, où $\lambda$ est un réel quelconque :
\[
(S') : \begin{cases} ax + by = c \\ (a' + \lambda a)x + (b' + \lambda b)y = c' + \lambda c \end{cases}
\]
\textbf{Étape 1.} Si $(x\,;y)$ est solution de $(S)$, alors les deux égalités $ax+by=c$ et $a'x+b'y=c'$ sont vraies. En multipliant la première par $\lambda$ et en l'ajoutant membre à membre à la seconde, l'égalité $(a'+\lambda a)x + (b'+\lambda b)y = c' + \lambda c$ est vraie aussi : $(x\,;y)$ est solution de $(S')$.

\textbf{Étape 2.} Réciproquement, si $(x\,;y)$ est solution de $(S')$, on retrouve $(L_2)$ en effectuant $(L_2 + \lambda L_1) - \lambda L_1$ : $(x\,;y)$ est donc solution de $(S)$.

\textbf{Conclusion.} $(S)$ et $(S')$ ont exactement les mêmes solutions : l'opération est réversible, c'est ce qui garantit l'équivalence. C'est ce point précis qui justifie toute la méthode des combinaisons linéaires.
\end{experience}

\courssub{La méthode de substitution}

L'intuition est naturelle : si l'une des équations nous dit que « $y$ vaut telle expression en $x$ », alors partout ailleurs où l'on voit $y$, on peut le \emph{remplacer} par cette expression. On se ramène ainsi à une équation à \emph{une seule} inconnue, que l'on sait résoudre depuis la classe de Troisième.

\begin{methode}[Résolution par substitution]
Pour résoudre $\begin{cases} ax + by = c \\ a'x + b'y = c' \end{cases}$ par substitution :
\begin{enumerate}
\item \textbf{Choisir} une équation et une inconnue à isoler : de préférence une inconnue dont le coefficient vaut $1$ ou $-1$ (cela évite les fractions).
\item \textbf{Exprimer} cette inconnue en fonction de l'autre dans l'équation choisie (par exemple $x = c - by$ si $a = 1$).
\item \textbf{Reporter} (substituer) cette expression dans \emph{l'autre} équation : on obtient une équation du premier degré à une seule inconnue.
\item \textbf{Résoudre} cette équation : on trouve la valeur de l'une des inconnues.
\item \textbf{Remonter} : remplacer cette valeur dans l'expression trouvée à l'étape 2 pour obtenir l'autre inconnue.
\item \textbf{Vérifier} le couple trouvé dans les \emph{deux} équations de départ, puis conclure en écrivant l'ensemble $S$ des solutions.
\end{enumerate}
\end{methode}

\begin{exemplebox}[Substitution, pas à pas]
Résolvons par substitution le système $(S) : \begin{cases} 3x + 2y = 26 \quad (L_1) \\ x + 4y = 28 \quad (L_2) \end{cases}$

\textbf{Étape 1 — Choix stratégique.} Dans $(L_2)$, le coefficient de $x$ vaut $1$ : c'est le candidat idéal, aucune fraction n'apparaîtra à l'étape suivante.

\textbf{Étape 2 — On isole.} De $(L_2)$ on tire : $x = 28 - 4y$.

\textbf{Étape 3 — On reporte dans $(L_1)$} (surtout pas dans $(L_2)$, voir l'attention ci-dessous !) :
\[
3(28 - 4y) + 2y = 26.
\]

\textbf{Étape 4 — On résout.}
\begin{align*}
3(28 - 4y) + 2y &= 26 \\
84 - 12y + 2y &= 26 \\
84 - 10y &= 26 \\
-10y &= 26 - 84 \\
-10y &= -58 \\
y &= \dfrac{-58}{-10} = \dfrac{29}{5}.
\end{align*}

\textbf{Étape 5 — On remonte.} On remplace $y = \dfrac{29}{5}$ dans l'expression de l'étape 2 :
\[
x = 28 - 4 \times \dfrac{29}{5} = \dfrac{140}{5} - \dfrac{116}{5} = \dfrac{24}{5}.
\]

\textbf{Étape 6 — Vérification.} Dans $(L_1)$ : $3\times\dfrac{24}{5} + 2\times\dfrac{29}{5} = \dfrac{72+58}{5} = \dfrac{130}{5} = 26$ \checkmark. Dans $(L_2)$ : $\dfrac{24}{5} + 4\times\dfrac{29}{5} = \dfrac{24+116}{5} = \dfrac{140}{5} = 28$ \checkmark.

\textbf{Conclusion :} $S = \left\lbrace \left(\dfrac{24}{5}\,;\dfrac{29}{5}\right) \right\rbrace$.
\end{exemplebox}

\begin{attention}[Deux pièges classiques de la substitution]
\begin{itemize}
\item \textbf{Reporter dans la mauvaise équation.} Ayant tiré $x = 28 - 4y$ de $(L_2)$, si l'on remplace $x$ dans $(L_2)$ elle-même, on obtient $(28 - 4y) + 4y = 28$, c'est-à-dire $28 = 28$ : une égalité toujours vraie qui ne donne \emph{aucune} information ! On a tourné en rond. \textbf{Règle d'or :} une expression tirée d'une équation doit être reportée dans \emph{l'autre} équation.
\item \textbf{Se fier à une solution « devinée ».} Un couple peut vérifier une équation et pas l'autre : par exemple $(6\,;4)$ vérifie $(L_1)$ car $3\times 6 + 2\times 4 = 26$, mais pas $(L_2)$ car $6 + 4\times 4 = 22 \neq 28$. Seule la vérification dans les \emph{deux} équations fait foi — et il est tout à fait normal qu'une solution soit fractionnaire.
\end{itemize}
\end{attention}

\courssub{La méthode des combinaisons linéaires (ou méthode d'addition)}

L'idée est tout aussi naturelle : si deux quantités égales sont ajoutées à deux quantités égales, les sommes sont égales. En multipliant astucieusement chaque équation par un réel bien choisi, on fait apparaître des coefficients \emph{opposés} devant une même inconnue ; en additionnant membre à membre, cette inconnue \emph{disparaît}.

\begin{methode}[Résolution par combinaisons linéaires]
Pour résoudre $\begin{cases} ax + by = c \quad (L_1)\\ a'x + b'y = c' \quad (L_2) \end{cases}$ :
\begin{enumerate}
\item \textbf{Choisir} l'inconnue à éliminer : celle dont les coefficients ont le plus petit multiple commun (souvent celle qui évite les gros nombres).
\item \textbf{Multiplier} $(L_1)$ et $(L_2)$ par des réels non nuls de sorte que les coefficients de l'inconnue choisie deviennent \emph{opposés} (par exemple $b'$ pour $(L_1)$ et $-b$ pour $(L_2)$).
\item \textbf{Additionner} membre à membre les deux nouvelles équations : l'inconnue choisie s'élimine ; on obtient une équation à une inconnue.
\item \textbf{Résoudre} cette équation.
\item \textbf{Recommencer} en éliminant l'autre inconnue, \emph{ou} remplacer la valeur trouvée dans l'une des équations de départ.
\item \textbf{Vérifier} dans les deux équations initiales, puis écrire $S$.
\end{enumerate}
\end{methode}

\begin{exemplebox}[Combinaisons linéaires sur le même système]
Reprenons $(S) : \begin{cases} 3x + 2y = 26 \quad (L_1) \\ x + 4y = 28 \quad (L_2) \end{cases}$

\textbf{Élimination de $x$.} Les coefficients de $x$ sont $3$ et $1$ : il suffit de multiplier $(L_2)$ par $-3$, puis d'additionner membre à membre avec $(L_1)$.
\[
\begin{array}{rcl}
(L_1) & : & 3x + 2y = 26 \\
-3(L_2) & : & -3x - 12y = -84 \\ \hline
(L_1) - 3(L_2) & : & -10y = -58
\end{array}
\]
d'où $y = \dfrac{-58}{-10} = \dfrac{29}{5}$.

\textbf{Élimination de $y$.} Multiplions $(L_1)$ par $-2$ et additionnons avec $(L_2)$ : les coefficients de $y$ deviennent $-4$ et $+4$.
\[
-2(L_1) + (L_2) : \quad (-6x + x) + (-4y + 4y) = -52 + 28 \iff -5x = -24 \iff x = \dfrac{24}{5}.
\]

\textbf{Vérification.} $3\times\dfrac{24}{5} + 2\times\dfrac{29}{5} = \dfrac{72+58}{5} = \dfrac{130}{5} = 26$ \checkmark{} ; $\dfrac{24}{5} + 4\times\dfrac{29}{5} = \dfrac{24+116}{5} = \dfrac{140}{5} = 28$ \checkmark.

\textbf{Conclusion :} $S = \left\lbrace \left(\dfrac{24}{5}\,;\dfrac{29}{5}\right) \right\rbrace$. Les deux méthodes donnent (heureusement !) le même résultat.
\end{exemplebox}

\begin{exoresolu}[Un problème de la vie courante]
\textbf{Énoncé.} Au marché Mokolo, Maman Béa achète $3$~kg de tomates et $2$~kg d'oignons pour \SI{2600}{FCFA}. Le lendemain, aux mêmes prix, elle achète $1$~kg de tomates et $4$~kg d'oignons pour \SI{2800}{FCFA}. Quel est le prix du kilogramme de chaque produit ?

\tcblower
\textbf{Solution détaillée.} Soit $x$ le prix d'un kg de tomates et $y$ celui d'un kg d'oignons (en FCFA). Le problème se traduit par :
\[
(S) : \begin{cases} 3x + 2y = 2600 \quad (L_1) \\ x + 4y = 2800 \quad (L_2) \end{cases}
\]
C'est exactement le système étudié ci-dessus, dont les seconds membres sont multipliés par $100$ ! Résolvons par combinaison : $(L_1) - 3(L_2)$ donne
\[
(3x + 2y) - 3(x + 4y) = 2600 - 8400 \iff -10y = -5800 \iff y = 580.
\]
En reportant dans $(L_2)$ : $x = 2800 - 4 \times 580 = 2800 - 2320 = 480$.

\textbf{Vérification :} $3 \times 480 + 2 \times 580 = 1440 + 1160 = 2600$ \checkmark{} ; $480 + 4 \times 580 = 480 + 2320 = 2800$ \checkmark.

\textbf{Conclusion :} le kilogramme de tomates coûte \SI{480}{FCFA} et celui d'oignons \SI{580}{FCFA}. Remarque au passage : les prix « ronds » du problème concret confirment la cohérence de nos calculs — les deux situations ne diffèrent que par une échelle de $100$ sur les seconds membres, donc les solutions sont aussi multipliées par $100$ : $\dfrac{24}{5}\times 100 = 480$ et $\dfrac{29}{5}\times 100 = 580$.
\end{exoresolu}

\courssub{Comment choisir sa stratégie ?}

Face à un système, quelques secondes d'observation permettent de gagner beaucoup de temps. L'organigramme suivant résume la démarche du bon élève.

\begin{popfigure}
\begin{center}
\begin{tikzpicture}[
  node distance=0.8cm,
  decision/.style={diamond, draw=poporange, fill=poporangeL, text width=3.8cm, align=center, inner sep=1pt, font=\small, aspect=2.4},
  block/.style={rectangle, rounded corners, draw=popblue, fill=popblueL, text width=4.4cm, align=center, font=\small, minimum height=0.9cm},
  start/.style={rectangle, rounded corners, draw=popgreen, fill=popgreenL, text width=5.2cm, align=center, font=\small\bfseries},
  fleche/.style={-{Stealth[length=2.5mm]}, thick, popdark}
]
\node[start, align=center] (depart) {Observer le système $\begin{cases} ax+by=c \\ a'x+b'y=c' \end{cases}$};
\node[decision, below=of depart] (coef1) {Un coefficient vaut-il $1$ ou $-1$ ?};
\node[block, below left=1.2cm and -1.2cm of coef1, align=center] (subst) {\textbf{Substitution}\\ Isoler cette inconnue, reporter dans l'autre équation};
\node[decision, below right=1.2cm and -1.2cm of coef1] (mult) {Un coefficient est-il multiple de l'autre ?};
\node[block, below left=1.1cm and -1.0cm of mult, align=center] (comb1) {\textbf{Combinaison simple}\\ Multiplier une seule équation, puis additionner};
\node[block, below right=1.1cm and -1.0cm of mult, align=center] (comb2) {\textbf{Combinaison générale}\\ Multiplier les deux équations, puis additionner};
\node[block, below=5.6cm of coef1] (verif) {\textbf{Toujours vérifier} le couple trouvé dans les deux équations, puis écrire $S$};
\draw[fleche] (depart) -- (coef1);
\draw[fleche] (coef1) -| node[pos=0.25, above, font=\small]{oui} (subst);
\draw[fleche] (coef1) -| node[pos=0.25, above, font=\small]{non} (mult);
\draw[fleche] (mult) -| node[pos=0.25, above, font=\small]{oui} (comb1);
\draw[fleche] (mult) -| node[pos=0.25, above, font=\small]{non} (comb2);
\draw[fleche] (subst.south) |- (verif.west);
\draw[fleche] (comb1.south) |- (verif.west);
\draw[fleche] (comb2.south) |- (verif.east);
\end{tikzpicture}
\end{center}
\legende{Arbre de décision : choisir entre substitution et combinaisons linéaires selon la forme des coefficients.}
\end{popfigure}

\begin{aretenir}[Quelle méthode privilégier ?]
\begin{itemize}
\item \textbf{Substitution} : idéale lorsqu'un coefficient vaut $1$ ou $-1$ (on isole sans fraction). Rapide, mais devient pénible si elle fait apparaître des fractions compliquées.
\item \textbf{Combinaisons linéaires} : idéale lorsque les coefficients sont « ronds » ou multiples l'un de l'autre, et \emph{indispensable} dès qu'aucun coefficient ne vaut $\pm 1$. C'est aussi elle qui se généralise le mieux aux systèmes de trois équations (pivot de Gauss, section 5).
\item Dans les deux cas, la \textbf{vérification finale} est systématique : elle coûte trente secondes et sauve des points au Probatoire.
\end{itemize}
\end{aretenir}

\courssub{Les cas dégénérés : quand les inconnues « disparaissent »}

Il arrive qu'en éliminant une inconnue, \emph{l'autre disparaisse aussi}. Ce n'est pas un échec : c'est un message. Deux situations peuvent se produire, correspondant exactement aux deux cas géométriques vus dans la section précédente (droites confondues ou strictement parallèles).

\begin{exemplebox}[Infinité de solutions : aboutir à $0 = 0$]
Résolvons $(S) : \begin{cases} 2x - y = 3 \quad (L_1) \\ 4x - 2y = 6 \quad (L_2) \end{cases}$

Effectuons $(L_2) - 2(L_1)$ :
\[
(4x - 2y) - 2(2x - y) = 6 - 2 \times 3 \iff 0 = 0.
\]
L'égalité $0 = 0$ est \emph{toujours vraie} : cela signifie que $(L_2)$ n'apporte aucune information nouvelle — c'est simplement le double de $(L_1)$. Les deux équations représentent la \emph{même droite}. Tout couple vérifiant la seule équation $2x - y = 3$ est solution. En exprimant $y$ en fonction de $x$ : $y = 2x - 3$. On écrit alors l'ensemble des solutions à l'aide d'un paramètre réel :
\[
S = \left\lbrace (x\,;\, 2x - 3),\; x \in \mathbb{R} \right\rbrace.
\]
Ainsi $(0\,;-3)$, $(1\,;-1)$, $(2\,;1)$, $(3\,;3)$... sont toutes des solutions : il y en a une infinité.
\end{exemplebox}

\begin{exemplebox}[Aucune solution : aboutir à $0 = k$ avec $k \neq 0$]
Résolvons $(S') : \begin{cases} 2x - y = 3 \quad (L_1) \\ 4x - 2y = 10 \quad (L_2) \end{cases}$

Effectuons $(L_2) - 2(L_1)$ :
\[
(4x - 2y) - 2(2x - y) = 10 - 6 \iff 0 = 4.
\]
L'égalité $0 = 4$ est \emph{fausse}, quelles que soient les valeurs de $x$ et $y$ : le système n'admet \textbf{aucune solution}. Géométriquement, les deux droites ont le même coefficient directeur ($2$) mais des ordonnées à l'origine différentes ($-3$ et $-5$) : elles sont strictement parallèles. On conclut :
\[
S = \emptyset.
\]
\end{exemplebox}

\begin{attention}[Bien rédiger la conclusion dans les cas dégénérés]
\begin{itemize}
\item Obtenir $0 = 0$ ne signifie \emph{pas} « tout nombre est solution » : seuls les couples $(x\,;y)$ vérifiant l'équation restante conviennent. Il faut écrire $S$ \textbf{à l'aide d'un paramètre}, par exemple $S = \{(x\,;2x-3),\, x \in \mathbb{R}\}$, et non « $S = \mathbb{R}$ » (faux : le couple $(0\,;0)$ n'est pas solution !).
\item Obtenir $0 = k$ avec $k \neq 0$ impose la conclusion $S = \emptyset$, rédigée explicitement.
\item Ne jamais conclure « impossible » sans avoir écrit l'égalité absurde obtenue : c'est elle qui justifie la réponse.
\end{itemize}
\end{attention}

\begin{savaistu}[Un peu d'histoire]
La méthode d'élimination par combinaisons était déjà décrite dans le \emph{Jiuzhang Suanshu} (« Les Neuf Chapitres sur l'art du calcul »), traité chinois datant d'environ 200 avant J.-C., où l'on résolvait des systèmes liés au partage des récoltes de grains. En Europe, c'est Carl Friedrich Gauss (1777--1855) qui systématisa la méthode d'élimination qui porte aujourd'hui son nom — et que tu rencontreras dans la section 5 pour les systèmes à trois inconnues. Au Cameroun, les mêmes calculs servent chaque jour : répartition des coûts dans les coopératives agricoles, tarification des forfaits téléphoniques, mélanges en agroalimentaire...
\end{savaistu}

\begin{exercice}[À toi de jouer]
\begin{enumerate}
\item Résoudre par substitution : $\begin{cases} x - 3y = 1 \\ 2x + 5y = 35 \end{cases}$
\item Résoudre par combinaisons linéaires : $\begin{cases} 5x + 3y = 44 \\ 7x - 2y = 12 \end{cases}$
\item Résoudre et interpréter géométriquement : $\begin{cases} 3x - 6y = 9 \\ x - 2y = 3 \end{cases}$ puis $\begin{cases} 3x - 6y = 9 \\ x - 2y = 5 \end{cases}$
\end{enumerate}
\end{exercice}

\begin{corrige}[Exercice 1]
\textbf{1.} De la première équation : $x = 1 + 3y$. On reporte dans la seconde : $2(1+3y) + 5y = 35 \iff 2 + 6y + 5y = 35 \iff 11y = 33 \iff y = 3$. Alors $x = 1 + 3\times 3 = 10$. Vérification : $10 - 9 = 1$ \checkmark{} ; $20 + 15 = 35$ \checkmark. $S = \{(10\,;3)\}$.

\textbf{2.} Éliminons $y$ : $2(L_1) + 3(L_2)$ donne $(10x + 6y) + (21x - 6y) = 88 + 36$, soit $31x = 124$, d'où $x = 4$. En reportant dans $(L_1)$ : $5 \times 4 + 3y = 44 \iff 3y = 24 \iff y = 8$. Vérification dans $(L_2)$ : $7\times 4 - 2\times 8 = 28 - 16 = 12$ \checkmark. $S = \{(4\,;8)\}$.

\textbf{3.} Premier système : $(L_1) - 3(L_2)$ donne $0 = 0$ : les deux équations représentent la même droite (droites confondues). De $x - 2y = 3$ on tire $y = \dfrac{x-3}{2}$, donc $S = \left\lbrace \left(x\,;\dfrac{x-3}{2}\right),\, x \in \mathbb{R} \right\rbrace$. Second système : $(L_1) - 3(L_2)$ donne $0 = 9 - 15$, c'est-à-dire $0 = -6$, égalité absurde : $S = \emptyset$ (droites strictement parallèles, de même coefficient directeur $\tfrac{1}{2}$ mais d'ordonnées à l'origine différentes).
\end{corrige}

\coursec{Déterminant d'un système $2\times 2$ et formules de Cramer}

Dans la section précédente, tu as appris à résoudre un système de deux équations à deux inconnues par \cle{substitution} et par \cle{combinaisons linéaires}. Ces deux méthodes sont efficaces, mais elles demandent de la vigilance : fractions, signes, substitutions maladroites... Une question naturelle se pose alors : n'existe-t-il pas un \emph{critère rapide} permettant de savoir, avant même de résoudre, si un système possède une unique solution ? Et mieux encore : n'existe-t-il pas des \emph{formules toutes faites} donnant directement $x$ et $y$ ?

La réponse est oui, et elle tient en un seul nombre, calculé à partir des quatre coefficients du système : le \cle{déterminant}. C'est l'un des outils les plus puissants et les plus élégants de l'algèbre linéaire, et il est exigible au Probatoire.

\courssub{Le déterminant d'un système $2\times 2$}

\textbf{Intuition.} Considère le système général
\[
(S) : \begin{cases} ax + by = c \\ a'x + b'y = c' \end{cases}
\]
où $a, b, a', b'$ sont des réels donnés (avec $(a,b)\neq(0,0)$ et $(a',b')\neq(0,0)$, pour que chaque ligne soit bien une équation de droite). Géométriquement, chaque équation représente une droite du plan. Le système a une unique solution exactement quand les deux droites sont \emph{sécantes}, c'est-à-dire quand elles n'ont pas la même direction. Or deux droites ont la même direction lorsque leurs coefficients sont proportionnels : $\dfrac{a}{a'}=\dfrac{b}{b'}$, ce qui s'écrit encore, en chassant les dénominateurs, $ab'-a'b=0$. Le nombre $ab'-a'b$ mesure donc, en quelque sorte, le \og défaut de parallélisme \fg{} des deux droites. C'est lui que l'on appelle déterminant.

\begin{definition}[Déterminant d'un système $2\times 2$]
Le \cle{déterminant} du système $(S) : \begin{cases} ax + by = c \\ a'x + b'y = c' \end{cases}$ est le nombre réel, noté $D$, défini par
\[
D = ab' - a'b.
\]
On l'écrit souvent sous forme de tableau encadré de barres verticales :
\[
D = \begin{vmatrix} a & b \\ a' & b' \end{vmatrix} = ab' - a'b.
\]
\end{definition}

\begin{attention}[Le déterminant ne dépend que des coefficients des inconnues]
Les nombres $c$ et $c'$ (les \emph{seconds membres}) n'interviennent \textbf{pas} dans le calcul de $D$. Le déterminant se calcule uniquement avec les coefficients de $x$ et de $y$, rangés dans le même ordre que dans le système. Veille aussi à écrire le système sous la \emph{forme standard} $ax+by=c$ avant d'identifier $a, b, a', b'$ : si le système est donné sous la forme $\begin{cases} 3y - 2x = 5 \\ x + 4y = 1 \end{cases}$, il faut d'abord réécrire la première équation $-2x+3y=5$, sinon tu permutes les coefficients et le résultat sera faux.
\end{attention}

\textbf{Moyen mnémotechnique : le calcul en croix.} Pour retenir la formule $D=ab'-a'b$, visualise le tableau des coefficients et effectue les produits en diagonale : la diagonale \emph{descendante} compte positivement, la diagonale \emph{montante} compte négativement.

\begin{popfigure}
\begin{center}
\begin{tikzpicture}[scale=1.4]
\node (a) at (0,1) {$a$};
\node (b) at (1.5,1) {$b$};
\node (a2) at (0,0) {$a'$};
\node (b2) at (1.5,0) {$b'$};
\draw[->,very thick,popgreen] (a.south east) -- (b2.north west) node[midway,above right=1pt,popgreen] {$+\,a\times b'$};
\draw[->,very thick,poporange] (a2.north east) -- (b.south west) node[midway,below right=1pt,poporange] {$-\,a'\times b$};
\node at (0.75,-1) {$D \;=\; \textcolor{popgreen}{a\,b'} \;-\; \textcolor{poporange}{a'\,b}$};
\end{tikzpicture}
\end{center}
\legende{Calcul en croix du déterminant : produit de la diagonale descendante (positif) moins produit de la diagonale montante (négatif).}
\end{popfigure}

\begin{exemplebox}[Entraînement au calcul rapide de déterminants]
Calcule mentalement, puis vérifie :
\begin{enumerate}
\item $\begin{vmatrix} 3 & 2 \\ 1 & 5 \end{vmatrix} = 3\times 5 - 1\times 2 = 15-2 = 13$ ;
\item $\begin{vmatrix} 4 & -2 \\ 3 & 7 \end{vmatrix} = 4\times 7 - 3\times(-2) = 28+6 = 34$ \quad (attention au double signe !) ;
\item $\begin{vmatrix} -5 & 6 \\ -2 & 3 \end{vmatrix} = (-5)\times 3 - (-2)\times 6 = -15+12 = -3$ ;
\item $\begin{vmatrix} \frac{1}{2} & 4 \\ 3 & \frac{2}{3} \end{vmatrix} = \frac{1}{2}\times\frac{2}{3} - 3\times 4 = \frac{1}{3}-12 = -\frac{35}{3}$ \quad (les fractions ne changent rien à la règle) ;
\item $\begin{vmatrix} 2 & 6 \\ 1 & 3 \end{vmatrix} = 2\times 3 - 1\times 6 = 0$ \quad (coefficients proportionnels : droites parallèles ou confondues).
\end{enumerate}
\end{exemplebox}

\courssub{Théorème fondamental : condition d'unicité de la solution}

Voici le résultat central de cette section. Sa démonstration est \textbf{exigible} : elle repose sur une simple combinaison linéaire, que tu maîtrises déjà depuis la section précédente.

\begin{propriete}[Théorème — existence et unicité (démonstration exigible)]
Le système $(S) : \begin{cases} ax + by = c \\ a'x + b'y = c' \end{cases}$ admet une \textbf{unique solution} si et seulement si son déterminant est non nul :
\[
D = ab'-a'b \neq 0 \iff (S) \text{ admet un unique couple solution } (x\,;\,y).
\]
\end{propriete}

\begin{experience}[Démonstration guidée du théorème]
On part du système $(S) : \begin{cases} ax + by = c \quad (L_1)\\ a'x + b'y = c' \quad (L_2) \end{cases}$.

\textbf{Étape 1 : éliminer $y$ par combinaison linéaire.} Multiplions $(L_1)$ par $b'$ et $(L_2)$ par $b$, puis soustrayons membre à membre : on effectue la combinaison $b'(L_1)-b(L_2)$. Les termes en $y$ disparaissent :
\[
b'(ax+by) - b(a'x+b'y) = b'c - bc' \iff (ab'-a'b)\,x = cb'-c'b.
\]
En notant $D=ab'-a'b$, on obtient l'équation \og réduite \fg{} :
\[
D\,x = cb'-c'b. \qquad (\star)
\]

\textbf{Étape 2 : discussion selon la valeur de $D$.}

\emph{Cas 1 : $D\neq 0$.} L'équation $(\star)$ donne immédiatement $x=\dfrac{cb'-c'b}{D}$, valeur unique. En éliminant $x$ de la même façon (combinaison $a(L_2)-a'(L_1)$), on obtient $D\,y = ac'-a'c$, d'où une valeur unique de $y$. Le système admet donc \textbf{une unique solution}.

\emph{Cas 2 : $D=0$.} L'équation $(\star)$ devient $0\times x = cb'-c'b$.
\begin{itemize}
\item Si $cb'-c'b\neq 0$ : l'égalité $0 = cb'-c'b$ est \emph{impossible} : le système n'a \textbf{aucune solution}.
\item Si $cb'-c'b=0$ : l'égalité $0=0$ est toujours vraie. Comme $D=0$, les coefficients $(a,b)$ et $(a',b')$ sont proportionnels ; la condition $cb'-c'b=0$ signifie que les seconds membres respectent la même proportionnalité : les deux équations sont alors \emph{équivalentes} et représentent la \emph{même droite}. Le système a une \textbf{infinité de solutions}.
\end{itemize}

\textbf{Conclusion.} Le système a une unique solution si et seulement si $D\neq 0$. $\blacksquare$
\end{experience}

\courssub{Les formules de Cramer}

La démonstration précédente fait mieux que prouver l'existence d'une solution : elle la \emph{calcule} explicitement. Introduisons deux nouveaux déterminants, obtenus en remplaçant dans $D$ une colonne par la colonne des seconds membres.

\begin{definition}[Déterminants $D_x$ et $D_y$]
Pour le système $(S)$, on pose :
\[
D_x = \begin{vmatrix} c & b \\ c' & b' \end{vmatrix} = cb'-c'b
\qquad\text{et}\qquad
D_y = \begin{vmatrix} a & c \\ a' & c' \end{vmatrix} = ac'-a'c.
\]
$D_x$ s'obtient en remplaçant la colonne des coefficients de $x$ par les seconds membres ; $D_y$ en remplaçant la colonne des coefficients de $y$.
\end{definition}

\begin{propriete}[Formules de Cramer]
Si $D\neq 0$, l'unique solution du système $(S)$ est donnée par les \cle{formules de Cramer} :
\[
x = \frac{D_x}{D} = \frac{cb'-c'b}{ab'-a'b}
\qquad\text{et}\qquad
y = \frac{D_y}{D} = \frac{ac'-a'c}{ab'-a'b}.
\]
\end{propriete}

\begin{aretenir}[Ce qu'il faut retenir]
\begin{itemize}
\item $D\neq 0$ : unique solution, donnée par $x=\dfrac{D_x}{D}$ et $y=\dfrac{D_y}{D}$.
\item $D=0$ : les formules de Cramer sont \emph{inutilisables} ; il faut discuter (aucune solution ou infinité).
\item $D_x$ : on remplace la colonne de $x$ par les seconds membres. $D_y$ : on remplace la colonne de $y$.
\end{itemize}
\end{aretenir}

\begin{methode}[Résoudre un système par le déterminant en 4 étapes]
\begin{enumerate}
\item \textbf{Écrire} le système sous la forme standard $\begin{cases} ax+by=c \\ a'x+b'y=c' \end{cases}$ et identifier $a,b,c,a',b',c'$.
\item \textbf{Calculer} $D=ab'-a'b$. Si $D=0$, s'arrêter et discuter séparément (voir plus bas).
\item \textbf{Calculer} $D_x=cb'-c'b$ et $D_y=ac'-a'c$.
\item \textbf{Conclure} : $x=\dfrac{D_x}{D}$, $y=\dfrac{D_y}{D}$, puis \emph{vérifier} en substituant dans les deux équations.
\end{enumerate}
\end{methode}

\begin{exoresolu}[Résolution par les formules de Cramer]
Résoudre dans $\mathbb{R}^2$ le système $\begin{cases} 5x-3y=1 \\ 2x+y=7 \end{cases}$.
\tcblower
\textbf{Étape 1.} Le système est sous forme standard : $a=5$, $b=-3$, $c=1$, $a'=2$, $b'=1$, $c'=7$.

\textbf{Étape 2.} Déterminant principal :
\[
D=\begin{vmatrix} 5 & -3 \\ 2 & 1 \end{vmatrix}=5\times 1-2\times(-3)=5+6=11.
\]
$D=11\neq 0$ : le système admet une unique solution.

\textbf{Étape 3.} Déterminants auxiliaires :
\[
D_x=\begin{vmatrix} 1 & -3 \\ 7 & 1 \end{vmatrix}=1\times 1-7\times(-3)=1+21=22,
\qquad
D_y=\begin{vmatrix} 5 & 1 \\ 2 & 7 \end{vmatrix}=5\times 7-2\times 1=35-2=33.
\]

\textbf{Étape 4.} Formules de Cramer :
\[
x=\frac{D_x}{D}=\frac{22}{11}=2,
\qquad
y=\frac{D_y}{D}=\frac{33}{11}=3.
\]
\emph{Vérification} : $5\times 2-3\times 3=10-9=1$ \checkmark{} et $2\times 2+3=7$ \checkmark.

\textbf{Conclusion} : $\mathcal{S}=\{(2\,;\,3)\}$.
\end{exoresolu}

\courssub{Le cas $D=0$ : discussion et interprétation géométrique}

\begin{attention}[Piège classique au Probatoire]
Les formules de Cramer ne s'appliquent \textbf{que si $D\neq 0$}. Diviser par $D$ quand $D=0$ est une erreur grave (division par zéro !). Lorsque $D=0$, il faut revenir aux équations et comparer : les deux équations sont-elles \emph{incompatibles} (droites strictement parallèles, aucune solution) ou \emph{équivalentes} (droites confondues, infinité de solutions) ? Le critère pratique : si $D=0$, les coefficients sont proportionnels ; il suffit alors de vérifier si les seconds membres respectent la même proportionnalité, ce qui revient à tester le signe de $D_x$ (ou de $D_y$).
\end{attention}

\begin{exemplebox}[Discussion quand $D=0$]
\textbf{(a)} $\begin{cases} 2x+4y=6 \\ x+2y=3 \end{cases}$ : $D=2\times 2-1\times 4=0$. La deuxième équation multipliée par $2$ redonne exactement la première : les équations sont équivalentes. Il y a une \textbf{infinité de solutions} : tous les couples $(x\,;\,y)$ avec $y=\dfrac{3-x}{2}$, par exemple $(1\,;\,1)$, $(3\,;\,0)$, $(5\,;\,-1)$\dots

\textbf{(b)} $\begin{cases} 2x+4y=6 \\ x+2y=5 \end{cases}$ : $D=0$ toujours, mais cette fois $2\times 5=10\neq 6$ : la proportionnalité n'est pas respectée par les seconds membres. Soustraire deux fois la deuxième équation de la première donnerait $0=-4$, ce qui est impossible : \textbf{aucune solution}. Les droites sont strictement parallèles.
\end{exemplebox}

Le tableau suivant résume toute la situation, avec l'interprétation géométrique associée (chaque équation du système représente une droite du plan).

\begin{center}
\begin{tabularx}{\linewidth}{|c|X|X|c|}
\hline
\rowcolor{popblueL} \textbf{Signe de $D$} & \textbf{Position des droites} & \textbf{Ensemble des solutions} & \textbf{Illustration} \\
\hline
$D\neq 0$ & Droites \textbf{sécantes} en un point & Unique solution $(x\,;\,y)$, donnée par Cramer &
\begin{tikzpicture}[scale=0.32,baseline]
\draw[popmuted] (-2,0)--(2,0); \draw[popmuted] (0,-1.6)--(0,1.6);
\draw[thick,popblue] (-1.6,-1.3)--(1.6,1.5);
\draw[thick,poporange] (-1.6,1.2)--(1.6,-1.1);
\fill[popink] (0.12,0.26) circle (0.13);
\end{tikzpicture} \\
\hline
$D=0$ et $D_x\neq 0$ & Droites \textbf{strictement parallèles} & Aucune solution : $\mathcal{S}=\emptyset$ &
\begin{tikzpicture}[scale=0.32,baseline]
\draw[popmuted] (-2,0)--(2,0); \draw[popmuted] (0,-1.6)--(0,1.6);
\draw[thick,popblue] (-1.7,-0.6)--(1.7,0.9);
\draw[thick,poporange] (-1.7,-1.5)--(1.7,0);
\end{tikzpicture} \\
\hline
$D=0$ et $D_x=0$ & Droites \textbf{confondues} & Infinité de solutions (toute la droite) &
\begin{tikzpicture}[scale=0.32,baseline]
\draw[popmuted] (-2,0)--(2,0); \draw[popmuted] (0,-1.6)--(0,1.6);
\draw[thick,popblue] (-1.7,-0.8)--(1.7,1.1);
\draw[thick,poporange,dashed] (-1.6,-0.72)--(1.8,1.18);
\end{tikzpicture} \\
\hline
\end{tabularx}
\end{center}

\begin{attention}[Rédaction : ne pas confondre les cas]
Quand $D=0$, ne te contente pas d'écrire \og pas de solution unique \fg. Le correcteur attend une conclusion précise : soit $\mathcal{S}=\emptyset$, soit une infinité de solutions que l'on décrit (par exemple $\mathcal{S}=\left\{\left(x\,;\,\tfrac{3-x}{2}\right),\ x\in\mathbb{R}\right\}$). Et n'oublie jamais la vérification finale quand tu appliques Cramer : elle prend dix secondes et élimine la plupart des erreurs de signe.
\end{attention}

\begin{savaistu}[Gabriel Cramer et ses formules]
Gabriel Cramer (1704--1752) était un mathématicien suisse, professeur à Genève. En 1750, dans son \emph{Introduction à l'analyse des lignes courbes algébriques}, il publia les formules qui portent aujourd'hui son nom, valables non seulement pour deux équations, mais pour des systèmes de $n$ équations à $n$ inconnues. Ironie de l'histoire : l'Écossais Colin Maclaurin connaissait ces formules avant lui ! Aujourd'hui, les déterminants sont partout : en informatique graphique pour calculer des aires et tester des positions, en économie pour résoudre des modèles de marché, et dans les algorithmes de traitement du signal utilisés par les réseaux de téléphonie mobile. Le petit calcul en croix que tu viens d'apprendre est la porte d'entrée de l'algèbre linéaire, l'un des outils mathématiques les plus utilisés au monde.
\end{savaistu}

\begin{exercice}[À toi de jouer]
\begin{enumerate}
\item Résoudre par les formules de Cramer : $\begin{cases} 3x+2y=8 \\ x-4y=-2 \end{cases}$.
\item Résoudre et discuter : $\begin{cases} 6x-3y=9 \\ 2x-y=3 \end{cases}$ puis $\begin{cases} 6x-3y=9 \\ 2x-y=5 \end{cases}$.
\item Déterminer le réel $m$ pour que le système $\begin{cases} mx+2y=4 \\ 3x+y=5 \end{cases}$ admette une unique solution.
\end{enumerate}
\end{exercice}

\begin{corrige}[Exercice]
\textbf{1.} $D=3\times(-4)-1\times 2=-12-2=-14\neq 0$. $D_x=8\times(-4)-(-2)\times 2=-32+4=-28$, $D_y=3\times(-2)-1\times 8=-6-8=-14$. Donc $x=\dfrac{-28}{-14}=2$ et $y=\dfrac{-14}{-14}=1$. Vérification : $3\times 2+2\times 1=8$ \checkmark{} et $2-4\times 1=-2$ \checkmark. $\mathcal{S}=\{(2\,;\,1)\}$.

\textbf{2.} Dans les deux cas, $D=6\times(-1)-2\times(-3)=-6+6=0$. Premier système : la deuxième équation multipliée par $3$ redonne la première ($6x-3y=9$) : droites confondues, infinité de solutions $\mathcal{S}=\{(x\,;\,2x-3),\ x\in\mathbb{R}\}$. Deuxième système : $3\times 5=15\neq 9$ : la proportionnalité n'est pas respectée par les seconds membres, les droites sont strictement parallèles, $\mathcal{S}=\emptyset$.

\textbf{3.} $D=m\times 1-3\times 2=m-6$. Le système admet une unique solution si et seulement si $D\neq 0$, c'est-à-dire $m\neq 6$.
\end{corrige}

Tu disposes maintenant de trois méthodes pour résoudre un système dans $\mathbb{R}^2$ : substitution, combinaisons linéaires et déterminant. Mais que se passe-t-il lorsque les seconds membres contiennent une \emph{lettre} inconnue, un paramètre, et que l'on demande de discuter selon ses valeurs ? C'est l'objet de la section suivante.

\coursec{Systèmes dans $\mathbb{R}^2$ à second membre dépendant d'un paramètre}

Dans la section précédente, tu as appris à résoudre un système de deux équations à deux inconnues grâce au déterminant et aux formules de Cramer. Mais dans la vie réelle — et au Probatoire ! — les données ne sont pas toujours figées. Imagine un commerçant du marché Mokolo qui sait que la recette de la journée dépend du cours des produits : l'équation devient $2x + y = m$, où $m$ est un \emph{paramètre} qui peut varier. Résoudre le système, c'est alors donner la solution \emph{en fonction de $m$}, puis discuter ce qui se passe pour chaque valeur de $m$. C'est exactement l'objet de cette section.

\courssub{Qu'est-ce qu'un système à paramètre ?}

\begin{definition}[Système linéaire à paramètre dans $\mathbb{R}^2$]
Un \cle{système linéaire à paramètre} est un système de la forme
\[
\begin{cases} ax + by = c \\ a'x + b'y = c' \end{cases}
\]
dans lequel certains coefficients (le plus souvent $c$ et $c'$, parfois aussi $a$, $b$, $a'$, $b'$) dépendent d'un nombre réel $m$ appelé \cle{paramètre}. Résoudre le système, c'est déterminer, \textbf{pour chaque valeur de $m$}, l'ensemble $S$ des couples $(x\,;\,y)$ solutions.
\end{definition}

L'idée fondamentale est la suivante : le paramètre $m$ n'est \emph{pas} une inconnue. C'est un nombre fixé, mais dont on ne connaît pas la valeur. On le manipule donc exactement comme un nombre ordinaire ($3$, $-5$, $\sqrt{2}$...) dans tous les calculs : on peut l'ajouter, le soustraire, le multiplier. La seule prudence à avoir concerne la \textbf{division} : on ne peut diviser par une expression contenant $m$ que si l'on est certain qu'elle n'est pas nulle — sinon, on traite le cas problématique à part.

\begin{aretenir}[Principe directeur]
On résout le système en traitant $m$ comme une constante. On obtient en général une solution $(x\,;\,y)$ exprimée en fonction de $m$. Puis on \cle{discute} : on examine les valeurs particulières de $m$ pour lesquelles le raisonnement précédent ne s'applique pas (division par zéro, déterminant nul).
\end{aretenir}

\courssub{Cas où le déterminant ne dépend pas de $m$}

C'est la situation la plus simple, et c'est celle visée par le programme lorsqu'il parle de systèmes « dont le second membre dépend d'un paramètre ». Le déterminant
\[
D = \begin{vmatrix} a & b \\ a' & b' \end{vmatrix} = ab' - a'b
\]
ne porte que sur les coefficients des inconnues : il ne contient donc \textbf{pas} $m$. Deux cas se présentent alors.

\begin{propriete}[Système à second membre paramétré]
Soit le système $\begin{cases} ax + by = c(m) \\ a'x + b'y = c'(m) \end{cases}$ dont le déterminant $D = ab' - a'b$ est indépendant de $m$.
\begin{itemize}
\item Si $D \neq 0$ : le système admet, \textbf{pour toute valeur de $m$}, une solution unique donnée par les formules de Cramer, qui s'expriment en fonction de $m$ :
\[ x = \frac{c(m)\,b' - b\,c'(m)}{D} \qquad \text{et} \qquad y = \frac{a\,c'(m) - c(m)\,a'}{D}. \]
\item Si $D = 0$ : les deux équations sont proportionnelles du côté des inconnues ; selon la valeur de $m$, le système admet une infinité de solutions (équations compatibles) ou aucune solution (équations incompatibles). Ce cas se discute valeur par valeur.
\end{itemize}
\end{propriete}

\begin{methode}[Résoudre un système à paramètre]
\begin{enumerate}
\item \textbf{Identifier} le paramètre et vérifier s'il apparaît dans le déterminant $D$ ou seulement dans les seconds membres.
\item \textbf{Calculer $D$}. Si $D \neq 0$, résoudre par substitution, combinaison linéaire ou Cramer en traitant $m$ comme un nombre fixe ; exprimer $x$ et $y$ en fonction de $m$.
\item \textbf{Discuter} les valeurs de $m$ qui annulent $D$ (ou qui annulent un coefficient par lequel on a divisé) : reprendre le système pour chacune de ces valeurs et conclure (infinité de solutions ou ensemble vide).
\item \textbf{Conclure} en résumant clairement l'ensemble $S$ pour chaque cas.
\end{enumerate}
\end{methode}

\begin{exemplebox}[Résolution et problème inverse]
Résolvons le système $(S_m) : \begin{cases} 2x + y = m \\ x - y = 3 \end{cases}$ puis déterminons $m$ pour que la solution vérifie $x = 2y$.

\textbf{Étape 1 : le déterminant.} $D = \begin{vmatrix} 2 & 1 \\ 1 & -1 \end{vmatrix} = 2\times(-1) - 1\times 1 = -3 \neq 0$. Le système admet donc une solution unique pour tout réel $m$.

\textbf{Étape 2 : résolution par combinaison.} Additionnons les deux équations : $3x = m + 3$, d'où $x = \dfrac{m+3}{3}$. La deuxième équation donne $y = x - 3 = \dfrac{m+3}{3} - 3 = \dfrac{m+3-9}{3} = \dfrac{m - 6}{3}$. Donc
\[
S = \left\{ \left( \frac{m+3}{3}\ ;\ \frac{m-6}{3} \right) \right\}.
\]

\textbf{Étape 3 : problème inverse.} On impose $x = 2y$ :
\[
\frac{m+3}{3} = 2\times\frac{m-6}{3} \iff m + 3 = 2m - 12 \iff m = 15.
\]
Vérification : pour $m = 15$, $x = \dfrac{18}{3} = 6$ et $y = \dfrac{9}{3} = 3$, et l'on a bien $x = 2y$. La valeur cherchée est $m = 15$.
\end{exemplebox}

\begin{popfigure}
\begin{center}
\begin{tikzpicture}
\begin{axis}[
  width=0.85\linewidth, height=7cm,
  axis lines=middle, xlabel={$x$}, ylabel={$y$},
  xmin=-1, xmax=7, ymin=-4, ymax=4,
  xtick={1,2,3,4,5,6}, ytick={-3,-2,-1,1,2,3},
  legend style={at={(0.02,0.98)}, anchor=north west, font=\small},
  grid=both, grid style={popline!40},
]
\addplot[popcurve, domain=-1:7, samples=2, popink, very thick]{x-3};
\addlegendentry{$(D') : x - y = 3$}
\addplot[popcurve, domain=-1:7, samples=2, popblue]{(0-x)/2};
\addlegendentry{$(D_0) : x+2y=0$}
\addplot[popcurve, domain=-1:7, samples=2, poporange]{(3-x)/2};
\addlegendentry{$(D_3) : x+2y=3$}
\addplot[popcurve, domain=-1:7, samples=2, popgreen]{(6-x)/2};
\addlegendentry{$(D_6) : x+2y=6$}
\addplot[popcurve, domain=-1:7, samples=2, poppink]{(9-x)/2};
\addlegendentry{$(D_9) : x+2y=9$}
\addplot[only marks, mark=*, mark size=2pt, popdark] coordinates {(2,-1) (3,0) (4,1) (5,2)};
\end{axis}
\end{tikzpicture}
\end{center}
\legende{Les droites $(D_m) : x + 2y = m$ forment une famille de droites parallèles (même coefficient directeur $-\tfrac12$). Chacune coupe la droite fixe $(D') : x - y = 3$ en un point unique (les points noirs, de coordonnées $\left(\tfrac{m+6}{3}\,;\,\tfrac{m-3}{3}\right)$) : la solution du système $\{x+2y=m \,;\, x-y=3\}$ se \emph{déplace} le long de $(D')$ lorsque $m$ varie.}
\end{popfigure}

\courssub{Cas où le déterminant dépend de $m$ : la discussion complète}

La situation devient plus riche lorsque le paramètre apparaît aussi dans les coefficients des inconnues. Le déterminant $D(m)$ est alors une expression dépendant de $m$, et tout se joue sur son annulation :

\begin{itemize}
\item pour les valeurs de $m$ telles que $D(m) \neq 0$ : solution unique (formules de Cramer) ;
\item pour les valeurs de $m$ telles que $D(m) = 0$ : il faut \textbf{reprendre le système à la main} pour chacune de ces valeurs et conclure.
\end{itemize}

\begin{attention}[Le piège classique du Probatoire]
Ne jamais conclure « le système a une solution unique » sans avoir vérifié que $D \neq 0$ ! Lorsque $D$ dépend de $m$, il existe presque toujours des valeurs de $m$ qui annulent $D$. Oublier de discuter ces valeurs séparément, c'est perdre l'essentiel des points de la question. De même, diviser par $m - 1$ sans traiter le cas $m = 1$ est une faute de raisonnement.
\end{attention}

\begin{exoresolu}[Discussion complète selon les valeurs de $m$]
On considère le système $(S_m) : \begin{cases} x + my = 2 \\ mx + y = m \end{cases}$ où $m$ est un paramètre réel.
\begin{enumerate}
\item Calculer le déterminant $D(m)$ du système et le factoriser.
\item Résoudre $(S_m)$ lorsque $D(m) \neq 0$.
\item Discuter les cas $m = 1$ et $m = -1$, puis conclure.
\end{enumerate}
\tcblower
\textbf{1. Déterminant.}
\[
D(m) = \begin{vmatrix} 1 & m \\ m & 1 \end{vmatrix} = 1\times 1 - m\times m = 1 - m^2 = (1 - m)(1 + m).
\]
$D(m) = 0 \iff m = 1$ ou $m = -1$.

\textbf{2. Cas où $m \neq 1$ et $m \neq -1$.} Le système est de Cramer. Les formules donnent :
\[
x = \frac{\begin{vmatrix} 2 & m \\ m & 1 \end{vmatrix}}{1 - m^2} = \frac{2 - m^2}{1 - m^2}, \qquad
y = \frac{\begin{vmatrix} 1 & 2 \\ m & m \end{vmatrix}}{1 - m^2} = \frac{m - 2m}{1 - m^2} = \frac{-m}{1 - m^2} = \frac{m}{m^2 - 1}.
\]
Donc $S = \left\{ \left( \dfrac{2 - m^2}{1 - m^2}\ ;\ \dfrac{m}{m^2 - 1} \right) \right\}$.

\emph{Vérification sur une valeur simple :} pour $m = 0$, le système s'écrit $\begin{cases} x = 2 \\ y = 0 \end{cases}$, et nos formules donnent bien $x = \dfrac{2}{1} = 2$ et $y = 0$.

\textbf{3. Cas particuliers.}

\emph{Cas $m = 1$ :} le système devient $\begin{cases} x + y = 2 \\ x + y = 1 \end{cases}$. Les deux équations sont incompatibles (une même somme ne peut valoir à la fois $2$ et $1$) : $S = \emptyset$.

\emph{Cas $m = -1$ :} le système devient $\begin{cases} x - y = 2 \\ -x + y = -1 \end{cases}$, c'est-à-dire, en multipliant la deuxième équation par $-1$, $\begin{cases} x - y = 2 \\ x - y = 1 \end{cases}$. Là encore, les équations sont incompatibles : $S = \emptyset$.

\textbf{Conclusion :}
\begin{itemize}
\item si $m \in \mathbb{R} \setminus \{-1\,;\,1\}$, solution unique : $S = \left\{ \left( \dfrac{2 - m^2}{1 - m^2}\ ;\ \dfrac{m}{m^2 - 1} \right) \right\}$ ;
\item si $m = 1$ ou $m = -1$ : $S = \emptyset$.
\end{itemize}
\end{exoresolu}

\begin{attention}[Infinité de solutions ou aucune ?]
Quand $D(m) = 0$, les deux issues sont possibles : infinité de solutions (équations proportionnelles, \emph{y compris les seconds membres}) ou ensemble vide (premiers membres proportionnels mais pas les seconds membres). On ne peut pas le deviner : il faut \textbf{substituer la valeur de $m$ dans le système} et comparer les deux équations obtenues. Dans l'exemple ci-dessus, pour $m = 1$, les premiers membres sont identiques ($x + y$) mais les seconds membres ($2$ et $1$) diffèrent : d'où l'impossibilité.
\end{attention}

\courssub{Problèmes inverses : imposer une condition sur la solution}

Un type de question très fréquent au Probatoire consiste à déterminer $m$ pour que la solution $(x\,;\,y)$ vérifie une condition donnée : $x = y$, $x + y = 5$, $x > 0$, etc. La démarche est toujours la même :

\begin{methode}[Problème inverse]
\begin{enumerate}
\item Résoudre le système : obtenir $x(m)$ et $y(m)$ en fonction de $m$.
\item Traduire la condition imposée en une équation (ou inéquation) d'inconnue $m$.
\item Résoudre cette équation, puis \textbf{vérifier} que les valeurs trouvées ne font pas partie des valeurs exclues par la discussion (celles qui annulent $D$).
\end{enumerate}
\end{methode}

\begin{exemplebox}[Condition d'inégalité]
Reprenons le système $\begin{cases} 2x + y = m \\ x - y = 3 \end{cases}$ dont la solution est $x = \dfrac{m+3}{3}$, $y = \dfrac{m-6}{3}$. Déterminons $m$ pour que la solution vérifie $x > 0$ et $y \geq 0$ :
\[
x > 0 \iff m + 3 > 0 \iff m > -3 \ ; \qquad y \geq 0 \iff m - 6 \geq 0 \iff m \geq 6.
\]
Les deux conditions doivent être satisfaites simultanément : on prend l'intersection des deux ensembles de valeurs. La réponse est $m \geq 6$, c'est-à-dire $m \in [6\,;\,+\infty[$.
\end{exemplebox}

\begin{savaistu}[Pourquoi les paramètres sont partout]
Les systèmes à paramètre modélisent toutes les situations où une donnée varie : le prix du carburant à la station Tradex, le tarif du forfait MTN ou Orange, le débit d'eau de la CAMWATER... En économie, on cherche souvent le point d'équilibre entre l'offre et la demande ; quand une taxe $m$ modifie l'équation de l'offre, l'équilibre devient une fonction de $m$, exactement comme dans nos exercices. Savoir discuter selon $m$, c'est savoir prévoir comment l'équilibre se déplace !
\end{savaistu}

\begin{exercice}[Pour s'entraîner]
On considère le système $(S_m) : \begin{cases} mx + y = 1 \\ x + my = m \end{cases}$ où $m \in \mathbb{R}$.
\begin{enumerate}
\item Montrer que $D(m) = m^2 - 1$.
\item Résoudre $(S_m)$ pour $m \notin \{-1\,;\,1\}$.
\item Discuter les cas $m = 1$ et $m = -1$.
\item Déterminer $m$ pour que le système admette une solution vérifiant $x = y$.
\end{enumerate}
\end{exercice}

\begin{corrige}[Exercice 1]
\textbf{1.} $D(m) = \begin{vmatrix} m & 1 \\ 1 & m \end{vmatrix} = m\times m - 1\times 1 = m^2 - 1 = (m-1)(m+1)$.

\textbf{2.} Pour $m \notin \{-1\,;\,1\}$, les formules de Cramer donnent :
\[
x = \frac{\begin{vmatrix} 1 & 1 \\ m & m \end{vmatrix}}{m^2-1} = \frac{m - m}{m^2-1} = 0, \qquad
y = \frac{\begin{vmatrix} m & 1 \\ 1 & m \end{vmatrix}}{m^2-1} = \frac{m^2 - 1}{m^2 - 1} = 1.
\]
Donc $S = \{(0\,;\,1)\}$ : remarquablement, la solution ne dépend pas de $m$ ! Vérification : $m\times 0 + 1 = 1$ et $0 + m\times 1 = m$ : les deux équations sont bien satisfaites.

\textbf{3.} \emph{Cas $m = 1$ :} le système devient $\begin{cases} x + y = 1 \\ x + y = 1 \end{cases}$ : les deux équations sont identiques, donc $S = \{(x\,;\,y) \in \mathbb{R}^2 \mid x + y = 1\}$, une infinité de solutions (tous les points de la droite d'équation $x + y = 1$).

\emph{Cas $m = -1$ :} le système devient $\begin{cases} -x + y = 1 \\ x - y = -1 \end{cases}$, soit deux fois la même équation $y = x + 1$ : $S = \{(x\,;\,x+1) \mid x \in \mathbb{R}\}$, une infinité de solutions.

\textbf{4.} Pour $m \notin \{-1\,;\,1\}$, la solution unique est $(0\,;\,1)$ : on n'a jamais $x = y$. Pour $m = 1$, le couple $\left(\tfrac12\,;\,\tfrac12\right)$ est solution et vérifie $x = y$. Pour $m = -1$, la condition $x = y$ jointe à $y = x + 1$ donnerait $x = x + 1$, impossible : aucune solution ne vérifie $x = y$. Conclusion : la condition $x = y$ n'est réalisable que pour $m = 1$.
\end{corrige}

\begin{aretenir}[Bilan de la section]
\begin{itemize}
\item Un paramètre $m$ se manipule comme un nombre fixe ; on exprime la solution en fonction de $m$.
\item Si le déterminant $D$ ne contient pas $m$ et $D \neq 0$ : solution unique pour tout $m$, donnée par les formules de Cramer.
\item Si $D$ dépend de $m$ : on résout pour $D(m) \neq 0$, puis on traite \textbf{séparément} chaque valeur annulant $D$ en la substituant dans le système (infinité de solutions ou ensemble vide).
\item Pour un problème inverse, on traduit la condition sur $(x\,;\,y)$ en une équation (ou inéquation) en $m$, sans oublier de vérifier la compatibilité avec la discussion.
\end{itemize}
\end{aretenir}

\coursec{Systèmes linéaires dans $\mathbb{R}^3$ : substitution, combinaisons et pivot de Gauss}

Dans la section précédente, tu as appris à discuter un système de deux équations à deux inconnues dont le second membre dépend d'un paramètre : tu as vu combien la méthode du déterminant est puissante dans $\mathbb{R}^2$. Mais la vie courante — et le Probatoire — impose souvent \emph{trois} inconnues : trois prix, trois quantités de ciment, de sable et de gravier pour un chantier à Yaoundé, trois durées de communication sur les réseaux MTN, Orange et Nexttel. Or, le programme officiel est formel : \textbf{le déterminant d'un système $3\times 3$ est hors programme}. Il va donc falloir résoudre autrement, par \cle{élimination}. C'est l'objet de cette section, qui te donnera trois outils complémentaires : la substitution, les combinaisons linéaires, et la reine des méthodes, le \cle{pivot de Gauss}.

\courssub{Équations linéaires à trois inconnues et systèmes $3\times 3$}

\begin{definition}[Équation linéaire à trois inconnues]
Une \cle{équation linéaire à trois inconnues} $x$, $y$, $z$ est une équation de la forme
\[ ax + by + cz = d, \]
où $a$, $b$, $c$, $d$ sont des réels donnés, avec $(a,b,c) \neq (0,0,0)$. Une \cle{solution} de cette équation est un triplet $(x\,;y\,;z)$ de réels qui vérifie l'égalité.
\end{definition}

\begin{definition}[Système linéaire $3\times 3$]
Un \cle{système linéaire de trois équations à trois inconnues} est un système de la forme
\[ (S) : \begin{cases} ax + by + cz = d \\ a'x + b'y + c'z = d' \\ a''x + b''y + c''z = d'' \end{cases} \]
où les coefficients $a, b, c, \dots, d''$ sont des réels donnés. \cle{Résoudre} $(S)$, c'est déterminer l'ensemble $\mathcal{S}$ de tous les triplets $(x\,;y\,;z)$ de $\mathbb{R}^3$ qui vérifient \emph{simultanément} les trois équations.
\end{definition}

\begin{attention}[Le déterminant dans $\mathbb{R}^3$ est hors programme]
Contrairement aux systèmes $2\times 2$, \textbf{aucune formule de Cramer n'est exigible en Première} pour les systèmes $3\times 3$. Toute résolution au Probatoire se fait par \cle{élimination} : substitution, combinaisons linéaires ou pivot de Gauss. N'écris donc jamais de déterminant $3\times 3$ dans une copie de Première.
\end{attention}

\courssub{Interprétation géométrique (admise)}

Dans le plan, une équation $ax+by=c$ représente une droite, et résoudre un système $2\times 2$ revient à chercher l'intersection de deux droites. Dans l'espace, muni d'un repère $(O\,;\vec{i},\vec{j},\vec{k})$, la situation est analogue mais plus riche.

\begin{propriete}[Interprétation géométrique — admise]
Dans un repère de l'espace, l'ensemble des points $M(x\,;y\,;z)$ tels que $ax+by+cz=d$ (avec $(a,b,c)\neq(0,0,0)$) est un \cle{plan}. Résoudre un système $3\times 3$ revient donc à chercher l'\cle{intersection de trois plans} de l'espace.
\end{propriete}

Quatre configurations sont possibles pour l'intersection de trois plans :
\begin{itemize}
\item les trois plans se coupent en \textbf{un unique point} : le système admet une \textbf{solution unique} ;
\item les trois plans se coupent selon \textbf{une droite} : le système admet une \textbf{infinité de solutions} décrites avec un paramètre ;
\item les trois plans sont confondus : infinité de solutions décrites avec \textbf{deux paramètres} ;
\item les plans n'ont \textbf{aucun point commun} (par exemple deux plans strictement parallèles, ou trois plans formant un « prisme ») : le système est \textbf{incompatible}, $\mathcal{S}=\emptyset$.
\end{itemize}

\begin{popfigure}
\begin{tikzpicture}[scale=1.1]
% Plan 1
\fill[popblue!25,opacity=0.8] (0,0) -- (3.2,0.6) -- (4.2,2.2) -- (1,1.6) -- cycle;
\draw[popblue] (0,0) -- (3.2,0.6) -- (4.2,2.2) -- (1,1.6) -- cycle;
% Plan 2
\fill[poporange!30,opacity=0.8] (0.6,2.4) -- (3.8,2.8) -- (4.2,1.1) -- (1,0.7) -- cycle;
\draw[poporange] (0.6,2.4) -- (3.8,2.8) -- (4.2,1.1) -- (1,0.7) -- cycle;
% Plan 3
\fill[popgreen!25,opacity=0.8] (1.4,-0.3) -- (2.6,3.1) -- (3.6,3.0) -- (2.4,-0.4) -- cycle;
\draw[popgreen] (1.4,-0.3) -- (2.6,3.1) -- (3.6,3.0) -- (2.4,-0.4) -- cycle;
% Point commun
\fill[popink] (2.35,1.45) circle (2.2pt);
\node[popink,right] at (2.45,1.45) {solution unique};
\node at (5.6,1.4) {\small $\mathcal{S}=\{(x_0\,;y_0\,;z_0)\}$};
\end{tikzpicture}
\hfill
\begin{tikzpicture}[scale=1.1]
% Deux plans parallèles
\fill[popblue!25,opacity=0.8] (0,0) -- (3.2,0.6) -- (4.2,2.2) -- (1,1.6) -- cycle;
\draw[popblue] (0,0) -- (3.2,0.6) -- (4.2,2.2) -- (1,1.6) -- cycle;
\fill[poporange!30,opacity=0.8] (0,1.4) -- (3.2,2.0) -- (4.2,3.6) -- (1,3.0) -- cycle;
\draw[poporange] (0,1.4) -- (3.2,2.0) -- (4.2,3.6) -- (1,3.0) -- cycle;
% Plan transversal
\fill[popgreen!25,opacity=0.8] (1.4,-0.3) -- (2.6,3.9) -- (3.6,3.8) -- (2.4,-0.4) -- cycle;
\draw[popgreen] (1.4,-0.3) -- (2.6,3.9) -- (3.6,3.8) -- (2.4,-0.4) -- cycle;
\node[popink] at (2.1,-0.8) {aucun point commun};
\node at (5.6,1.4) {\small $\mathcal{S}=\emptyset$};
\end{tikzpicture}
\legende{À gauche : trois plans se coupant en un point unique (solution unique). À droite : deux plans strictement parallèles, le système est incompatible.}
\end{popfigure}

\courssub{La méthode de substitution en cascade}

L'idée est identique à celle du cas $2\times 2$, mais elle se fait « en cascade » : on descend d'un étage à la fois.

\begin{methode}[Substitution en cascade]
\begin{enumerate}
\item Dans une équation, \textbf{isoler} une inconnue (choisir celle dont le coefficient vaut $1$ ou $-1$ pour éviter les fractions), par exemple $x = d - by - cz$.
\item \textbf{Reporter} cette expression dans les deux autres équations : on obtient un système de \textbf{deux équations à deux inconnues} $y$ et $z$.
\item \textbf{Résoudre} ce système $2\times 2$ par les méthodes connues (substitution, combinaison, déterminant).
\item \textbf{Remonter} : remplacer les valeurs trouvées dans l'expression de l'inconnue isolée.
\end{enumerate}
\end{methode}

\begin{exemplebox}[Substitution en cascade]
Résolvons $(S) : \begin{cases} x + y + z = 6 & (L_1)\\ x - y + 2z = 3 & (L_2)\\ 2x + y - z = 3 & (L_3) \end{cases}$

$(L_1)$ donne $x = 6 - y - z$. On reporte dans $(L_2)$ et $(L_3)$ :
\begin{align*}
(L_2) &: (6-y-z) - y + 2z = 3 \iff -2y + z = -3,\\
(L_3) &: 2(6-y-z) + y - z = 3 \iff -y - 3z = -9.
\end{align*}
On résout le système $2\times 2$ : $\begin{cases} -2y + z = -3 \\ -y - 3z = -9 \end{cases}$.
De la première équation, $z = 2y - 3$ ; en reportant dans la deuxième : $-y -3(2y-3) = -9 \iff -7y + 9 = -9 \iff -7y = -18 \iff y = \dfrac{18}{7}$. Les fractions apparaissent ! C'est le défaut de la substitution quand les coefficients ne s'y prêtent pas. Terminons : $z = 2\times\dfrac{18}{7} - 3 = \dfrac{36-21}{7} = \dfrac{15}{7}$, puis $x = 6 - \dfrac{18}{7} - \dfrac{15}{7} = \dfrac{42-33}{7} = \dfrac{9}{7}$.
\[
\mathcal{S} = \left\{\left(\tfrac{9}{7}\,;\tfrac{18}{7}\,;\tfrac{15}{7}\right)\right\}.
\]
Vérification dans $(L_2)$ : $\dfrac{9}{7} - \dfrac{18}{7} + \dfrac{30}{7} = \dfrac{21}{7} = 3$. Exact.
\end{exemplebox}

La substitution fonctionne toujours, mais elle peut traîner des fractions. Les combinaisons linéaires, elles, évitent souvent ce désagrément.

\courssub{La méthode des combinaisons linéaires}

Le principe repose sur un fait simple : certaines transformations laissent l'ensemble des solutions inchangé. Ce sont les \cle{opérations élémentaires} sur les lignes.

\begin{propriete}[Opérations élémentaires sur les lignes — admise]
On ne change pas l'ensemble des solutions d'un système si l'on effectue l'une des opérations suivantes :
\begin{itemize}
\item \textbf{échanger} deux équations : $(L_i) \leftrightarrow (L_j)$ ;
\item \textbf{multiplier} une équation par un réel non nul : $(L_i) \leftarrow k(L_i)$ avec $k \neq 0$ ;
\item \textbf{remplacer} une équation $(L_i)$ par une combinaison $k(L_i) + k'(L_j)$ avec $k \neq 0$ et $j \neq i$ : $(L_i) \leftarrow k(L_i) + k'(L_j)$.
\end{itemize}
\end{propriete}

\begin{methode}[Combinaisons linéaires]
\begin{enumerate}
\item Choisir une inconnue à éliminer (par exemple $x$).
\item Combiner les équations deux à deux pour faire disparaître cette inconnue : on obtient un système $2\times 2$ en $y$ et $z$.
\item Résoudre ce système $2\times 2$ par combinaison ou déterminant.
\item Reporter les valeurs trouvées dans une équation du départ pour obtenir la dernière inconnue.
\end{enumerate}
\end{methode}

\begin{exemplebox}[Combinaisons linéaires]
Résolvons $(S) : \begin{cases} x + y + z = 6 & (L_1)\\ 2x - y + z = 3 & (L_2)\\ x + 2y - z = 2 & (L_3) \end{cases}$

Éliminons $x$ entre $(L_1)$ et $(L_2)$, puis entre $(L_1)$ et $(L_3)$ :
\begin{align*}
(L_2) - 2(L_1) &: (2x-y+z) - (2x+2y+2z) = 3 - 12 \iff -3y - z = -9,\\
(L_3) - (L_1) &: (x+2y-z) - (x+y+z) = 2 - 6 \iff y - 2z = -4.
\end{align*}
On résout $\begin{cases} -3y - z = -9 \\ y - 2z = -4 \end{cases}$. Multiplions la deuxième équation par $3$ : $3y - 6z = -12$, puis ajoutons membre à membre à la première : $-7z = -21$, donc $z = 3$. Alors $y = 2z - 4 = 2$, et $(L_1)$ donne $x = 6 - 2 - 3 = 1$.
\[
\mathcal{S} = \{(1\,;2\,;3)\}.
\]
\end{exemplebox}

\courssub{La méthode du pivot de Gauss}

Le pivot de Gauss n'est rien d'autre que la méthode des combinaisons linéaires \emph{organisée de façon systématique} : au lieu de choisir les combinaisons au feeling, on procède toujours de la même manière, ce qui rend la méthode sûre, rapide et programmable sur machine.

\begin{definition}[Système triangulaire (échelonné)]
Un système $3\times 3$ est dit \cle{triangulaire} (ou \cle{échelonné}) s'il a la forme
\[
\begin{cases} ax + by + cz = d \\ \phantom{ax +{}} b'y + c'z = d' \\ \phantom{ax + by +{}} c''z = d'' \end{cases} \quad \text{avec } a \neq 0,\ b' \neq 0,\ c'' \neq 0.
\]
Un tel système se résout par \cle{remontée} : la dernière équation donne $z$, la deuxième donne alors $y$, la première donne enfin $x$. Il admet une solution unique.
\end{definition}

\begin{methode}[Pivot de Gauss en 4 étapes]
\begin{enumerate}
\item \textbf{Choisir le pivot} : dans la première colonne, choisir un coefficient non nul (de préférence $1$ ou $-1$) ; la ligne qui le contient est la \cle{ligne pivot}, notée $(L_1)$. On peut échanger deux lignes pour placer un bon pivot en haut.
\item \textbf{Éliminer} : sans toucher à la ligne pivot, remplacer chaque autre ligne $(L_i)$ par une combinaison du type $(L_i) \leftarrow \alpha(L_i) + \beta(L_1)$ qui annule le coefficient de $x$. On obtient un bloc de deux équations en $y$ et $z$.
\item \textbf{Répéter} : sur ce bloc $2\times 2$, choisir un nouveau pivot (coefficient de $y$ non nul) et éliminer $y$ dans la dernière ligne. Le système est maintenant triangulaire.
\item \textbf{Remonter} : résoudre de bas en haut : $z$, puis $y$, puis $x$.
\end{enumerate}
\end{methode}

\begin{popfigure}
\begin{tikzpicture}[scale=0.9]
% Etape 1 : systeme plein
\node at (1.5,3.6) {\small Système de départ};
\draw[popmuted] (0,0) rectangle (3,3);
\foreach \i in {0.5,1.5,2.5} {\foreach \j in {0.5,1.5,2.5} {\fill[popblue!60] (\i,\j) circle (4pt);}}
\draw[->,popink,thick] (3.4,1.5) -- (4.4,1.5) node[midway,above]{\small éliminer $x$};
% Etape 2
\begin{scope}[xshift=4.8cm]
\draw[popmuted] (0,0) rectangle (3,3);
\foreach \j in {0.5,1.5,2.5} {\fill[popblue!60] (0.5,\j) circle (4pt);}
\foreach \j in {0.5,1.5} {\fill[poporange!80] (1.5,\j) circle (4pt); \fill[poporange!80] (2.5,\j) circle (4pt);}
\draw[popgreen,thick] (0.2,2.2) rectangle (0.8,2.8);
\node[popgreen] at (0.5,3.25) {\small pivot};
\end{scope}
\draw[->,popink,thick] (8.2,1.5) -- (9.2,1.5) node[midway,above]{\small éliminer $y$};
% Etape 3 : triangulaire
\begin{scope}[xshift=9.6cm]
\draw[popmuted] (0,0) rectangle (3,3);
\foreach \j in {0.5,1.5,2.5} {\fill[popblue!60] (0.5,\j) circle (4pt);}
\foreach \j in {0.5,1.5} {\fill[poporange!80] (1.5,\j) circle (4pt);}
\fill[poppink!90] (2.5,0.5) circle (4pt);
\draw[popgreen,thick] (0.2,2.2) rectangle (0.8,2.8);
\draw[popgreen,thick] (1.2,1.2) rectangle (1.8,1.8);
\draw[popgreen,thick] (2.2,0.2) rectangle (2.8,0.8);
\node at (1.5,-0.5) {\small système triangulaire};
\end{scope}
\end{tikzpicture}
\legende{Structure en escalier visée par le pivot de Gauss : chaque point représente un coefficient non nul, les pivots sont encadrés en vert. On « vide » progressivement le bas de la première colonne, puis de la deuxième.}
\end{popfigure}

\begin{exoresolu}[Résolution complète par pivot de Gauss]
Résoudre dans $\mathbb{R}^3$ le système $(S) : \begin{cases} x + y + z = 6 & (L_1)\\ 2x - y + z = 3 & (L_2)\\ x + 2y - z = 2 & (L_3) \end{cases}$
\tcblower
\textbf{Étape 1 — Choix du pivot.} Le coefficient de $x$ dans $(L_1)$ vaut $1$ : c'est un pivot idéal. La ligne pivot est $(L_1)$, on n'y touche plus.

\textbf{Étape 2 — Élimination de $x$ dans $(L_2)$ et $(L_3)$.}
\[
(L_2) \leftarrow (L_2) - 2(L_1) : \quad -3y - z = -9,
\qquad
(L_3) \leftarrow (L_3) - (L_1) : \quad y - 2z = -4.
\]
Le système devient :
\[
\begin{cases} x + y + z = 6 & (L_1)\\ -3y - z = -9 & (L_2')\\ y - 2z = -4 & (L_3') \end{cases}
\]

\textbf{Étape 3 — Nouveau pivot et élimination de $y$.} Dans le bloc $\{(L_2'),(L_3')\}$, prenons le coefficient $-3$ de $y$ dans $(L_2')$ comme pivot. On élimine $y$ dans $(L_3')$ :
\[
(L_3') \leftarrow 3(L_3') + (L_2') : \quad 3(y - 2z) + (-3y - z) = 3\times(-4) + (-9) \iff -7z = -21.
\]
Le système est maintenant triangulaire :
\[
\begin{cases} x + y + z = 6 \\ -3y - z = -9 \\ -7z = -21 \end{cases}
\]

\textbf{Étape 4 — Remontée.}
\[
z = \frac{-21}{-7} = 3 ; \qquad -3y = -9 + z = -6 \implies y = 2 ; \qquad x = 6 - y - z = 1.
\]
\textbf{Vérification} : $1+2+3=6$ ; $2-2+3=3$ ; $1+4-3=2$. Tout est exact.
\[
\boxed{\mathcal{S} = \{(1\,;2\,;3)\}}
\]
\end{exoresolu}

\begin{attention}[Ne jamais modifier la ligne pivot pendant l'élimination]
L'erreur la plus fréquente consiste à transformer \emph{aussi} la ligne pivot lors d'une combinaison, par exemple en écrivant $(L_1) \leftarrow (L_1) - (L_2)$ « pour faire joli ». C'est interdit dans l'organisation du pivot : \textbf{la ligne pivot reste intacte}, ce sont les \emph{autres} lignes qu'on modifie grâce à elle. Si tu modifies le pivot, tu perds le zéro que tu venais de créer et le système ne devient jamais triangulaire. Autre piège : dans $(L_i) \leftarrow \alpha(L_i) + \beta(L_{\text{pivot}})$, il faut multiplier \emph{toute} la ligne, y compris le second membre.
\end{attention}

\courssub{Cas particuliers : systèmes incompatibles et indéterminés}

Pendant l'élimination, il peut arriver qu'une ligne se transforme en une égalité ne contenant plus aucune inconnue. Deux cas se présentent.

\begin{propriete}[Lecture des cas particuliers]
Au cours de la résolution d'un système $3\times 3$ :
\begin{itemize}
\item si une ligne devient $0 = k$ avec $k \neq 0$, le système est \cle{incompatible} : $\mathcal{S} = \emptyset$ ;
\item si une ligne devient $0 = 0$, cette équation disparaît : le système équivaut à un système de deux équations à trois inconnues, qui admet une \cle{infinité de solutions} que l'on décrit à l'aide d'un paramètre (deux lignes $0=0$ : deux paramètres).
\end{itemize}
\end{propriete}

\begin{exemplebox}[Système incompatible]
$(S) : \begin{cases} x + y + z = 1 & (L_1)\\ 2x + 2y + 2z = 5 & (L_2)\\ x - y + z = 0 & (L_3) \end{cases}$

$(L_2) \leftarrow (L_2) - 2(L_1)$ donne $0x + 0y + 0z = 3$, c'est-à-dire $0 = 3$ : absurde. Donc $\mathcal{S} = \emptyset$. Géométriquement, les deux premiers plans sont strictement parallèles.
\end{exemplebox}

\begin{exemplebox}[Système indéterminé : infinité de solutions]
$(S) : \begin{cases} x + y + z = 3 & (L_1)\\ 2x + 2y + 2z = 6 & (L_2)\\ x - y + z = 1 & (L_3) \end{cases}$

$(L_2) \leftarrow (L_2) - 2(L_1)$ donne $0 = 0$ : l'équation $(L_2)$ était un multiple de $(L_1)$ et disparaît. Il reste $\begin{cases} x + y + z = 3 \\ x - y + z = 1 \end{cases}$. Par soustraction membre à membre, $2y = 2$ donc $y = 1$. Posons alors $z = t$ ($t \in \mathbb{R}$) : la première équation donne $x = 3 - 1 - t = 2 - t$. Donc
\[
\mathcal{S} = \{(2 - t\,;\,1\,;\,t),\ t \in \mathbb{R}\}.
\]
Géométriquement, les solutions forment une droite de l'espace, intersection de deux plans distincts.
\end{exemplebox}

\courssub{Complément pour la Terminale : l'écriture matricielle (hors programme strict)}

\begin{savaistu}[Un système peut s'écrire $AX = B$]
On peut condenser un système $3\times 3$ en une seule égalité \cle{matricielle} :
\[
\underbrace{\begin{pmatrix} a & b & c \\ a' & b' & c' \\ a'' & b'' & c'' \end{pmatrix}}_{A}
\underbrace{\begin{pmatrix} x \\ y \\ z \end{pmatrix}}_{X}
=
\underbrace{\begin{pmatrix} d \\ d' \\ d'' \end{pmatrix}}_{B}
\qquad\text{c'est-à-dire}\qquad AX = B.
\]
Par exemple, notre système fil rouge s'écrit $\begin{pmatrix} 1 & 1 & 1 \\ 2 & -1 & 1 \\ 1 & 2 & -1 \end{pmatrix}\begin{pmatrix} x \\ y \\ z \end{pmatrix} = \begin{pmatrix} 6 \\ 3 \\ 2 \end{pmatrix}$. Cette écriture, étudiée en détail en Terminale, est celle qu'utilisent les ordinateurs : pour eux, résoudre un système, c'est manipuler la matrice $A$.
\end{savaistu}

\begin{savaistu}[Gauss, le prince des mathématiciens]
La méthode du pivot porte le nom de \cle{Carl Friedrich Gauss} (1777--1855), mathématicien allemand surnommé « le prince des mathématiciens », qui la systématisa pour ses calculs d'astronomie (elle était connue des Chinois dix-huit siècles plus tôt !). Aujourd'hui encore, elle est au cœur des logiciels de calcul scientifique : les prévisions météorologiques, qui simulent l'atmosphère du globe — y compris les pluies de la saison humide au Cameroun —, reviennent à résoudre des systèmes de \emph{millions} d'équations par des variantes du pivot de Gauss.
\end{savaistu}

\begin{exercice}[À toi de jouer]
Résoudre dans $\mathbb{R}^3$ par pivot de Gauss le système :
\[
\begin{cases} x + 2y + z = 8 \\ 2x - y + 3z = 9 \\ 3x + y - z = 2 \end{cases}
\]
\end{exercice}

\begin{corrige}[Exercice 1]
Pivot $1$ en $(L_1)$. $(L_2) \leftarrow (L_2) - 2(L_1)$ : $-5y + z = -7$. $(L_3) \leftarrow (L_3) - 3(L_1)$ : $-5y - 4z = -22$. Nouveau pivot $-5$ devant $y$ : $(L_3') \leftarrow (L_3') - (L_2')$ : $-5z = -15$, donc $z = 3$. Remontée : $-5y = -7 - z = -10$ donc $y = 2$ ; puis $x = 8 - 2y - z = 8 - 4 - 3 = 1$. Vérification : $1 + 4 + 3 = 8$ ; $2 - 2 + 9 = 9$ ; $3 + 2 - 3 = 2$. Donc $\mathcal{S} = \{(1\,;2\,;3)\}$.
\end{corrige}

\begin{aretenir}[L'essentiel de la section]
\begin{itemize}
\item Un système linéaire $3\times 3$ représente l'intersection de trois plans de l'espace : solution unique, infinité de solutions, ou ensemble vide.
\item Le déterminant $3\times 3$ est \textbf{hors programme} : on résout par élimination (substitution, combinaisons, pivot de Gauss).
\item Le pivot de Gauss : choisir un pivot non nul, éliminer cette inconnue dans les autres lignes \emph{sans toucher à la ligne pivot}, répéter pour obtenir un système triangulaire, puis remonter.
\item Une ligne $0 = k$ ($k\neq 0$) signifie $\mathcal{S} = \emptyset$ ; une ligne $0 = 0$ signifie une infinité de solutions, décrites avec un paramètre.
\item Toujours \textbf{vérifier} la solution en la reportant dans les trois équations de départ.
\end{itemize}
\end{aretenir}

\coursec{Résolution de problèmes concrets par les systèmes}

Dans la section précédente, tu as appris à résoudre des systèmes linéaires dans $\mathbb{R}^3$ par substitution, par combinaisons linéaires et par le pivot de Gauss. Tu disposes donc maintenant de toute la « boîte à outils » technique. Mais à quoi servent ces outils dans la vie réelle ? C'est précisément l'objet de cette dernière section : transformer une situation concrète — un marché à Mokolo, un tarif de transport, un partage d'héritage — en un système d'équations, le résoudre, puis revenir à la réalité avec une réponse claire. C'est cette compétence de \cle{modélisation} qui est évaluée au Probatoire dans les problèmes de la vie courante.

\courssub{La démarche de modélisation en quatre étapes}

\textbf{Intuition.} Un problème concret est écrit en français, avec des phrases, des prix, des personnes. Un système d'équations est écrit en langage mathématique, avec des lettres et des égalités. Modéliser, c'est \emph{traduire} d'une langue à l'autre, résoudre dans la langue mathématique, puis \emph{retraduire} la réponse en français.

\begin{methode}[Les 4 étapes de la modélisation d'un problème concret]
\begin{enumerate}
\item \textbf{Choix des inconnues.} Lis attentivement l'énoncé et identifie les quantités inconnues. Pose des lettres explicites : $x$ = prix d'un ticket adulte (en FCFA), $y$ = prix d'un ticket enfant (en FCFA), etc. Précise toujours l'\cle{unité} et l'\cle{ensemble} auquel appartient chaque inconnue ($\mathbb{N}$, $\mathbb{R}_+$, ...).
\item \textbf{Mise en équations.} Traduis chaque information de l'énoncé par une équation. Il faut \emph{autant d'équations indépendantes que d'inconnues} pour espérer une solution unique.
\item \textbf{Résolution.} Résous le système par la méthode de ton choix : substitution, combinaisons linéaires, déterminant (dans $\mathbb{R}^2$) ou pivot de Gauss (dans $\mathbb{R}^3$).
\item \textbf{Interprétation et vérification.} Reviens au contexte : la solution est-elle positive, entière si nécessaire, compatible avec l'énoncé ? Vérifie en remplaçant dans les données de départ, puis rédige une phrase de conclusion avec les unités.
\end{enumerate}
\end{methode}

\begin{popfigure}
\begin{center}
\begin{tikzpicture}[node distance=0.9cm and 1.6cm,
etape/.style={rectangle, rounded corners=4pt, draw=popblue, fill=popblueL, thick, minimum height=1cm, text width=3.1cm, align=center, font=\small},
fleche/.style={-{Stealth[length=3mm]}, thick, popdark}]
\node[etape, align=center] (sit) {\textbf{Situation réelle}\\ (énoncé en français)};
\node[etape, right=of sit, align=center] (inc) {\textbf{1. Inconnues}\\ $x$, $y$, $z$ + unités};
\node[etape, right=of inc, align=center] (sys) {\textbf{2. Système}\\ traduction en équations};
\node[etape, below=of sys, align=center] (sol) {\textbf{3. Solution}\\ résolution algébrique};
\node[etape, left=of sol, align=center] (int) {\textbf{4. Interprétation}\\ vérification + conclusion};
\draw[fleche] (sit) -- (inc);
\draw[fleche] (inc) -- (sys);
\draw[fleche] (sys) -- (sol);
\draw[fleche] (sol) -- (int);
\draw[fleche, poporange, dashed] (int.north) -- node[left, font=\small\itshape, text=poporange]{incohérence : revoir le modèle} (sit.south);
\draw[fleche, popgreen] (int.west) to[bend right=20] node[below, font=\small\itshape, text=popgreen]{solution cohérente : réponse finale} ($(sit.south west)+(0,-0.4)$);
\end{tikzpicture}
\end{center}
\legende{La boucle de modélisation : si la solution mathématique est incohérente avec la réalité, on revient à l'énoncé pour corriger la traduction.}
\end{popfigure}

\begin{attention}[Une solution mathématique n'est pas toujours une solution du problème]
Le système peut avoir une solution unique $(x\,;y)$ parfaitement correcte algébriquement, mais \textbf{inacceptable concrètement} :
\begin{itemize}
\item un prix ou une quantité \textbf{négatif} ($x=-300$ FCFA) ;
\item une \textbf{fraction de personne} ou d'objet indivisible ($y=3{,}5$ élèves, $z=\tfrac{7}{2}$ poules) ;
\item une valeur qui contredit une donnée de l'énoncé.
\end{itemize}
Dans ce cas, il y a une \cle{erreur de modélisation} : équation mal traduite, donnée mal recopiée. On revient à l'étape 1 ou 2. Au Probatoire, conclure « le prix est de $-300$ FCFA » sans commentaire coûte des points !
\end{attention}

\courssub{Problèmes de commerce et de tarification (deux inconnues)}

\textbf{Intuition.} Au marché, on connaît souvent le montant total payé pour deux achats différents, mais pas le prix unitaire de chaque article. Deux achats = deux équations ; deux prix inconnus = deux inconnues. C'est un système $2\times 2$.

\begin{exoresolu}[Tarifs au cinéma — retour sur la situation d'accroche]
\textbf{Énoncé.} À l'entrée d'une salle de spectacle à Yaoundé, un groupe de 2 adultes et 3 enfants paie \SI{5800}{FCFA}. Un autre groupe de 1 adulte et 4 enfants paie \SI{5400}{FCFA}. Quel est le prix d'un ticket adulte et d'un ticket enfant ?
\tcblower
\textbf{Solution détaillée.}

\emph{Étape 1 — Inconnues.} Soit $x$ le prix d'un ticket adulte et $y$ le prix d'un ticket enfant, en FCFA ($x>0$, $y>0$).

\emph{Étape 2 — Mise en équations.}
\[
\begin{cases} 2x+3y=5800 & (L_1)\\ x+4y=5400 & (L_2) \end{cases}
\]

\emph{Étape 3 — Résolution} par combinaisons linéaires : $2L_2-L_1$ donne
\[2(x+4y)-(2x+3y)=2\times 5400-5800 \iff 5y=5000 \iff y=1000.\]
En reportant dans $(L_2)$ : $x=5400-4\times 1000=1400$.

Vérifions le déterminant : $\Delta=2\times 4-3\times 1=5\neq 0$, la solution est bien unique.

\emph{Étape 4 — Vérification et conclusion.} $2\times 1400+3\times 1000=2800+3000=5800$ \checkmark\ ; $1400+4\times 1000=5400$ \checkmark. Les valeurs sont positives et plausibles.

\textbf{Conclusion : le ticket adulte coûte \SI{1400}{FCFA} et le ticket enfant \SI{1000}{FCFA}.}
\end{exoresolu}

\begin{attention}[Erreur fréquente : mélanger les lignes]
Dans l'énoncé « 2 adultes et 3 enfants paient 5800 », c'est bien $2x+3y=5800$ et non $3x+2y=5800$. Prends l'habitude de souligner dans l'énoncé ce qui correspond à $x$ et ce qui correspond à $y$ avant d'écrire les équations.
\end{attention}

\begin{exemplebox}[Forfaits téléphoniques]
Un opérateur propose deux services : les appels facturés $a$ FCFA la minute et les SMS à $b$ FCFA l'unité. Amina consomme 10 minutes et 20 SMS pour \SI{1500}{FCFA} ; Bouba consomme 15 minutes et 5 SMS pour \SI{1375}{FCFA}. On obtient
\[\begin{cases} 10a+20b=1500\\ 15a+5b=1375 \end{cases}\]
$\Delta=10\times 5-20\times 15=-250$. Par les formules de Cramer : $a=\dfrac{1500\times 5-20\times 1375}{-250}=\dfrac{-20000}{-250}=80$ et $b=\dfrac{10\times 1375-1500\times 15}{-250}=\dfrac{-8750}{-250}=35$. La minute coûte \SI{80}{FCFA} et le SMS \SI{35}{FCFA}.
\end{exemplebox}

\courssub{Problèmes de partage, d'âges et de moyennes}

\textbf{Intuition.} « La somme de deux nombres vaut..., leur différence vaut... », « le père a trois fois l'âge de son fils, et dans 10 ans... » : chaque phrase de comparaison donne une équation. Attention aux décalages de temps : dans $n$ années, \emph{tout le monde} vieillit de $n$ ans.

\begin{exoresolu}[Problème d'âges]
\textbf{Énoncé.} Un père a 40 ans et son fils a 10 ans. Dans combien d'années l'âge du père sera-t-il le double de celui du fils ?
\tcblower
\textbf{Solution.} Soit $n$ le nombre d'années cherché ($n\in\mathbb{N}$). Dans $n$ années, le père aura $40+n$ ans et le fils $10+n$ ans. La condition s'écrit :
\[40+n=2(10+n)\iff 40+n=20+2n\iff n=20.\]
Vérification : dans 20 ans, le père aura 60 ans et le fils 30 ans, et $60=2\times 30$ \checkmark.

\textbf{Conclusion : c'est dans 20 ans que l'âge du père sera le double de celui du fils.}
\end{exoresolu}

\begin{exemplebox}[Moyenne au Probatoire]
Dans une classe, la moyenne des 25 garçons est 11 et la moyenne des 15 filles est 13. Si $S_g$ et $S_f$ sont les sommes des notes, alors $S_g=25\times 11=275$ et $S_f=15\times 13=195$, d'où la moyenne de la classe : $\dfrac{275+195}{40}=\dfrac{470}{40}=11{,}75$. Les moyennes se traduisent naturellement en équations linéaires sur les sommes.
\end{exemplebox}

\courssub{Problèmes à trois inconnues}

\textbf{Intuition.} Dès qu'une situation met en jeu trois quantités inconnues liées par trois informations indépendantes, on obtient un système $3\times 3$, que l'on résout par substitution ou par le pivot de Gauss (vu à la section précédente). Rappel : le déterminant $3\times 3$ est hors programme.

\begin{exoresolu}[Répartition d'une somme entre trois personnes]
\textbf{Énoncé.} Une prime de \SI{100000}{FCFA} est répartie entre trois employés Alima, Bala et Célestin. Bala reçoit le double d'Alima, et Célestin reçoit \SI{10000}{FCFA} de plus que Bala. Combien reçoit chacun ?
\tcblower
\textbf{Solution détaillée.}

\emph{Étape 1 — Inconnues.} Soient $a$, $b$, $c$ les parts respectives d'Alima, Bala et Célestin, en FCFA.

\emph{Étape 2 — Mise en équations.}
\[\begin{cases} a+b+c=100000 & (L_1)\\ b=2a & (L_2)\\ c=b+10000 & (L_3) \end{cases}\]

\emph{Étape 3 — Résolution} par substitution : on reporte $(L_2)$ et $(L_3)$ dans $(L_1)$ :
\begin{align*}
a+2a+(2a+10000)&=100000\\
5a&=90000\\
a&=18000.
\end{align*}
D'où $b=2\times 18000=36000$ et $c=36000+10000=46000$.

\emph{Étape 4 — Vérification.} $18000+36000+46000=100000$ \checkmark ; $36000=2\times 18000$ \checkmark ; $46000=36000+10000$ \checkmark. Les trois parts sont positives.

\textbf{Conclusion : Alima reçoit \SI{18000}{FCFA}, Bala \SI{36000}{FCFA} et Célestin \SI{46000}{FCFA}.}
\end{exoresolu}

\begin{exemplebox}[Composition d'une ration alimentaire]
Un éleveur de la région de l'Ouest veut préparer \SI{100}{kg} de mélange à partir de maïs, de soja et de son de blé. Les contraintes : la quantité de soja doit être le tiers de celle du maïs, et le son représente \SI{10}{kg} de plus que le soja. En notant $m$, $s$, $b$ les masses en kg :
\[\begin{cases} m+s+b=100\\ s=\dfrac{m}{3}\\ b=s+10 \end{cases}\]
La substitution donne $m+\tfrac{m}{3}+\tfrac{m}{3}+10=100$, soit $\tfrac{5m}{3}=90$, donc $m=54$ kg, $s=18$ kg, $b=28$ kg. Vérification : $54+18+28=100$ \checkmark.
\end{exemplebox}

\begin{methode}[Rédiger la réponse finale (attendu au Probatoire)]
\begin{enumerate}
\item Annonce les inconnues avec leurs unités dès le début de la copie.
\item Justifie chaque équation par une courte phrase (« le premier achat donne... »).
\item Présente la résolution proprement (méthode nommée, calculs alignés).
\item Termine par une \textbf{phrase de conclusion complète, en français, avec les unités} : « Le prix d'un cahier est de \SI{350}{FCFA} et celui d'un stylo de \SI{150}{FCFA}. » Une réponse du type « $x=350$ et $y=150$ » sans phrase est incomplète.
\end{enumerate}
\end{methode}

\begin{exercice}[Transport en commun]
Une agence de voyage Yaoundé--Douala vend des tickets VIP et des tickets classiques. Un départ avec 8 VIP et 20 classiques rapporte \SI{108000}{FCFA} ; un autre avec 5 VIP et 30 classiques rapporte \SI{120000}{FCFA}. Détermine le prix de chaque type de ticket, puis vérifie la cohérence de ta réponse.
\end{exercice}

\begin{corrige}[Exercice : Transport en commun]
Soient $v$ et $c$ les prix (en FCFA) des tickets VIP et classique.
\[\begin{cases} 8v+20c=108000\\ 5v+30c=120000 \end{cases}\]
On simplifie les équations en divisant la première par 4 et la seconde par 5 :
\[\begin{cases} 2v+5c=27000 & (L_1')\\ v+6c=24000 & (L_2') \end{cases}\]
Méthode par substitution : d'après $(L_2')$, $v=24000-6c$.
On reporte dans $(L_1')$ :
\[2(24000-6c)+5c=27000 \iff 48000-12c+5c=27000 \iff -7c=-21000 \iff c=3000.\]
Alors $v=24000-6\times 3000=6000$.

\emph{Vérification :} $8\times 6000+20\times 3000=48000+60000=108000$ \checkmark ; $5\times 6000+30\times 3000=30000+90000=120000$ \checkmark. Les prix sont des entiers positifs, cohérents.

\textbf{Conclusion : le ticket VIP coûte \SI{6000}{FCFA} et le ticket classique \SI{3000}{FCFA}.}
\end{corrige}

\begin{savaistu}[Les systèmes linéaires partout autour de nous]
Les systèmes d'équations linéaires sont l'un des outils les plus utilisés des mathématiques appliquées :
\begin{itemize}
\item en \textbf{économie}, l'équilibre entre l'offre et la demande d'un produit se trouve en résolvant un système : le prix d'équilibre est l'intersection de deux droites ;
\item en \textbf{agriculture}, les zootechniciens composent les rations alimentaires du bétail en résolvant des systèmes (quantités de protéines, de fibres, coût minimal) — c'est le début de la \emph{programmation linéaire}, inventée pendant la Seconde Guerre mondiale ;
\item dans les \textbf{réseaux de transport et de télécommunications} (MTN, Orange, CAMTEL), les flux de données ou de véhicules aux carrefours obéissent à des systèmes linéaires parfois gigantesques, résolus par ordinateur grâce à... la méthode du pivot de Gauss que tu as apprise !
\end{itemize}
\end{savaistu}

\begin{aretenir}[L'essentiel de la section]
\begin{itemize}
\item Modéliser, c'est traduire une situation réelle en système : \textbf{inconnues} $\to$ \textbf{équations} $\to$ \textbf{résolution} $\to$ \textbf{interprétation}.
\item Il faut autant d'équations indépendantes que d'inconnues.
\item La solution doit être \textbf{vérifiée dans l'énoncé} et être \textbf{concrètement acceptable} (positive, entière si le contexte l'exige).
\item La réponse finale est une \textbf{phrase en français avec les unités}.
\end{itemize}
\end{aretenir}

\newpage

\coursec{Activité d'intégration}

\courssub{Objectifs de l'activité}

À l'issue de cette activité d'intégration, tu seras capable de :
\begin{itemize}
\item traduire une situation concrète de la vie courante en un système d'équations linéaires à trois inconnues ;
\item choisir, parmi la substitution, les combinaisons linéaires et le pivot de Gauss, la méthode la plus adaptée à la structure du système ;
\item résoudre complètement un système $3\times 3$ par le pivot de Gauss, en rédigeant chaque étape ;
\item vérifier la cohérence d'une solution (valeurs entières, positives, compatibles avec les contraintes du contexte) ;
\item introduire un paramètre dans un second membre ou un coefficient et discuter l'existence de solutions réalistes selon les valeurs de ce paramètre ;
\item comparer l'efficacité de deux méthodes de résolution sur un même problème.
\end{itemize}

\courssub{Situation}

Le club de mathématiques du lycée bilingue de Yaoundé organise sa traditionnelle fête de fin d'année, qui réunira élèves, enseignants et anciens membres du club. Le bureau du club, dirigé par sa présidente Amina, élève de Première C, dispose d'un budget de \num{150000} FCFA, intégralement prélevé sur la caisse du club alimentée par les cotisations et une subvention de l'APEE.

Après concertation, le bureau décide d'acheter exactement trois types d'articles au marché Mokolo :
\begin{itemize}
\item des \textbf{boissons} (bouteilles de jus et d'eau), à \num{1000} FCFA l'unité ;
\item des \textbf{brochettes} de bœuf, à \num{1500} FCFA l'unité ;
\item des \textbf{galettes} (puff-puff), à \num{2000} FCFA la barquette.
\end{itemize}

Le trésorier du club impose trois contraintes précises :
\begin{enumerate}
\item le nombre total d'articles achetés doit être exactement $120$ ;
\item la dépense totale doit être exactement de \num{150000} FCFA (le budget est juste bouclé) ;
\item l'expérience des années passées montre qu'il faut \textbf{deux fois plus de boissons que de barquettes de galettes} pour que personne ne reste assoiffé.
\end{enumerate}

\textbf{Problème :} combien d'articles de chaque type le bureau doit-il acheter ?

\textbf{Variante proposée par le professeur :} le vendeur de brochettes annonce que le prix unitaire des brochettes pourrait varier selon l'approvisionnement. On note alors $p$ (en FCFA) le prix unitaire d'une brochette. Pour quelles valeurs de $p$ le problème garde-t-il une solution réaliste (nombres d'articles strictement positifs) ?

\courssub{Consignes}

Travail en groupes de quatre élèves. Chaque groupe rédige sa démarche complète sur une feuille qui sera présentée au tableau.

\begin{enumerate}
\item \textbf{Modélisation.} Choisir trois inconnues $x$, $y$, $z$ et préciser clairement ce qu'elles représentent. Traduire chacune des trois contraintes par une équation linéaire. Écrire le système $(S)$ obtenu.
\item \textbf{Résolution.} Résoudre le système $(S)$ par la méthode du \cle{pivot de Gauss}, en détaillant chaque opération élémentaire sur les équations.
\item \textbf{Vérification et interprétation.} Vérifier que la solution trouvée satisfait les trois équations. Les valeurs obtenues sont-elles des entiers positifs ? Conclure par une phrase répondant au problème du bureau du club.
\item \textbf{Variante avec paramètre.} Reprendre le système en remplaçant le prix \num{1500} FCFA des brochettes par un paramètre $p$. Exprimer $x$, $y$ et $z$ en fonction de $p$. Déterminer pour quelles valeurs de $p$ les trois quantités restent strictement positives. Le prix initial de \num{1500} FCFA appartient-il à cet ensemble ?
\item \textbf{Comparaison des méthodes.} Résoudre à nouveau le système initial par \cle{substitution} (exprimer $x$ en fonction de $z$ grâce à la troisième contrainte, puis se ramener à un système $2\times 2$). Quelle méthode te semble la plus rapide ici, et pourquoi ?
\item \textbf{Présentation.} Chaque groupe expose sa démarche en cinq minutes ; la classe établit ensuite une synthèse collective sur le choix de la méthode la plus efficace.
\end{enumerate}

\courssub{Ressources à mobiliser}

\begin{aretenir}[Boîte à outils de la leçon]
\begin{itemize}
\item \textbf{Modélisation :} une contrainte du type « total », « dépense », « deux fois plus de... que... » se traduit par une équation linéaire en les inconnues choisies.
\item \textbf{Pivot de Gauss :} par des opérations élémentaires $L_i \leftarrow L_i - k\,L_j$, on élimine successivement les inconnues pour obtenir un système triangulaire, que l'on résout par remontée.
\item \textbf{Substitution :} on exprime une inconnue en fonction des autres à l'aide d'une équation, puis on reporte dans les équations restantes.
\item \textbf{Système à paramètre :} on résout formellement en fonction du paramètre, puis on traduit les conditions de réalisme (ici $x>0$, $y>0$, $z>0$) par des inéquations sur le paramètre, en veillant au signe des dénominateurs.
\item \textbf{Vérification :} toute solution d'un problème concret doit être contrôlée dans les équations \emph{et} dans le contexte (entiers, positivité).
\end{itemize}
\end{aretenir}

\courssub{Démarche et corrigé commenté}

\begin{exoresolu}[Le budget de la fête du club Maths]
\textbf{Énoncé.} Avec les données de la situation : (1) modéliser le problème par un système de trois équations à trois inconnues ; (2) le résoudre par le pivot de Gauss ; (3) vérifier et interpréter ; (4) discuter la variante où le prix des brochettes est un paramètre $p$ ; (5) comparer avec la méthode de substitution.
\tcblower
\textbf{1. Modélisation.}

Posons : $x$ le nombre de boissons, $y$ le nombre de brochettes, $z$ le nombre de barquettes de galettes.

\begin{itemize}
\item Contrainte de quantité totale : $x+y+z=120$.
\item Contrainte budgétaire (en FCFA) : $1000x+1500y+2000z=150000$, que l'on simplifie en divisant les deux membres par $500$ : $2x+3y+4z=300$.
\item Contrainte de proportion : « deux fois plus de boissons que de galettes » se traduit par $x=2z$, c'est-à-dire $x-2z=0$. (Test mental : s'il y a $10$ barquettes de galettes, il faut $20$ boissons, et on a bien $20=2\times 10$.)
\end{itemize}

Le système à résoudre est donc :
\[
(S)\quad\begin{cases} x+y+z=120 & (L_1)\\ 2x+3y+4z=300 & (L_2)\\ x-2z=0 & (L_3)\end{cases}
\]

\textbf{2. Résolution par le pivot de Gauss.}

\emph{Étape 1 : éliminer $x$ dans $L_2$ et $L_3$ grâce au pivot $x$ de $L_1$.}

$L_2 \leftarrow L_2-2L_1$ : $(2x-2x)+(3y-2y)+(4z-2z)=300-2\times 120$, soit $y+2z=60$.

$L_3 \leftarrow L_3-L_1$ : $(x-x)+(0-y)+(-2z-z)=0-120$, soit $-y-3z=-120$, c'est-à-dire $y+3z=120$.

Le système devient :
\[
\begin{cases} x+y+z=120 & (L_1)\\ y+2z=60 & (L'_2)\\ y+3z=120 & (L'_3)\end{cases}
\]

\emph{Étape 2 : éliminer $y$ dans $L'_3$ grâce au pivot $y$ de $L'_2$.}

$L'_3 \leftarrow L'_3-L'_2$ : $(y-y)+(3z-2z)=120-60$, soit $z=60$.

Le système est maintenant triangulaire :
\[
\begin{cases} x+y+z=120\\ y+2z=60\\ z=60\end{cases}
\]

\emph{Étape 3 : remontée.} De $z=60$, on tire $y=60-2z=60-120=-60$, puis $x=120-y-z=120-(-60)-60=120$.

Le système $(S)$ admet donc l'unique solution $(x\,;y\,;z)=(120\,;-60\,;60)$.

\textbf{3. Vérification et interprétation.}

Vérifions dans les trois équations initiales :
\begin{itemize}
\item $x+y+z=120-60+60=120$ \checkmark ;
\item $2x+3y+4z=2(120)+3(-60)+4(60)=240-180+240=300$ \checkmark ;
\item $x-2z=120-2(60)=0$ \checkmark.
\end{itemize}

Le triplet $(120\,;-60\,;60)$ est bien l'unique solution \emph{mathématique} du système, \textbf{mais elle n'est pas réaliste} : on ne peut pas acheter $-60$ brochettes ! Les trois contraintes du trésorier sont donc \emph{incompatibles dans la réalité}, bien que le système admette une unique solution. C'est une découverte importante : un système peut avoir une solution mathématique unique sans que le problème concret ait de solution.

Le bureau doit donc assouplir au moins une contrainte. Par exemple, s'il remplace la contrainte « $x=2z$ » par « il faut autant de boissons que de brochettes et de galettes réunies », soit $x=y+z$, le système devient :
\[
\begin{cases} x+y+z=120\\ 2x+3y+4z=300\\ x-y-z=0\end{cases}
\]
La troisième équation donne $x=y+z$ ; en reportant dans la première : $2x=120$, donc $x=60$ et $y+z=60$. La deuxième équation donne alors $120+3y+4z=300$, soit $3y+4z=180$. Avec $y=60-z$ : $3(60-z)+4z=180$, d'où $180+z=180$, c'est-à-dire $z=0$, puis $y=60$. La solution $(60\,;60\,;0)$ est bien positive, mais elle est \emph{dégénérée} pour la fête : elle ne prévoit aucune galette ! En pratique, le bureau décide finalement d'augmenter légèrement le budget ou d'ajuster les quantités : c'est toute la richesse de la discussion en groupe.

\textbf{4. Variante avec paramètre $p$.}

Remplaçons le prix des brochettes par $p$ FCFA ($p>0$). En exprimant la contrainte budgétaire en milliers de francs (on divise par $1000$) et en posant $q=\dfrac{p}{1000}$, le système devient :
\[
(S_p)\quad\begin{cases} x+y+z=120\\ x+q\,y+2z=150\\ x-2z=0\end{cases}
\]

De $(L_3)$ : $x=2z$. En reportant dans $(L_1)$ : $y=120-x-z=120-3z$. En reportant dans $(L_2)$ :
\[
2z+q(120-3z)+2z=150 \iff 120q+(4-3q)z=150.
\]
Pour $4-3q\neq 0$, c'est-à-dire $p\neq \dfrac{4000}{3}$, on obtient :
\[
z=\dfrac{150-120q}{4-3q}=\dfrac{150000-120p}{4000-3p},\qquad x=2z=\dfrac{300000-240p}{4000-3p},
\]
\[
y=120-3z=\dfrac{120(4000-3p)-3(150000-120p)}{4000-3p}=\dfrac{480000-360p-450000+360p}{4000-3p}=\dfrac{30000}{4000-3p}.
\]
\emph{Contrôle :} pour $p=1500$, on a $4000-3p=-500$, d'où $y=\dfrac{30000}{-500}=-60$, $z=\dfrac{150000-180000}{-500}=60$ et $x=120$ : on retrouve exactement le résultat de la question 2.

\emph{Conditions de réalisme :} il faut simultanément $x>0$, $y>0$ et $z>0$.
\begin{itemize}
\item $y>0 \iff \dfrac{30000}{4000-3p}>0 \iff 4000-3p>0 \iff p<\dfrac{4000}{3}\approx 1333,33$ ;
\item sous cette condition, le dénominateur de $z$ est positif, donc $z>0 \iff 150000-120p>0 \iff p<1250$ ;
\item $x=2z>0$ est alors automatiquement vérifié.
\end{itemize}
La condition la plus restrictive est $p<1250$. Le problème garde donc une solution réaliste si et seulement si :
\[
0<p<1250 \text{ (FCFA)}.
\]
(Pour que les quantités soient de plus entières, il faudrait que $p$ rende $z$ entier ; on s'en tient ici à la positivité.) Le prix annoncé de \num{1500} FCFA n'appartient pas à cet intervalle : il est \textbf{trop élevé} pour que les trois contraintes soient simultanément satisfaites. Conclusion concrète et utile : le bureau du club doit négocier le prix des brochettes \emph{strictement en dessous de \num{1250} FCFA}, ou modifier une autre contrainte.

\textbf{5. Comparaison avec la substitution.}

Par substitution : $(L_3)$ donne immédiatement $x=2z$. En reportant dans $(L_1)$ : $2z+y+z=120$, soit $y+3z=120$ ; et dans $(L_2)$ : $2(2z)+3y+4z=300$, soit $3y+8z=300$. On obtient le système $2\times 2$ :
\[
\begin{cases} y+3z=120\\ 3y+8z=300\end{cases}
\]
De la première équation, $y=120-3z$ ; en reportant dans la deuxième : $3(120-3z)+8z=300$, soit $360-9z+8z=300$, d'où $z=60$, puis $y=120-180=-60$ et $x=2z=120$. On retrouve la même solution, plus rapidement, car la contrainte $x=2z$ rend la substitution immédiate.

\emph{Conclusion de la comparaison :} quand une équation exprime directement une inconnue en fonction d'une autre (coefficient $1$), la substitution est la plus efficace ; le pivot de Gauss reste la méthode la plus systématique et la plus sûre lorsqu'aucune équation ne se prête à une substitution simple.
\end{exoresolu}

\courssub{Bilan de l'activité}

Cette activité a permis de mettre en évidence plusieurs compétences essentielles de la leçon :
\begin{itemize}
\item \textbf{La modélisation} est une étape à part entière : le choix des inconnues et la traduction fidèle des contraintes (« deux fois plus de... que... » donne $x=2z$ et non $z=2x$ !) conditionnent toute la suite.
\item \textbf{Le pivot de Gauss} fournit une méthode systématique pour tout système $3\times 3$ : élimination successive, forme triangulaire, remontée.
\item \textbf{Une solution mathématique n'est pas toujours une solution du problème concret} : ici, l'unique solution $(120\,;-60\,;60)$ est mathématiquement exacte mais physiquement impossible. La vérification dans le contexte (entiers positifs) est donc indispensable.
\item \textbf{L'introduction d'un paramètre} transforme le problème en une discussion : résoudre formellement, puis traduire les contraintes de réalisme par des inéquations sur $p$. On a découvert que le problème n'a de sens que pour $0<p<1250$ FCFA.
\item \textbf{Le choix de la méthode} dépend de la structure du système : la substitution est redoutable quand une équation isole une inconnue ; le pivot de Gauss est universel. Un bon élève de Première sait passer de l'une à l'autre.
\end{itemize}

\begin{attention}[Erreurs fréquentes repérées pendant l'activité]
\begin{itemize}
\item Traduire « deux fois plus de boissons que de galettes » par $z=2x$ au lieu de $x=2z$ : toujours tester la relation sur un exemple numérique simple.
\item Oublier de simplifier l'équation budgétaire ($2x+3y+4z=300$ au lieu de $1000x+1500y+2000z=150000$), ce qui alourdit tous les calculs.
\item Conclure « il faut acheter $-60$ brochettes » sans s'alarmer : une quantité négative signifie que les contraintes sont incompatibles dans la réalité, pas que le calcul est fini.
\item Dans la discussion en $p$, oublier que le dénominateur $4000-3p$ peut être négatif, ce qui renverse le sens des inégalités, et oublier d'exclure la valeur $p=\dfrac{4000}{3}$ qui l'annule.
\end{itemize}
\end{attention}

\newpage

\coursec{Méthodes et exercices résolus}

\courssub{Savoir-faire 1 : Résoudre un système dans $\mathbb{R}^2$ par substitution}

\begin{methode}[Méthode de substitution]
Pour résoudre un système $\begin{cases} ax+by=c \\ a'x+b'y=c' \end{cases}$ par substitution :
\begin{enumerate}
\item \textbf{Choisir} l'équation où une inconnue a le coefficient le plus simple (idéalement $1$ ou $-1$) et \textbf{exprimer} cette inconnue en fonction de l'autre. Exemple : $x = \dfrac{c-by}{a}$.
\item \textbf{Substituer} cette expression dans l'autre équation : on obtient une équation à une seule inconnue.
\item \textbf{Résoudre} cette équation pour trouver la valeur de la deuxième inconnue.
\item \textbf{Reporter} cette valeur dans l'expression de l'étape 1 pour obtenir la première inconnue.
\item \textbf{Vérifier} le couple trouvé dans les DEUX équations initiales, puis conclure par $S = \{(x\,;y)\}$.
\end{enumerate}
\emph{Réflexe :} la substitution est la méthode à privilégier dès qu'un coefficient vaut $1$ ou $-1$.
\end{methode}

\begin{exoresolu}[Substitution — type Probatoire]
Résoudre dans $\mathbb{R}^2$ le système : $\begin{cases} 2x - y = 3 \quad (E_1)\\ 3x + 4y = 10 \quad (E_2) \end{cases}$
\tcblower
\textbf{Étape 1.} Dans $(E_1)$, le coefficient de $y$ est $-1$ : on exprime $y$ en fonction de $x$ :
\[ 2x - y = 3 \iff y = 2x - 3. \]
\textbf{Étape 2.} On substitue cette expression dans $(E_2)$ :
\[ 3x + 4(2x - 3) = 10. \]
\textbf{Étape 3.} On résout :
\begin{align*}
3x + 8x - 12 &= 10\\
11x &= 22\\
x &= 2.
\end{align*}
\textbf{Étape 4.} On reporte dans $y = 2x - 3$ : $y = 2\times 2 - 3 = 1$.
\textbf{Étape 5.} Vérification : $2(2) - 1 = 3$ \checkmark{} et $3(2) + 4(1) = 10$ \checkmark.
\textbf{Conclusion :} le système admet un unique couple solution : $S = \{(2\,;1)\}$.
\end{exoresolu}

\courssub{Savoir-faire 2 : Résoudre un système dans $\mathbb{R}^2$ par combinaisons linéaires}

\begin{methode}[Méthode des combinaisons linéaires (addition)]
Pour résoudre $\begin{cases} ax+by=c \quad (E_1)\\ a'x+b'y=c' \quad (E_2) \end{cases}$ :
\begin{enumerate}
\item \textbf{Multiplier} $(E_1)$ et $(E_2)$ par des nombres bien choisis de sorte que les coefficients d'une même inconnue deviennent opposés (on utilise le PPCM des coefficients).
\item \textbf{Additionner membre à membre} les deux nouvelles équations : une inconnue est éliminée.
\item \textbf{Résoudre} l'équation à une inconnue obtenue.
\item \textbf{Recommencer} en éliminant l'autre inconnue (ou reporter la valeur trouvée dans une équation initiale).
\item \textbf{Vérifier} dans les deux équations et conclure.
\end{enumerate}
\emph{Réflexe :} cette méthode est la plus rapide quand aucun coefficient ne vaut $\pm 1$, et c'est celle qui se généralise à $\mathbb{R}^3$ (pivot de Gauss).
\end{methode}

\begin{exoresolu}[Combinaisons linéaires — type Probatoire]
Résoudre dans $\mathbb{R}^2$ : $\begin{cases} 3x + 5y = 1 \quad (E_1)\\ 2x - 7y = 17 \quad (E_2) \end{cases}$
\tcblower
\textbf{Élimination de $x$.} On effectue la combinaison $2\times(E_1) - 3\times(E_2)$, qui fait apparaître $6x$ dans les deux équations :
\begin{align*}
2\times(E_1) : \quad 6x + 10y &= 2\\
3\times(E_2) : \quad 6x - 21y &= 51\\
\text{Soustraction membre à membre : } \quad 31y &= -49 \quad \text{donc} \quad y = -\frac{49}{31}.
\end{align*}
\textbf{Élimination de $y$.} Plutôt que de reporter cette fraction dans une équation (calculs lourds), on élimine l'autre inconnue par la combinaison $7\times(E_1) + 5\times(E_2)$ :
\begin{align*}
7\times(E_1) : \quad 21x + 35y &= 7\\
5\times(E_2) : \quad 10x - 35y &= 85\\
\text{Addition membre à membre : } \quad 31x &= 92 \quad \text{donc} \quad x = \frac{92}{31}.
\end{align*}
\textbf{Vérification :} $3\times\dfrac{92}{31} + 5\times\left(-\dfrac{49}{31}\right) = \dfrac{276 - 245}{31} = \dfrac{31}{31} = 1$ \checkmark ;
$2\times\dfrac{92}{31} - 7\times\left(-\dfrac{49}{31}\right) = \dfrac{184 + 343}{31} = \dfrac{527}{31} = 17$ \checkmark.
\textbf{Conclusion :} $S = \left\{\left(\dfrac{92}{31}\,;-\dfrac{49}{31}\right)\right\}$.
\emph{Remarque :} les solutions fractionnaires sont fréquentes au Probatoire ; il ne faut pas croire s'être trompé parce que le résultat n'est pas entier !
\end{exoresolu}

\courssub{Savoir-faire 3 : Résoudre par le déterminant (formules de Cramer)}

\begin{methode}[Déterminant et formules de Cramer]
Pour le système $\begin{cases} ax+by=c \\ a'x+b'y=c' \end{cases}$, on appelle \cle{déterminant} du système le nombre
\[ \Delta = \begin{vmatrix} a & b \\ a' & b' \end{vmatrix} = ab' - a'b. \]
\begin{enumerate}
\item \textbf{Calculer} $\Delta = ab' - a'b$.
\item \textbf{Discuter} :
\begin{itemize}
\item Si $\Delta \neq 0$ : solution unique donnée par les \cle{formules de Cramer} :
\[ x = \dfrac{\Delta_x}{\Delta} = \dfrac{cb' - c'b}{\Delta}, \qquad y = \dfrac{\Delta_y}{\Delta} = \dfrac{ac' - a'c}{\Delta}. \]
($\Delta_x$ s'obtient en remplaçant la colonne des coefficients de $x$ par le second membre ; de même pour $\Delta_y$.)
\item Si $\Delta = 0$ : le système a soit \textbf{aucune} solution, soit une \textbf{infinité} de solutions (droites strictement parallèles ou confondues). On tranche en comparant les rapports $\dfrac{a}{a'}$, $\dfrac{b}{b'}$, $\dfrac{c}{c'}$.
\end{itemize}
\item \textbf{Vérifier} et conclure.
\end{enumerate}
\emph{Réflexe :} le déterminant est l'outil roi dès que le système contient des lettres ou des paramètres.
\end{methode}

\begin{exoresolu}[Formules de Cramer — type Probatoire]
Résoudre dans $\mathbb{R}^2$ par la méthode du déterminant : $\begin{cases} 5x + 2y = 8 \\ 3x - y = 7 \end{cases}$
\tcblower
\textbf{Étape 1.} Déterminant principal :
\[ \Delta = \begin{vmatrix} 5 & 2 \\ 3 & -1 \end{vmatrix} = 5\times(-1) - 3\times 2 = -5 - 6 = -11. \]
\textbf{Étape 2.} $\Delta = -11 \neq 0$ : le système admet un unique couple solution.
\textbf{Étape 3.} Déterminants partiels :
\[ \Delta_x = \begin{vmatrix} 8 & 2 \\ 7 & -1 \end{vmatrix} = 8\times(-1) - 7\times 2 = -8 - 14 = -22 ; \]
\[ \Delta_y = \begin{vmatrix} 5 & 8 \\ 3 & 7 \end{vmatrix} = 5\times 7 - 3\times 8 = 35 - 24 = 11. \]
\textbf{Étape 4.} Formules de Cramer :
\[ x = \frac{\Delta_x}{\Delta} = \frac{-22}{-11} = 2, \qquad y = \frac{\Delta_y}{\Delta} = \frac{11}{-11} = -1. \]
\textbf{Vérification :} $5(2) + 2(-1) = 8$ \checkmark ; $3(2) - (-1) = 7$ \checkmark.
\textbf{Conclusion :} $S = \{(2\,;-1)\}$.
\end{exoresolu}

\courssub{Savoir-faire 4 : Systèmes de $\mathbb{R}^2$ dont le second membre dépend d'un paramètre}

\begin{methode}[Système à second membre paramétrique]
Pour résoudre $\begin{cases} ax+by=c(m) \\ a'x+b'y=c'(m) \end{cases}$ où $m$ est un paramètre réel :
\begin{enumerate}
\item \textbf{Calculer} le déterminant $\Delta = ab' - a'b$ : il ne dépend \textbf{pas} de $m$ (il est formé uniquement des coefficients des inconnues).
\item Si $\Delta \neq 0$ : \textbf{appliquer les formules de Cramer} ; la solution $(x\,;y)$ s'exprime alors \textbf{en fonction de $m$}.
\item Répondre ensuite aux questions annexes typiques : « déterminer $m$ pour que $x = y$ », « pour que $x > 0$ », etc. — on résout alors une équation ou une inéquation d'inconnue $m$.
\item Si $\Delta = 0$, discuter selon les valeurs de $m$ en comparant les équations (cas peu exigé, mais sachez dire : « droites parallèles ou confondues »).
\end{enumerate}
\emph{Point de vigilance :} le paramètre est au \textbf{second membre} ; ne le mélangez jamais avec les coefficients de $x$ et $y$ dans le calcul de $\Delta$.
\end{methode}

\begin{exoresolu}[Second membre paramétrique — type Probatoire]
On considère le système $(S_m)$ : $\begin{cases} 2x + y = m \\ x - 3y = m + 4 \end{cases}$ où $m$ est un réel.
\begin{enumerate}
\item Résoudre $(S_m)$ en fonction de $m$.
\item Déterminer $m$ pour que le couple solution vérifie $x = y$.
\end{enumerate}
\tcblower
\textbf{1. Résolution en fonction de $m$.}
Déterminant : $\Delta = \begin{vmatrix} 2 & 1 \\ 1 & -3 \end{vmatrix} = 2\times(-3) - 1\times 1 = -7 \neq 0$ : solution unique pour tout réel $m$.
\[ \Delta_x = \begin{vmatrix} m & 1 \\ m+4 & -3 \end{vmatrix} = -3m - (m+4) = -4m - 4 ; \]
\[ \Delta_y = \begin{vmatrix} 2 & m \\ 1 & m+4 \end{vmatrix} = 2(m+4) - m = m + 8. \]
Par les formules de Cramer :
\[ x = \frac{\Delta_x}{\Delta} = \frac{-4m-4}{-7} = \frac{4m+4}{7}, \qquad y = \frac{\Delta_y}{\Delta} = \frac{m+8}{-7} = -\frac{m+8}{7}. \]
Donc $S = \left\{\left(\dfrac{4m+4}{7}\,;-\dfrac{m+8}{7}\right)\right\}$.

\textbf{2. Condition $x = y$.}
\begin{align*}
\frac{4m+4}{7} = -\frac{m+8}{7} &\iff 4m + 4 = -(m+8)\\
&\iff 4m + 4 = -m - 8\\
&\iff 5m = -12 \iff m = -\frac{12}{5}.
\end{align*}
\textbf{Conclusion :} pour $m = -\dfrac{12}{5}$, le couple solution vérifie $x = y$. On obtient alors $x = y = \dfrac{4\times\left(-\frac{12}{5}\right)+4}{7} = \dfrac{-\frac{28}{5}}{7} = -\dfrac{4}{5}$.
\end{exoresolu}

\courssub{Savoir-faire 5 : Résoudre un système dans $\mathbb{R}^3$ (substitution, combinaisons, pivot de Gauss)}

\begin{methode}[Pivot de Gauss dans $\mathbb{R}^3$]
Pour résoudre un système de 3 équations à 3 inconnues :
\begin{enumerate}
\item \textbf{Choisir un pivot} : une équation où le coefficient de $x$ est simple (idéalement $1$). La placer en première ligne si nécessaire.
\item \textbf{Éliminer $x$} des deux autres équations par combinaisons linéaires : $(E_2) \leftarrow a(E_2) - a'(E_1)$ et $(E_3) \leftarrow a(E_3) - a''(E_1)$. On obtient un sous-système de 2 équations en $y$ et $z$.
\item \textbf{Éliminer $y$} entre ces deux équations pour obtenir $z$ directement.
\item \textbf{Remonter} : reporter $z$ dans une équation en $(y\,;z)$ pour trouver $y$, puis $y$ et $z$ dans $(E_1)$ pour trouver $x$ (\emph{remontée triangulaire}).
\item \textbf{Vérifier} le triplet dans les TROIS équations initiales et conclure : $S = \{(x\,;y\,;z)\}$.
\end{enumerate}
\emph{Réflexes :} écrire les opérations effectuées sur les lignes à côté du système (ex. $E_2 \leftarrow 2E_2 - 3E_1$) ; garder les équations alignées par colonnes d'inconnues. Le déterminant $3\times 3$ est \textbf{hors programme} : on ne l'utilise pas.
\end{methode}

\begin{exoresolu}[Pivot de Gauss — type Probatoire]
Résoudre dans $\mathbb{R}^3$ : $\begin{cases} x + 2y - z = 2 \quad (E_1)\\ 2x - y + z = 5 \quad (E_2)\\ 3x + y + 2z = 9 \quad (E_3) \end{cases}$
\tcblower
\textbf{Étape 1.} Le coefficient de $x$ dans $(E_1)$ vaut $1$ : c'est notre pivot.
\textbf{Étape 2.} Élimination de $x$ :
\begin{itemize}
\item $(E_2') \leftarrow (E_2) - 2(E_1)$ : $\quad (2x - y + z) - (2x + 4y - 2z) = 5 - 4$, soit $-5y + 3z = 1$ ;
\item $(E_3') \leftarrow (E_3) - 3(E_1)$ : $\quad (3x + y + 2z) - (3x + 6y - 3z) = 9 - 6$, soit $-5y + 5z = 3$.
\end{itemize}
Système réduit : $\begin{cases} -5y + 3z = 1 \quad (E_2')\\ -5y + 5z = 3 \quad (E_3') \end{cases}$
\textbf{Étape 3.} Élimination de $y$ : $(E_3') - (E_2')$ donne $2z = 2$, donc $z = 1$.
\textbf{Étape 4.} Remontée : dans $(E_2')$ : $-5y + 3(1) = 1 \iff -5y = -2 \iff y = \dfrac{2}{5}$.
Dans $(E_1)$ : $x + 2\times\dfrac{2}{5} - 1 = 2 \iff x = 3 - \dfrac{4}{5} = \dfrac{11}{5}$.
\textbf{Étape 5.} Vérification :
\begin{itemize}
\item $(E_1)$ : $\dfrac{11}{5} + \dfrac{4}{5} - 1 = \dfrac{15}{5} - 1 = 2$ \checkmark
\item $(E_2)$ : $\dfrac{22}{5} - \dfrac{2}{5} + 1 = \dfrac{20}{5} + 1 = 5$ \checkmark
\item $(E_3)$ : $\dfrac{33}{5} + \dfrac{2}{5} + 2 = \dfrac{35}{5} + 2 = 9$ \checkmark
\end{itemize}
\textbf{Conclusion :} $S = \left\{\left(\dfrac{11}{5}\,;\dfrac{2}{5}\,;1\right)\right\}$.
\end{exoresolu}

\courssub{Savoir-faire 6 : Résoudre un problème concret à l'aide d'un système}

\begin{methode}[Mise en équations d'un problème concret]
\begin{enumerate}
\item \textbf{Choisir les inconnues} : nommer clairement ce que l'on cherche (ex. « soit $x$ le prix d'un cahier en FCFA »). Préciser l'ensemble ($\mathbb{N}$, $\mathbb{R}_+$\ldots).
\item \textbf{Traduire chaque information} de l'énoncé en une équation : une phrase = une équation. Relire l'énoncé pour vérifier qu'aucune donnée n'est inutilisée.
\item \textbf{Résoudre} le système par la méthode la plus adaptée.
\item \textbf{Contrôler la cohérence} : les solutions doivent être compatibles avec le contexte (prix positifs, nombres de personnes entiers\ldots).
\item \textbf{Rédiger la réponse} en phrase, avec l'unité, en répondant exactement à la question posée.
\end{enumerate}
\end{methode}

\begin{exoresolu}[Problème concret — contexte camerounais]
Au marché Mokolo, une maman achète $3$ kg de tomates et $2$ kg d'oignons pour \SI{3500}{FCFA}. Le lendemain, aux mêmes prix, elle achète $2$ kg de tomates et $5$ kg d'oignons pour \SI{4900}{FCFA}.
Déterminer le prix d'un kilogramme de tomates et celui d'un kilogramme d'oignons.
\tcblower
\textbf{Étape 1. Choix des inconnues.} Soit $x$ le prix d'un kg de tomates et $y$ celui d'un kg d'oignons (en FCFA, $x > 0$, $y > 0$).
\textbf{Étape 2. Mise en équations.}
\begin{itemize}
\item Premier achat : $3x + 2y = 3500$ ;
\item Deuxième achat : $2x + 5y = 4900$.
\end{itemize}
\textbf{Étape 3. Résolution} par le déterminant :
\[ \Delta = 3\times 5 - 2\times 2 = 15 - 4 = 11 \neq 0. \]
\[ \Delta_x = 3500\times 5 - 4900\times 2 = 17500 - 9800 = 7700 ; \]
\[ \Delta_y = 3\times 4900 - 2\times 3500 = 14700 - 7000 = 7700. \]
\[ x = \frac{\Delta_x}{\Delta} = \frac{7700}{11} = 700, \qquad y = \frac{\Delta_y}{\Delta} = \frac{7700}{11} = 700. \]
\textbf{Étape 4. Vérification :} $3(700) + 2(700) = 2100 + 1400 = 3500$ \checkmark ; $2(700) + 5(700) = 1400 + 3500 = 4900$ \checkmark. Les prix sont positifs : cohérent.
\textbf{Conclusion :} le kilogramme de tomates coûte \SI{700}{FCFA} et le kilogramme d'oignons coûte également \SI{700}{FCFA}.
\end{exoresolu}

\begin{exoresolu}[Problème concret à trois inconnues — type Probatoire]
Une librairie de Yaoundé vend trois articles : un stylo, un cahier et une règle.
\begin{itemize}
\item $1$ stylo, $1$ cahier et $1$ règle coûtent ensemble \SI{1000}{FCFA} ;
\item $2$ stylos et $3$ cahiers coûtent \SI{2000}{FCFA} ;
\item le cahier coûte \SI{100}{FCFA} de plus que la règle.
\end{itemize}
Déterminer le prix de chaque article.
\tcblower
\textbf{Étape 1. Inconnues.} Soient $x$ le prix du stylo, $y$ celui du cahier, $z$ celui de la règle (en FCFA).
\textbf{Étape 2. Mise en équations.}
\[ \begin{cases} x + y + z = 1000 \quad (E_1)\\ 2x + 3y = 2000 \quad (E_2)\\ y - z = 100 \quad (E_3) \end{cases} \]
\textbf{Étape 3. Résolution par substitution.} L'équation $(E_3)$ donne $z = y - 100$, et $(E_2)$ donne $x = \dfrac{2000 - 3y}{2}$. On reporte ces deux expressions dans $(E_1)$ :
\begin{align*}
\frac{2000 - 3y}{2} + y + (y - 100) &= 1000\\
2000 - 3y + 2y + 2y - 200 &= 2000 \quad (\text{on multiplie tout par } 2)\\
y + 1800 &= 2000\\
y &= 200.
\end{align*}
On en déduit : $z = y - 100 = 200 - 100 = 100$ et $x = \dfrac{2000 - 3\times 200}{2} = \dfrac{1400}{2} = 700$.
\textbf{Étape 4. Vérification :}
\begin{itemize}
\item $(E_1)$ : $700 + 200 + 100 = 1000$ \checkmark
\item $(E_2)$ : $2\times 700 + 3\times 200 = 1400 + 600 = 2000$ \checkmark
\item $(E_3)$ : $200 - 100 = 100$ \checkmark
\end{itemize}
Les trois prix sont positifs : cohérent avec le contexte.
\textbf{Conclusion :} le stylo coûte \SI{700}{FCFA}, le cahier \SI{200}{FCFA} et la règle \SI{100}{FCFA}.
\emph{Remarque pédagogique :} la vérification dans les trois équations initiales est le seul juge de paix ; ne la sacrifiez jamais, même sous pression le jour de l'examen.
\end{exoresolu}

\begin{attention}[Les 5 erreurs qui coûtent des points au Probatoire]
\begin{enumerate}
\item \textbf{Oublier de vérifier la solution dans TOUTES les équations.} Une vérification partielle ne prouve rien : le couple (ou triplet) doit satisfaire chaque équation du système \emph{initial}, pas les équations transformées.
\item \textbf{Se tromper dans le signe du déterminant.} $\Delta = ab' - a'b$ : c'est le produit de la diagonale principale \emph{moins} l'autre produit. Écrire $\Delta = ab' + a'b$ ou inverser l'ordre fausse tout le reste. De même, dans $\Delta_x$, on remplace la colonne des coefficients de $x$ par le second membre — pas celle de $y$.
\item \textbf{Diviser par $\Delta$ sans vérifier que $\Delta \neq 0$.} Les formules de Cramer n'ont de sens que si $\Delta \neq 0$. Annoncer d'abord « $\Delta = \ldots \neq 0$, donc le système admet un unique couple solution » : c'est une condition d'application exigée dans la rédaction.
\item \textbf{Faire disparaître le paramètre ou le mélanger aux coefficients.} Dans un système à second membre paramétrique, $m$ n'apparaît \emph{jamais} dans $\Delta$. La solution doit être donnée \emph{en fonction de $m$} : répondre par des nombres quand l'énoncé demande une expression en $m$ fait perdre tous les points de la question.
\item \textbf{Négliger la rédaction finale.} Trois fautes classiques : conclure par « $x = 2$ et $y = 1$ » sans écrire l'ensemble $S = \{(2\,;1)\}$ ; oublier les unités dans un problème concret (FCFA, kg, km) ; ne pas vérifier la cohérence avec le contexte (un nombre de billets négatif ou fractionnaire doit alerter, pas être recopié tel quel).
\end{enumerate}
\end{attention}

\newpage

\coursec{Exercices}

\courssub{Je vérifie mes connaissances}

\begin{exercice}[QCM : réflexes sur les systèmes]
Pour chaque question, une seule réponse est exacte. Entoure-la et justifie brièvement.
\begin{enumerate}
\item Le déterminant du système $\begin{cases} 2x - 3y = 5 \\ 4x + y = 1 \end{cases}$ est :
\begin{enumerate}
\item $-10$ \quad \item $14$ \quad \item $-14$ \quad \item $10$
\end{enumerate}
\item Si le déterminant d'un système de deux équations à deux inconnues est non nul, alors le système :
\begin{enumerate}
\item n'a aucune solution \quad \item a une unique solution \quad \item a une infinité de solutions \quad \item a deux solutions
\end{enumerate}
\item Le couple $(2 ; -1)$ est solution du système :
\begin{enumerate}
\item $\begin{cases} x + y = 1 \\ x - y = 1 \end{cases}$ \quad \item $\begin{cases} x + y = 1 \\ x - y = 3 \end{cases}$ \quad \item $\begin{cases} 2x + y = 5 \\ x + y = 3 \end{cases}$ \quad \item $\begin{cases} x - 2y = 0 \\ x + y = 3 \end{cases}$
\end{enumerate}
\item Géométriquement, un système de deux équations à deux inconnues ayant une infinité de solutions correspond à :
\begin{enumerate}
\item deux droites sécantes \quad \item deux droites strictement parallèles \quad \item deux droites confondues \quad \item deux droites perpendiculaires
\end{enumerate}
\end{enumerate}
\end{exercice}

\begin{exercice}[Vrai ou faux ?]
Réponds par \textbf{vrai} ou \textbf{faux} en justifiant chaque réponse par un argument ou un contre-exemple.
\begin{enumerate}
\item Le système $\begin{cases} x + 2y = 3 \\ 2x + 4y = 6 \end{cases}$ admet une infinité de solutions.
\item Le système $\begin{cases} 3x - y = 2 \\ 6x - 2y = 5 \end{cases}$ admet une unique solution.
\item Tout système linéaire de trois équations à trois inconnues admet au moins une solution.
\item La méthode du pivot de Gauss consiste à éliminer successivement des inconnues pour obtenir un système triangulaire.
\end{enumerate}
\end{exercice}

\begin{exercice}[Déterminants : calculs directs]
Calculer le déterminant de chacun des systèmes suivants et préciser, sans résoudre, si le système admet une unique solution.
\begin{enumerate}
\item $\begin{cases} 5x + 2y = 1 \\ 3x - y = 4 \end{cases}$ \qquad \item $\begin{cases} 4x - 6y = 7 \\ -2x + 3y = 1 \end{cases}$ \qquad \item $\begin{cases} \sqrt{2}\,x + y = 3 \\ x + \sqrt{2}\,y = 0 \end{cases}$
\end{enumerate}
\end{exercice}

\begin{exercice}[Solutions ou pas ?]
Parmi les couples et triplets proposés, indiquer lesquels sont solutions du système donné, en justifiant par un calcul.
\begin{enumerate}
\item Système $\begin{cases} 2x + y = 7 \\ x - 3y = 0 \end{cases}$ ; couples proposés : $(3 ; 1)$, $(1 ; 3)$, $(2 ; 3)$.
\item Système $\begin{cases} x + y + z = 6 \\ x - y + z = 2 \\ 2x + y - z = 1 \end{cases}$ ; triplets proposés : $(1 ; 2 ; 3)$, $(2 ; 1 ; 3)$, $(1 ; 3 ; 2)$.
\end{enumerate}
\end{exercice}

\courssub{Je m'entraîne}

\begin{exercice}[Trois méthodes, un même système]
On considère le système $(S) : \begin{cases} 3x - 2y = 7 \\ 5x + y = 16 \end{cases}$.
\begin{enumerate}
\item Résoudre $(S)$ par la méthode de \cle{substitution}.
\item Résoudre $(S)$ par la méthode des \cle{combinaisons linéaires}.
\item Résoudre $(S)$ par la méthode du \cle{déterminant} (formules de Cramer).
\item Comparer les trois méthodes : laquelle te semble la plus rapide ici ? Justifier.
\end{enumerate}
\end{exercice}

\begin{exercice}[Compter avant de résoudre]
Pour chacun des systèmes suivants, calculer le déterminant, en déduire le nombre de solutions, puis résoudre ceux qui admettent une solution unique.
\begin{enumerate}
\item $\begin{cases} x - 4y = 3 \\ 2x + 3y = 17 \end{cases}$ \qquad \item $\begin{cases} 6x - 9y = 2 \\ 4x - 6y = 5 \end{cases}$ \qquad \item $\begin{cases} 3x + 2y = 16 \\ 5x - y = 5 \end{cases}$
\end{enumerate}
\end{exercice}

\begin{exercice}[Pivot de Gauss dans $\mathbb{R}^3$]
On considère le système $(S) : \begin{cases} x - 2y + z = 1 \\ 2x + y - z = 1 \\ 3x - y + 2z = 6 \end{cases}$.
\begin{enumerate}
\item À l'aide du pivot de Gauss, transformer $(S)$ en un système triangulaire équivalent.
\item Résoudre le système triangulaire obtenu par remontée.
\item Vérifier que le triplet trouvé est bien solution des trois équations de $(S)$.
\end{enumerate}
\end{exercice}

\begin{exercice}[Substitution et combinaison dans $\mathbb{R}^3$]
Résoudre dans $\mathbb{R}^3$ les systèmes suivants, en indiquant clairement la méthode utilisée.
\begin{enumerate}
\item $\begin{cases} x + y + z = 6 \\ x - y = 1 \\ y + z = 5 \end{cases}$ \qquad \item $\begin{cases} 2x + y - z = 1 \\ x - y + z = 2 \\ x + 2y + z = 8 \end{cases}$
\end{enumerate}
\end{exercice}

\courssub{Je me prépare au Probatoire}

\begin{exercice}[La librairie de quartier]
Un commerçant de Douala achète des cahiers de $100$ pages à $500$~F l'unité et des cahiers de $200$ pages à $800$~F l'unité. Il achète au total $50$ cahiers et dépense exactement $\num{31000}$~F. On note $x$ le nombre de cahiers de $100$ pages et $y$ le nombre de cahiers de $200$ pages.
\begin{enumerate}
\item Traduire les données de l'énoncé par un système de deux équations à deux inconnues.
\item Résoudre ce système par la méthode de ton choix et donner le nombre de cahiers de chaque type.
\item Le commerçant revend chaque cahier de $100$ pages à $650$~F et chaque cahier de $200$ pages à $\num{1000}$~F. Calculer son bénéfice total sur cette vente.
\item Pour une commande suivante, il veut acheter $60$ cahiers pour un montant de $\num{36000}$~F, aux mêmes prix d'achat. Montrer, en calculant le déterminant du système correspondant, que ce nouveau problème admet une unique solution, puis la déterminer. Cette solution est-elle réaliste pour le commerçant ?
\end{enumerate}
\end{exercice}

\begin{exercice}[Le bénéfice des trois commerçantes]
Trois commerçantes du marché de Bafoussam, Amina, Brigitte et Chantal, se partagent un bénéfice de $\num{450000}$~F. Amina reçoit le double de la part de Brigitte, et Chantal reçoit $\num{50000}$~F de plus que Brigitte. On note $a$, $b$ et $c$ les parts respectives d'Amina, Brigitte et Chantal, en francs.
\begin{enumerate}
\item Montrer que le problème se traduit par le système $\begin{cases} a + b + c = \num{450000} \\ a - 2b = 0 \\ -b + c = \num{50000} \end{cases}$.
\item Résoudre ce système par la méthode du pivot de Gauss (on détaillera chaque étape de la triangulation).
\item En déduire la part de chaque commerçante.
\item Vérifier que les trois parts trouvées satisfont les trois conditions de l'énoncé.
\item \textbf{Bonus.} Un père a trois fois l'âge de son fils. Dans $12$ ans, il aura le double de l'âge de son fils. En notant $p$ et $f$ leurs âges actuels, traduire la situation par un système et déterminer leurs âges.
\end{enumerate}
\end{exercice}

\begin{exercice}[Discussion avec paramètre]
On considère le système $(S_m) : \begin{cases} mx + 2y = m + 1 \\ 2x + my = 4 \end{cases}$ où $m$ est un paramètre réel.
\begin{enumerate}
\item Calculer le déterminant $D(m)$ du système $(S_m)$ et vérifier que $D(m) = (m - 2)(m + 2)$.
\item Discuter, suivant les valeurs de $m$, le nombre de solutions de $(S_m)$. On étudiera en particulier les cas $m = 2$ et $m = -2$ en remplaçant $m$ par sa valeur dans le système.
\item Résoudre $(S_m)$ pour $m = 3$.
\item On considère maintenant le système $(T_m) : \begin{cases} x + y = m \\ 2x - y = 3 \end{cases}$.
\begin{enumerate}
\item Montrer que $(T_m)$ admet une unique solution pour tout réel $m$, et exprimer $x$ et $y$ en fonction de $m$.
\item Déterminer les valeurs de $m$ pour lesquelles la solution $(x ; y)$ vérifie $x > 0$ et $y > 0$.
\end{enumerate}
\end{enumerate}
\end{exercice}



\newpage

\coursec{Corrigés des exercices}

\begin{corrige}[Exercice 1 — QCM : réflexes sur les systèmes]
\begin{enumerate}
\item Réponse \textbf{b} : $D = 2\times 1 - (-3)\times 4 = 2 + 12 = 14 \neq 0$.
\item Réponse \textbf{b} : c'est le théorème de Cramer : pour un système linéaire $2\times 2$, déterminant non nul $\iff$ solution unique.
\item Réponse \textbf{b} : le couple $(2 ; -1)$ satisfait $\begin{cases} x+y=1 \\ x-y=3 \end{cases}$ car $2+(-1)=1$ et $2-(-1)=3$. Les autres systèmes ne sont pas vérifiés.
\item Réponse \textbf{c} : une infinité de solutions signifie que les deux équations représentent la même droite (droites confondues), donc proportionnalité des coefficients et des seconds membres.
\end{enumerate}
\end{corrige}

\begin{corrige}[Exercice 2 — Vrai ou faux ?]
\begin{enumerate}
\item \textbf{Vrai} : la deuxième équation est le double de la première : $2(x + 2y) = 2\times 3$. Les deux droites sont confondues ; l'ensemble solution est $\{(3-2y ; y) \mid y \in \mathbb{R}\}$.
\item \textbf{Faux} : le déterminant vaut $3\times(-2) - (-1)\times 6 = -6 + 6 = 0$. De plus, $2(3x - y) = 6x - 2y = 4$, incompatible avec $6x - 2y = 5$ : le système n'a \emph{aucune} solution (droites strictement parallèles distinctes).
\item \textbf{Faux} : contre-exemple : $\begin{cases} x + y + z = 1 \\ x + y + z = 2 \\ x - y = 0 \end{cases}$ n'a aucune solution, car les deux premières équations sont incompatibles ($1 \neq 2$).
\item \textbf{Vrai} : c'est exactement le principe du pivot de Gauss : par des opérations élémentaires sur les lignes (transvections), on élimine $x$ de $(E_2)$ et $(E_3)$, puis $y$ de $(E_3)$, ce qui donne un système triangulaire supérieur résolu par remontée.
\end{enumerate}
\end{corrige}

\begin{corrige}[Exercice 3 — Déterminants : calculs directs]
\begin{enumerate}
\item $D = \begin{vmatrix} 5 & 2 \\ 3 & -1 \end{vmatrix} = 5\times(-1) - 2\times 3 = -5 - 6 = -11 \neq 0$ : le système admet une solution unique.
\item $D = \begin{vmatrix} 4 & -6 \\ -2 & 3 \end{vmatrix} = 4\times 3 - (-6)\times(-2) = 12 - 12 = 0$ : pas de solution unique. Selon le second membre, soit les droites sont confondues (infinité de solutions), soit strictement parallèles (aucune solution). Dans le cas de l'exercice, les équations sont incompatibles, donc $\mathcal{S} = \emptyset$.
\item $D = \begin{vmatrix} \sqrt{2} & 1 \\ 1 & \sqrt{2} \end{vmatrix} = \sqrt{2}\times\sqrt{2} - 1\times 1 = 2 - 1 = 1 \neq 0$ : solution unique.
\end{enumerate}
\end{corrige}

\begin{corrige}[Exercice 4 — Solutions ou pas ?]
\begin{enumerate}
\item Système : $\begin{cases} 2x + y = 7 \\ x - 3y = 0 \end{cases}$.
\begin{itemize}
\item $(3 ; 1)$ : $2\times 3 + 1 = 7$ et $3 - 3\times 1 = 0$ : \textbf{solution}.
\item $(1 ; 3)$ : $2\times 1 + 3 = 5 \neq 7$ : rejeté.
\item $(2 ; 3)$ : $2\times 2 + 3 = 7$ mais $2 - 3\times 3 = -7 \neq 0$ : rejeté.
\end{itemize}
\item Système : $\begin{cases} x + y + z = 6 \\ x - y + z = 2 \\ 2x + y - z = 1 \end{cases}$.
\begin{itemize}
\item $(1 ; 2 ; 3)$ : $1+2+3=6$, $1-2+3=2$, $2+2-3=1$ : \textbf{solution}.
\item $(2 ; 1 ; 3)$ : $2-1+3=4 \neq 2$ : rejeté.
\item $(1 ; 3 ; 2)$ : $1-3+2=0 \neq 2$ : rejeté.
\end{itemize}
\end{enumerate}
\end{corrige}

\begin{corrige}[Exercice 5 — Trois méthodes, un même système]
Soit le système : $\begin{cases} 3x - 2y = 7 & (E_1) \\ 5x + y = 16 & (E_2) \end{cases}$

\begin{enumerate}
\item \textbf{Méthode par substitution.} $(E_2)$ donne $y = 16 - 5x$. On reporte dans $(E_1)$ :
\[ 3x - 2(16 - 5x) = 7 \iff 3x - 32 + 10x = 7 \iff 13x = 39 \iff x = 3. \]
Alors $y = 16 - 5\times 3 = 1$.
\item \textbf{Méthode par combinaisons linéaires.}
\begin{itemize}
\item Pour éliminer $y$ : $(E_1) + 2(E_2)$ : $(3x-2y) + 2(5x+y) = 7 + 32 \iff 13x = 39 \iff x = 3$.
\item Pour éliminer $x$ : $5(E_1) - 3(E_2)$ : $5(3x-2y) - 3(5x+y) = 35 - 48 \iff -13y = -13 \iff y = 1$.
\end{itemize}
\item \textbf{Méthode de Cramer (déterminants).}
$D = \begin{vmatrix} 3 & -2 \\ 5 & 1 \end{vmatrix} = 3\times 1 - (-2)\times 5 = 3 + 10 = 13 \neq 0$.
\[ D_x = \begin{vmatrix} 7 & -2 \\ 16 & 1 \end{vmatrix} = 7\times 1 - (-2)\times 16 = 7 + 32 = 39 \quad\Rightarrow\quad x = \frac{39}{13} = 3. \]
\[ D_y = \begin{vmatrix} 3 & 7 \\ 5 & 16 \end{vmatrix} = 3\times 16 - 7\times 5 = 48 - 35 = 13 \quad\Rightarrow\quad y = \frac{13}{13} = 1. \]
\item \textbf{Commentaire.} La substitution est la plus rapide ici car le coefficient de $y$ dans $(E_2)$ vaut $1$. Cramer est systématique mais exige trois déterminants ; les combinaisons sont efficaces quand les coefficients s'éliminent facilement (ici $y$ s'élimine directement).
\end{enumerate}
Vérification : $3\times 3 - 2\times 1 = 7$ et $5\times 3 + 1 = 16$. $\mathcal{S} = \{(3 ; 1)\}$.
\end{corrige}

\begin{corrige}[Exercice 6 — Compter avant de résoudre (discussion sans paramètre)]
\begin{enumerate}
\item $\begin{cases} x - 4y = 3 \\ 2x + 3y = 17 \end{cases}$.
$D = \begin{vmatrix} 1 & -4 \\ 2 & 3 \end{vmatrix} = 1\times 3 - (-4)\times 2 = 3 + 8 = 11 \neq 0$ : solution unique.
\[ x = \frac{\begin{vmatrix} 3 & -4 \\ 17 & 3 \end{vmatrix}}{11} = \frac{9 + 68}{11} = 7, \qquad
   y = \frac{\begin{vmatrix} 1 & 3 \\ 2 & 17 \end{vmatrix}}{11} = \frac{17 - 6}{11} = 1. \]
$\mathcal{S} = \{(7 ; 1)\}$.

\item $\begin{cases} 6x - 9y = 2 \\ 4x - 6y = 5 \end{cases}$.
$D = \begin{vmatrix} 6 & -9 \\ 4 & -6 \end{vmatrix} = 6\times(-6) - (-9)\times 4 = -36 + 36 = 0$ : pas de solution unique.
On compare les coefficients : $\frac{6}{4} = \frac{-9}{-6} = \frac{3}{2}$ mais $\frac{2}{5} \neq \frac{3}{2}$. Les droites sont strictement parallèles : \textbf{aucune solution}, $\mathcal{S} = \emptyset$.

\item $\begin{cases} 3x + 2y = 16 \\ 5x - y = 5 \end{cases}$.
$D = \begin{vmatrix} 3 & 2 \\ 5 & -1 \end{vmatrix} = 3\times(-1) - 2\times 5 = -3 - 10 = -13 \neq 0$ : solution unique.
\[ x = \frac{\begin{vmatrix} 16 & 2 \\ 5 & -1 \end{vmatrix}}{-13} = \frac{-16 - 10}{-13} = 2, \qquad
   y = \frac{\begin{vmatrix} 3 & 16 \\ 5 & 5 \end{vmatrix}}{-13} = \frac{15 - 80}{-13} = 5. \]
$\mathcal{S} = \{(2 ; 5)\}$.
\end{enumerate}
\end{corrige}

\begin{corrige}[Exercice 7 — Pivot de Gauss dans $\mathbb{R}^3$]
Soit le système :
\[ (S) \begin{cases} x - 2y + z = 1 & (E_1) \\ 2x + y - z = 1 & (E_2) \\ 3x - y + 2z = 6 & (E_3) \end{cases} \]

\begin{enumerate}
\item \textbf{Élimination (descente).}
\begin{itemize}
\item On élimine $x$ de $(E_2)$ et $(E_3)$ :
\[ (E_2) \leftarrow (E_2) - 2(E_1) : \quad (2x+y-z) - 2(x-2y+z) = 1-2 \;\Rightarrow\; 5y - 3z = -1 \quad (E_2') \]
\[ (E_3) \leftarrow (E_3) - 3(E_1) : \quad (3x-y+2z) - 3(x-2y+z) = 6-3 \;\Rightarrow\; 5y - z = 3 \quad (E_3') \]
\item On élimine $y$ de $(E_3')$ :
\[ (E_3) \leftarrow (E_3') - (E_2') : \quad (5y-z) - (5y-3z) = 3 - (-1) \;\Rightarrow\; 2z = 4 \quad (E_3'') \]
\end{itemize}
Système triangulaire équivalent obtenu :
\[ \begin{cases} x - 2y + z = 1 & (E_1) \\ 5y - 3z = -1 & (E_2') \\ 2z = 4 & (E_3'') \end{cases} \]

\item \textbf{Remontée.}
$2z = 4 \;\Rightarrow\; z = 2$.
$5y - 3(2) = -1 \;\Rightarrow\; 5y = 5 \;\Rightarrow\; y = 1$.
$x - 2(1) + 2 = 1 \;\Rightarrow\; x = 1$.

\item \textbf{Vérification dans le système d'origine.}
$E_1 : 1 - 2 + 2 = 1$ \checkmark \quad
$E_2 : 2 + 1 - 2 = 1$ \checkmark \quad
$E_3 : 3 - 1 + 4 = 6$ \checkmark
$\mathcal{S} = \{(1 ; 1 ; 2)\}$.
\end{enumerate}
\end{corrige}

\begin{corrige}[Exercice 8 — Substitution et combinaisons dans $\mathbb{R}^3$]
On résout deux systèmes différents pour illustrer chaque méthode.

\textbf{1. Système favorable à la substitution :}
\[ (S_1) \begin{cases} x + y + z = 6 & (E_1) \\ x - y = 1 & (E_2) \\ y + z = 5 & (E_3) \end{cases} \]
\begin{itemize}
\item $(E_2)$ donne $x = y + 1$.
\item Report dans $(E_1)$ : $(y+1) + y + z = 6 \;\Rightarrow\; 2y + z = 5 \quad (E_1')$.
\item Avec $(E_3)$ : $y + z = 5$, on soustrait : $(E_1') - (E_3)$ donne $y = 0$.
\item Alors $z = 5$ (par $(E_3)$) et $x = 1$ (par $(E_2)$).
\end{itemize}
Vérification : $E_1: 1+0+5=6$ ; $E_2: 1-0=1$ ; $E_3: 0+5=5$. $\mathcal{S}_1 = \{(1 ; 0 ; 5)\}$.

\textbf{2. Système favorable aux combinaisons linéaires :}
\[ (S_2) \begin{cases} x + y - z = 0 & (E_1) \\ 2x - y + z = 3 & (E_2) \\ x + 2y + z = 8 & (E_3) \end{cases} \]
\begin{itemize}
\item $(E_1) + (E_2)$ élimine $y$ et $z$ : $3x = 3 \;\Rightarrow\; x = 1$.
\item On reporte $x=1$ dans $(E_1)$ et $(E_3)$ :
\[ (E_1) \;\Rightarrow\; 1 + y - z = 0 \;\Rightarrow\; y - z = -1 \quad (E_1') \]
\[ (E_3) \;\Rightarrow\; 1 + 2y + z = 8 \;\Rightarrow\; 2y + z = 7 \quad (E_3') \]
\item $(E_1') + (E_3')$ : $3y = 6 \;\Rightarrow\; y = 2$.
\item Alors $z = y + 1 = 3$ (par $(E_1')$).
\end{itemize}
Vérification : $E_1: 1+2-3=0$ ; $E_2: 2-2+3=3$ ; $E_3: 1+4+3=8$. $\mathcal{S}_2 = \{(1 ; 2 ; 3)\}$.
\end{corrige}

\begin{corrige}[Exercice 9 — La librairie de quartier]
\begin{enumerate}
\item Soit $x$ le nombre de cahiers de 100 pages et $y$ le nombre de cahiers de 200 pages.
Nombre total : $x + y = 50$.
Dépense totale (en francs CFA) : $500x + 800y = \num{31000}$.
\[ (S) \begin{cases} x + y = 50 & (E_1) \\ 500x + 800y = \num{31000} & (E_2) \end{cases} \]

\item On simplifie $(E_2)$ en divisant par $100$ : $5x + 8y = 310 \quad (E_2')$.
Combinaison : $(E_2') - 5(E_1)$ : $(5x+8y) - 5(x+y) = 310 - 250 \;\Rightarrow\; 3y = 60 \;\Rightarrow\; y = 20$.
Alors $x = 50 - 20 = 30$.
Le commerçant achète \textbf{30 cahiers de 100 pages} et \textbf{20 cahiers de 200 pages}.
Vérification : $500\times 30 + 800\times 20 = \num{15000} + \num{16000} = \num{31000}$~F.

\item Bénéfice unitaire : $650 - 500 = 150$~F (100 pages) ; $\num{1000} - 800 = 200$~F (200 pages).
Bénéfice total : $30\times 150 + 20\times 200 = \num{4500} + \num{4000} = \num{8500}$~F.

\item Nouveau système : $\begin{cases} x + y = 60 \\ 500x + 800y = \num{36000} \end{cases}$.
Déterminant : $D = \begin{vmatrix} 1 & 1 \\ 500 & 800 \end{vmatrix} = 800 - 500 = 300 \neq 0$ : solution unique.
Simplification : $5x + 8y = 360$. $(E_2) - 5(E_1)$ : $3y = 60 \;\Rightarrow\; y = 20$, puis $x = 40$.
La solution $(40 ; 20)$, formée d'entiers naturels, est réaliste.
\end{enumerate}
\textbf{Barème indicatif (sur 4 pts)} : mise en équations 1 pt ; résolution 1 pt ; bénéfice 1 pt ; déterminant et discussion 1 pt.\\
\textbf{Points de vigilance} : bien simplifier les équations avant de calculer (diviser par 100) ; vérifier que la solution est constituée d'entiers naturels (contexte concret) ; conclure par une phrase interprétant le résultat.
\end{corrige}

\begin{corrige}[Exercice 10 — Le bénéfice des trois commerçantes]
\begin{enumerate}
\item Soient $a$, $b$, $c$ les parts (en francs CFA) d'Amina, Brigitte, Chantal.
Partage total : $a + b + c = \num{450000} \quad (E_1)$.
« Amina reçoit le double de Brigitte » : $a = 2b \;\Leftrightarrow\; a - 2b = 0 \quad (E_2)$.
« Chantal reçoit $\num{50000}$~F de plus que Brigitte » : $c = b + \num{50000} \;\Leftrightarrow\; -b + c = \num{50000} \quad (E_3)$.

\item \textbf{Triangulation (pivot de Gauss).}
\begin{itemize}
\item On élimine $a$ : $(E_2') \leftarrow (E_1) - (E_2)$ : $(a+b+c) - (a-2b) = \num{450000} - 0 \;\Rightarrow\; 3b + c = \num{450000}$.
\item Système intermédiaire :
\[ \begin{cases} a + b + c = \num{450000} & (E_1) \\ 3b + c = \num{450000} & (E_2') \\ -b + c = \num{50000} & (E_3) \end{cases} \]
\item On élimine $c$ : $(E_3') \leftarrow (E_2') - (E_3)$ : $(3b+c) - (-b+c) = \num{450000} - \num{50000} \;\Rightarrow\; 4b = \num{400000}$.
\item Système triangulaire :
\[ \begin{cases} a + b + c = \num{450000} \\ 3b + c = \num{450000} \\ 4b = \num{400000} \end{cases} \]
\end{itemize}

\item \textbf{Remontée.}
$4b = \num{400000} \;\Rightarrow\; b = \num{100000}$.
$c = \num{450000} - 3b = \num{450000} - \num{300000} = \num{150000}$ (ou $c = b + \num{50000} = \num{150000}$).
$a = \num{450000} - b - c = \num{450000} - \num{100000} - \num{150000} = \num{200000}$ (ou $a = 2b = \num{200000}$).

\item \textbf{Réponse.} Amina : $\num{200000}$~F ; Brigitte : $\num{100000}$~F ; Chantal : $\num{150000}$~F.

\item \textbf{Vérification.}
$\num{200000} + \num{100000} + \num{150000} = \num{450000}$ \checkmark
$\num{200000} = 2\times \num{100000}$ \checkmark
$\num{150000} = \num{100000} + \num{50000}$ \checkmark

\item \textbf{Bonus (âge du père et du fils).}
Soit $p$ l'âge du père, $f$ l'âge du fils.
$p = 3f \;\Leftrightarrow\; p - 3f = 0$.
Dans 12 ans : $p + 12 = 2(f + 12) \;\Leftrightarrow\; p - 2f = 12$.
Système : $\begin{cases} p - 3f = 0 \\ p - 2f = 12 \end{cases}$.
Soustraction : $(p-2f) - (p-3f) = 12 - 0 \;\Rightarrow\; f = 12$.
Alors $p = 3\times 12 = 36$.
Le père a $36$ ans, le fils $12$ ans. Vérification : dans $12$ ans, $48 = 2\times 24$.
\end{enumerate}
\textbf{Barème indicatif (sur 5 pts)} : traduction 1 pt ; triangulation détaillée 2 pts ; remontée et parts 1 pt ; vérification et bonus 1 pt.\\
\textbf{Points de vigilance} : écrire explicitement chaque opération élémentaire $(E_i) \leftarrow \dots$ ; vérifier la solution dans les \emph{trois} équations d'origine ; ne pas oublier l'unité (francs) dans la conclusion.
\end{corrige}

\begin{corrige}[Exercice 11 — Discussion avec paramètre]
\begin{enumerate}
\item Soit le système $(S_m)$ :
\[ \begin{cases} mx + 2y = m+1 \\ 2x + my = 4 \end{cases} \]
Déterminant : $D(m) = \begin{vmatrix} m & 2 \\ 2 & m \end{vmatrix} = m\times m - 2\times 2 = m^2 - 4 = (m-2)(m+2)$.

\item \textbf{Discussion selon $m$.}
\begin{itemize}
\item \textbf{Cas 1 : $m \neq 2$ et $m \neq -2$.} Alors $D(m) \neq 0$. D'après le théorème de Cramer, $(S_m)$ admet une \textbf{unique solution}.
\item \textbf{Cas 2 : $m = 2$.} Le système devient :
\[ \begin{cases} 2x + 2y = 3 \\ 2x + 2y = 4 \end{cases} \]
Les deux équations sont incompatibles (membres de gauche identiques, seconds membres différents) : \textbf{aucune solution} ($\mathcal{S} = \emptyset$). Droites strictement parallèles.
\item \textbf{Cas 3 : $m = -2$.} Le système devient :
\[ \begin{cases} -2x + 2y = -1 \\ 2x - 2y = 4 \end{cases} \;\Leftrightarrow\; \begin{cases} -2x + 2y = -1 \\ -2x + 2y = -4 \end{cases} \]
Incompatibles ($-1 \neq -4$) : \textbf{aucune solution} ($\mathcal{S} = \emptyset$). Droites strictement parallèles.
\end{itemize}
\textbf{Synthèse :} $\mathcal{S} \neq \emptyset \iff m \notin \{-2; 2\}$ (et alors solution unique). Pour $m \in \{-2; 2\}$, $\mathcal{S} = \emptyset$.

\item \textbf{Résolution pour $m = 3$.}
$D(3) = 9 - 4 = 5 \neq 0$.
Formules de Cramer :
\[ x = \frac{\begin{vmatrix} 4 & 2 \\ 4 & 3 \end{vmatrix}}{5} = \frac{12 - 8}{5} = \frac{4}{5}, \qquad
   y = \frac{\begin{vmatrix} 3 & 4 \\ 2 & 4 \end{vmatrix}}{5} = \frac{12 - 8}{5} = \frac{4}{5}. \]
Vérification : $3\times\frac{4}{5} + 2\times\frac{4}{5} = \frac{20}{5} = 4 = m+1$ ; $2\times\frac{4}{5} + 3\times\frac{4}{5} = \frac{20}{5} = 4$.
$\mathcal{S} = \left\{\left(\frac{4}{5} ; \frac{4}{5}\right)\right\}$.

\item \textbf{Système $(T_m)$ :} $\begin{cases} x + y = m \\ 2x - y = 3 \end{cases}$.
\begin{enumerate}
\item $D = \begin{vmatrix} 1 & 1 \\ 2 & -1 \end{vmatrix} = -1 - 2 = -3 \neq 0$ pour tout $m$ : solution unique toujours.
Addition des deux équations : $3x = m + 3 \;\Rightarrow\; x = \dfrac{m+3}{3}$.
Puis $y = m - x = m - \dfrac{m+3}{3} = \dfrac{2m - 3}{3}$.
\item $x > 0 \iff m+3 > 0 \iff m > -3$.
$y > 0 \iff 2m - 3 > 0 \iff m > \dfrac{3}{2}$.
Les deux conditions sont réalisées simultanément si et seulement si $m > \dfrac{3}{2}$.
\end{enumerate}
\end{enumerate}
\textbf{Barème indicatif (sur 5 pts)} : déterminant factorisé 1 pt ; discussion des trois cas 1,5 pt ; résolution pour $m = 3$ 1 pt ; solution de $(T_m)$ en fonction de $m$ 1 pt ; inéquations et conclusion 0,5 pt.\\
\textbf{Points de vigilance} : citer le théorème « $D \neq 0 \iff$ solution unique » ; pour les valeurs annulant le déterminant, \emph{remplacer} $m$ dans le système avant de conclure (ne jamais se contenter de $D = 0$) ; conclure la discussion par une synthèse claire des cas (tableau de valeurs ou phrase récapitulative).
\end{corrige}

\newpage

\coursec{Fiche bilan}

\begin{popfigure}
\begin{tikzpicture}[mindmap, grow cyclic, every node/.style={concept, text=white, font=\small\bfseries}, root concept/.append style={concept color=popdark, font=\large\bfseries}, level 1/.append style={level distance=3.6cm, sibling angle=51.4}, level 2/.append style={level distance=2.2cm, font=\footnotesize\bfseries}]
\node[root concept]{Systèmes linéaires}
  child[concept color=popblue]{ node {Dans $\mathbb{R}^2$ : définitions}
    child { node {Couple solution} }
    child { node {2 droites du plan} }
  }
  child[concept color=poporange]{ node {Substitution et combinaisons}
    child { node {Isoler une inconnue} }
    child { node {Éliminer une inconnue} }
  }
  child[concept color=popgreen]{ node {Déterminant et Cramer}
    child { node {$D=ab'-a'b$} }
    child { node {$x=\dfrac{D_x}{D}$, $y=\dfrac{D_y}{D}$} }
  }
  child[concept color=popgold]{ node {Systèmes à paramètre}
    child { node {Discuter selon $m$} }
  }
  child[concept color=poppurple]{ node {Dans $\mathbb{R}^3$}
    child { node {Pivot de Gauss} }
    child { node {Système triangulaire} }
  }
  child[concept color=poppink]{ node {Interprétation géométrique}
    child { node {Sécantes : 1 solution} }
    child { node {Parallèles : 0 ou $\infty$} }
  }
  child[concept color=popmuted]{ node {Problèmes concrets}
    child { node {Mise en équations} }
    child { node {Vérification} }
  };
\end{tikzpicture}
\legende{Carte mentale de la leçon : systèmes d'équations linéaires.}
\end{popfigure}

\begin{bilan}
\textbf{Définitions clés}
\begin{itemize}
\item Une \cle{équation linéaire} à deux inconnues est de la forme $ax+by=c$ ; ses solutions sont les couples $(x\,;y)$ qui la vérifient.
\item Un \cle{système linéaire} $2\times 2$ : $\begin{cases} ax+by=c \\ a'x+b'y=c' \end{cases}$. Le résoudre, c'est trouver tous les couples $(x\,;y)$ vérifiant \textbf{simultanément} les deux équations.
\item Un système $3\times 3$ fait intervenir trois inconnues $x$, $y$, $z$ ; sa solution (quand elle est unique) est un \textbf{triplet} $(x\,;y\,;z)$.
\item Le \cle{déterminant} du système $2\times 2$ est $D=ab'-a'b$.
\end{itemize}

\textbf{Formules à connaître par c\oe ur}
\begin{center}
\begin{tabularx}{\linewidth}{|l|X|X|}\hline
\rowcolor{popblueL} \textbf{Objet} & \textbf{Formule} & \textbf{Conclusion} \\ \hline
Déterminant & $D=\begin{vmatrix} a & b \\ a' & b' \end{vmatrix}=ab'-a'b$ & $D\neq 0$ : solution unique \\ \hline
Formules de Cramer & $x=\dfrac{D_x}{D}=\dfrac{cb'-c'b}{D}$ ; \quad $y=\dfrac{D_y}{D}=\dfrac{ac'-a'c}{D}$ & valables \textbf{uniquement} si $D\neq 0$ \\ \hline
Cas $D=0$ & droites parallèles & $0$ solution (strictement parallèles) ou infinité (confondues) \\ \hline
\end{tabularx}
\end{center}

\textbf{Méthodes express}
\begin{itemize}
\item \cle{Substitution} : on isole une inconnue dans une équation (celle dont le coefficient vaut $1$ ou $-1$ de préférence), on la remplace dans l'autre, on résout, puis on revient en arrière.
\item \cle{Combinaisons linéaires} : on multiplie les équations par des coefficients bien choisis pour éliminer une inconnue par addition membre à membre.
\item \cle{Pivot de Gauss} (dans $\mathbb{R}^3$) : on élimine $x$ des équations (2) et (3), puis $y$ de l'équation (3) ; on obtient un système \textbf{triangulaire} que l'on résout par remontée ($z$, puis $y$, puis $x$).
\item \cle{Système à paramètre} $m$ : on calcule $D(m)$ ; on discute : $D(m)\neq 0$ $\Rightarrow$ solution unique (Cramer) ; $D(m)=0$ $\Rightarrow$ on remplace $m$ par sa valeur et on étudie directement le système obtenu.
\item \cle{Problème concret} : choix des inconnues $\to$ traduction en équations $\to$ résolution $\to$ \textbf{vérification} et réponse avec unités.
\end{itemize}
\end{bilan}

\begin{aretenir}[Checklist avant le Probatoire]
\begin{itemize}
\item[$\square$] Je sais reconnaître un système linéaire et identifier ses coefficients $a$, $b$, $c$, $a'$, $b'$, $c'$ sans me tromper de signe.
\item[$\square$] Je sais résoudre un système $2\times 2$ par substitution.
\item[$\square$] Je sais résoudre un système $2\times 2$ par combinaisons linéaires.
\item[$\square$] Je sais calculer le déterminant $D=ab'-a'b$ et conclure sur l'unicité de la solution.
\item[$\square$] Je connais les formules de Cramer $x=\dfrac{D_x}{D}$ et $y=\dfrac{D_y}{D}$ et je sais les appliquer quand $D\neq 0$.
\item[$\square$] Je sais interpréter géométriquement les trois cas : droites sécantes, strictement parallèles, confondues.
\item[$\square$] Je sais discuter un système dont le second membre dépend d'un paramètre $m$ (valeurs de $m$ annulant $D$ traitées à part).
\item[$\square$] Je sais appliquer le pivot de Gauss pour triangulariser un système $3\times 3$.
\item[$\square$] Je sais résoudre un système triangulaire par remontée : $z$, puis $y$, puis $x$.
\item[$\square$] Je sais mettre un problème concret en équations et vérifier ma solution dans l'énoncé.
\item[$\square$] Je rédige proprement : j'écris l'ensemble solution $S=\{(x_0\,;y_0)\}$ ou $S=\{(x_0\,;y_0\,;z_0)\}$.
\item[$\square$] Je vérifie toujours ma solution en la substituant dans \textbf{toutes} les équations du système.
\end{itemize}
\end{aretenir}

\findecours

\end{document}
